
\title{\textbf{Mathematical Foundations, Architecture Diagrams,\\and State Analysis of Modern LLM Training Methods}}
\author{Chief Strategist J}
\date{\today}
\maketitle
\tableofcontents
\newpage

%% ============================================================
\part{Foundation Training Methods}
%% ============================================================

\section{Pretraining}

\subsection{Mathematical Foundation}
Causal language modeling models the probability of a sequence $x = (x_1, \ldots, x_T)$ as the product of autoregressive conditional probabilities:
\begin{equation}
  P(x) = \prod_{t=1}^T P(x_t \mid x_{<t})
\end{equation}
The hidden representation $h_t^{(l)}$ at layer $l$ and token position $t$ is computed using multi-head causal self-attention. The causal attention mask matrix $M \in \mathbb{R}^{T \times T}$ is defined as:
\begin{equation}
  M_{i,j} = \begin{cases} 0 & \text{if } j \le i \\ -\infty & \text{if } j > i \end{cases}
\end{equation}
This mask is added directly to the scaled query-key dot products:
\begin{equation}
  \text{Attention}(Q, K, V) = \text{softmax}\left( \frac{Q K^\top}{\sqrt{d_k}} + M \right) V
\end{equation}
where $Q = X W_Q$, $K = X W_K$, $V = X W_V$. This prevents information leakage from future tokens. At the output layer, the logits $z_t \in \mathbb{R}^{|\mathcal{V}|}$ are computed via $z_t = h_t^{(L)} W_{\text{vocab}}$. To maintain numerical stability, softmax is computed by subtracting the maximum logit:
\begin{equation}
  P(x_t = v \mid x_{<t}) = \frac{\exp(z_{t,v} - \max_k z_{t,k})}{\sum_{j \in \mathcal{V}} \exp(z_{t,j} - \max_k z_{t,k})}
\end{equation}
The causal model parameters $\theta$ are trained by minimizing the cross-entropy loss over a dataset $\mathcal{D}$:
\begin{equation}
  \mathcal{L}_{\text{PT}}(\theta) = -\frac{1}{|\mathcal{D}|}\sum_{x \in \mathcal{D}} \sum_{t=1}^T \log P_\theta(x_t \mid x_{<t})
\end{equation}
Updates are made using the AdamW optimizer. Let $g_t = \nabla_\theta \mathcal{L}_t$ be the gradient at step $t$. The update equations are:
\begin{align}
  m_t &= \beta_1 m_{t-1} + (1-\beta_1) g_t \\
  v_t &= \beta_2 v_{t-1} + (1-\beta_2) g_t^2 \\
  \hat{m}_t &= \frac{m_t}{1-\beta_1^t}, \quad \hat{v}_t = \frac{v_t}{1-\beta_2^t} \\
  \theta_t &= \theta_{t-1} - \eta \left( \frac{\hat{m}_t}{\sqrt{\hat{v}_t} + \epsilon} + \lambda \theta_{t-1} \right)
\end{align}
where $\eta$ is the learning rate, $\lambda$ is the weight decay rate, $\beta_1, \beta_2$ are the exponential decay rates, and $\epsilon$ prevents division by zero. A cosine learning rate schedule with linear warmup is employed:
\begin{equation}
  \eta_t = \eta_{\text{min}} + \frac{1}{2}(\eta_{\text{max}} - \eta_{\text{min}})\left(1 + \cos\left(\frac{t - T_{\text{warmup}}}{T_{\text{max}} - T_{\text{warmup}}}\pi\right)\right)
\end{equation}

\subsection{Architecture Flow Diagram}

\begin{figure}[h]
\centering
\begin{adjustbox}{max width=\linewidth}
\begin{tikzpicture}[node distance=0.9cm and 1.4cm]
  \node[datblk] (corpus) {\textbf{Raw Text} \\ (Corpus $\mathcal{D}$)};
  \node[grayblk, right=of corpus] (tok) {\textbf{Tokenise} \\ ($x_{1..T}$)};
  \node[frzblk, right=of tok] (fwd) {\textbf{Transformer} \\ (Forward Pass)};
  \node[lossblk, right=of fwd] (loss) {\textbf{Cross-Entropy} \\ (Loss $\mathcal{L}$)};
  \node[trnblk, below=0.8cm of loss] (bwd) {\textbf{Backprop +} \\ (Adam Step)};
  \node[mrgblk, left=3.6cm of bwd] (model) {\textbf{Base Model} \\ ($\theta^*$)};

  \draw[arr] (corpus)--(tok);
  \draw[arr] (tok)--(fwd);
  \draw[arr] (fwd)--(loss);
  \draw[arr] (loss)--(bwd);
  \draw[arr] (bwd)--node[below,font=\tiny]{update $\theta$}(model);
  \draw[arr] (model.north) |- ($(fwd.south)+(0,-0.3)$) -- (fwd.south);
\end{tikzpicture}
\end{adjustbox}
\caption{Pretraining pipeline: corpus sampling, forward pass, loss computation, and weight update loop.}
\end{figure}

\subsection{State Transition Table}
\begin{table}[h]\centering
\begin{tabular}{lll}\toprule
\textbf{Property} & \textbf{Before} & \textbf{After}\\\midrule
Parameter source & Random init $\theta_0 \sim \mathcal{N}(0,\sigma^2)$ & Converged weights $\theta^*$\\
Training corpus & Not consumed & Trillions of tokens processed\\
Perplexity & Very high ($>10^3$) & Low ($<20$ on held-out data)\\
World knowledge & None & Broad language \& factual knowledge\\
requires\_grad & True (all params) & True (all params)\\\bottomrule
\end{tabular}
\caption{State properties before and after pretraining.}
\end{table}

\newpage
%% ------------------------------------------------------------
\section{Continued Pretraining (CPT)}

\subsection{Mathematical Foundation}
Continued Pretraining (CPT) initializes parameters with a pretrained checkpoint $\theta_{\text{base}}$ and continues autoregressive training. To mitigate catastrophic forgetting, training batches are sampled from a mixture of the new domain-specific corpus $\mathcal{D}_{\text{new}}$ and a replay buffer of general pretraining data $\mathcal{D}_{\text{replay}}$:
\begin{equation}
  x \sim \lambda \mathcal{D}_{\text{new}} + (1 - \lambda) \mathcal{D}_{\text{replay}}
\end{equation}
where $\lambda \in [0,1]$ is the mixing ratio. The CPT objective minimizes the cross-entropy loss over this mixed distribution:
\begin{equation}
  \mathcal{L}_{\text{CPT}}(\theta) = -\mathbb{E}_{x \sim \mathcal{D}_{\text{mix}}} \sum_{t=1}^T \log P_\theta(x_t \mid x_{<t})
\end{equation}
The parameter updates follow the AdamW optimizer starting from $\theta_0 = \theta_{\text{base}}$, typically using a warm-restart cosine learning rate schedule:
\begin{equation}
  \eta_t = \eta_{\text{min}} + \frac{1}{2}(\eta_{\text{max}} - \eta_{\text{min}})\left(1 + \cos\left(\frac{t}{T_{\text{max}}}\pi\right)\right)
\end{equation}
where the peak learning rate $\eta_{\text{max}}$ is set to a lower value than during the initial pretraining phase to preserve general capabilities.

\subsection{Architecture Flow Diagram}

\begin{figure}[h]
\centering
\begin{adjustbox}{max width=\linewidth}
\begin{tikzpicture}[node distance=0.9cm and 1.3cm]
  \node[frzblk] (base) {\textbf{Pretrained} \\ (Base $\theta_{\text{base}}$)};
  \node[datblk, above right=0.6cm and 1.4cm of base] (new) {\textbf{New Corpus} \\ ($\mathcal{D}_{\text{new}}$)};
  \node[grayblk, below right=0.6cm and 1.4cm of base] (replay) {\textbf{Replay Buffer} \\ ($\mathcal{D}_{\text{replay}}$)};
  \node[grayblk, right=3.2cm of base] (mix) {\textbf{Batch} \\ (Mixer $\lambda$)};
  \node[lossblk, right=of mix] (loss) {\textbf{LM Loss} \\ ($\mathcal{L}_{\text{CPT}}$)};
  \node[trnblk, right=of loss] (upd) {\textbf{Gradient} \\ (Update)};
  \node[mrgblk, below=1.1cm of upd] (out) {\textbf{Adapted} \\ (Model $\theta'$)};

  \draw[arr] (new.east) -| (mix.north);
  \draw[arr] (replay.east) -| (mix.south);
  \draw[arr] (base) -- (mix);
  \draw[arr] (mix) -- (loss);
  \draw[arr] (loss) -- (upd);
  \draw[arr] (upd) -- (out);
\end{tikzpicture}
\end{adjustbox}
\caption{CPT mixes new corpus and replay buffer batches to update a pretrained checkpoint.}
\end{figure}

\subsection{State Transition Table}
\begin{table}[h]\centering
\begin{tabular}{lll}\toprule
\textbf{Property} & \textbf{Before} & \textbf{After}\\\midrule
Starting checkpoint & Pretrained $\theta_{\text{base}}$ & $\theta' = \theta_{\text{base}} + \Delta\theta$\\
New domain knowledge & Absent & Incorporated\\
Original knowledge & Retained & Partially retained via replay\\
Learning rate schedule & Finished & Warm-restart applied\\
Perplexity on new domain & High & Significantly reduced\\\bottomrule
\end{tabular}
\caption{State properties before and after continued pretraining.}
\end{table}

\newpage
%% ------------------------------------------------------------
\section{Domain Adaptive Pretraining (DAPT)}

\subsection{Mathematical Foundation}
Domain Adaptive Pretraining (DAPT) adapts a general-domain pretrained model $\theta_{\text{base}}$ to a target domain (e.g., medical, legal, or code) by training on a large, unlabelled domain-specific corpus $\mathcal{D}_{\text{domain}}$.
The objective function is the causal language modeling loss over the domain dataset:
\begin{equation}
  \mathcal{L}_{\text{DAPT}}(\theta) = -\mathbb{E}_{x \sim \mathcal{D}_{\text{domain}}} \sum_{t=1}^T \log P_\theta(x_t \mid x_{<t})
\end{equation}
where parameters are initialized at $\theta_0 = \theta_{\text{base}}$. DAPT uses a lower learning rate than initial pretraining to prevent catastrophic disruption of base semantic features while allowing the model to adapt vocabulary representations and learn domain-specific concepts.

\subsection{Architecture Flow Diagram}

\begin{figure}[h]
\centering
\begin{adjustbox}{max width=\linewidth}
\begin{tikzpicture}[node distance=0.9cm and 1.2cm]
  \node[datblk] (domain) {\textbf{Domain Corpus} \\ (Medical/Legal/Code)};
  \node[grayblk, right=of domain] (clean) {\textbf{Clean \&} \\ (Deduplicate)};
  \node[frzblk, right=of clean] (base) {\textbf{Base Model} \\ ($\theta_{\text{base}}$)};
  \node[lossblk, right=of base] (loss) {\textbf{Domain LM} \\ (Loss $\mathcal{L}$)};
  \node[trnblk, below=0.9cm of loss] (step) {\textbf{Adam Step} \\ (low LR)};
  \node[mrgblk, left=2.2cm of step] (dapt) {\textbf{DAPT Model} \\ ($\theta_{\text{domain}}$)};

  \draw[arr] (domain)--(clean);
  \draw[arr] (clean)--(base);
  \draw[arr] (base)--(loss);
  \draw[arr] (loss)--(step);
  \draw[arr] (step)--node[below,font=\tiny]{converged}(dapt);
  \draw[arr,dashed,blue!60] (dapt.north) -- ++(0,0.5) -| (base.south) node[midway,above,font=\tiny]{checkpoint};
\end{tikzpicture}
\end{adjustbox}
\caption{DAPT pipeline: domain corpus is cleaned, fed into the base model, and trained with standard LM loss.}
\end{figure}

\subsection{State Transition Table}
\begin{table}[h]\centering
\begin{tabular}{lll}\toprule
\textbf{Property} & \textbf{Before} & \textbf{After}\\\midrule
Domain vocabulary familiarity & Low & High\\
Domain perplexity & High ($\rho \gg 1$) & Low (near 1)\\
General language ability & Strong & Marginally reduced\\
Downstream task (same domain) & Weak & Significantly stronger\\
Training data type & General web text & Domain-specific unlabelled\\\bottomrule
\end{tabular}
\caption{State properties before and after domain adaptive pretraining.}
\end{table}

\newpage
%% ------------------------------------------------------------
\section{Task Adaptive Pretraining (TAPT)}

\subsection{Mathematical Foundation}
Task Adaptive Pretraining (TAPT) optimizes a pretrained or domain-adapted checkpoint $\theta_{\text{base}}$ on a task-specific unlabelled dataset $\mathcal{D}_{\text{task}}$:
\begin{equation}
  \mathcal{L}_{\text{TAPT}}(\theta) = -\mathbb{E}_{x \sim \mathcal{D}_{\text{task}}} \sum_{t=1}^T \log P_\theta(x_t \mid x_{<t})
\end{equation}
Since the task dataset size $|\mathcal{D}_{\text{task}}|$ is typically much smaller than a general or domain-specific corpus, TAPT uses very few epochs, a small learning rate, and regularizers (such as dropout) to prevent overfitting before final supervised fine-tuning.

\subsection{Architecture Flow Diagram}

\begin{figure}[h]
\centering
\begin{adjustbox}{max width=\linewidth}
\begin{tikzpicture}[node distance=0.8cm and 1.2cm]
  \node[datblk] (labeled) {\textbf{Labelled Task} \\ (Data (small))};
  \node[grayblk, above=0.7cm of labeled] (unlabeled) {\textbf{Unlabelled Task-} \\ (Adjacent Text)};
  \node[frzblk, right=2.0cm of labeled] (dapt) {\textbf{DAPT / Base} \\ (Model $\theta$)};
  \node[lossblk, right=of dapt] (loss) {\textbf{Task LM} \\ (Loss)};
  \node[trnblk, right=of loss] (update) {\textbf{Update} \\ ($\theta \to \theta_{\text{task}}$)};
  \node[mrgblk, below=0.9cm of update] (sft) {\textbf{SFT} \\ (next step)};

  \draw[arr] (unlabeled.east) -| (dapt.north);
  \draw[arr] (labeled.east) -- node[above,font=\tiny]{strip labels} (dapt.west);
  \draw[arr] (dapt)--(loss);
  \draw[arr] (loss)--(update);
  \draw[arr] (update)--(sft);
\end{tikzpicture}
\end{adjustbox}
\caption{TAPT uses task-adjacent unlabelled text to bridge base model and supervised fine-tuning.}
\end{figure}

\subsection{State Transition Table}
\begin{table}[h]\centering
\begin{tabular}{lll}\toprule
\textbf{Property} & \textbf{Before} & \textbf{After}\\\midrule
Task vocabulary alignment & Partial & Strong\\
Corpus size & Millions of docs & Thousands of docs\\
Training epochs & 1--3 & 10--100 (small corpus)\\
Downstream label efficiency & Requires more labels & Requires fewer labels\\
Transfer to other tasks & Good & Reduced (narrow focus)\\\bottomrule
\end{tabular}
\caption{State properties before and after task adaptive pretraining.}
\end{table}

\newpage
%% ============================================================
\part{Supervised Fine-Tuning Methods}
%% ============================================================

\section{Supervised Fine-Tuning (SFT)}

\subsection{Mathematical Foundation}
SFT trains the model on structured prompt-response pairs $(x, y) \sim \mathcal{D}_{\text{SFT}}$ using teacher forcing. Let $s = [x ; y]$ be the concatenated token sequence of length $M + U$, where $x = (s_1, \ldots, s_M)$ is the prompt and $y = (s_{M+1}, \ldots, s_{M+U})$ is the response. The key math lies in applying a target label mask $L \in \{-100, \mathcal{V}\}^{M+U}$ defined as:
\begin{equation}
  L_t = \begin{cases} -100 & \text{if } t \le M \\ s_t & \text{if } t > M \end{cases}
\end{equation}
The cross-entropy loss function ignores targets set to $-100$, directing gradient updates exclusively to the response generation tokens:
\begin{equation}
  \mathcal{L}_{\text{SFT}}(\theta) = -\frac{1}{U} \sum_{t=M+1}^{M+U} \log P_\theta(s_t \mid s_{<t})
\end{equation}
This loss ignores the prompt distribution $P(x)$, preventing the model from degrading its output generation capability to match arbitrary prompts.
The gradients only backpropagate through response tokens:
\begin{equation}
  \nabla_\theta \mathcal{L}_{\text{SFT}} = -\frac{1}{U} \sum_{t=M+1}^{M+U} \nabla_\theta \log P_\theta(s_t \mid s_{<t})
\end{equation}

\subsection{Architecture Flow Diagram}

\begin{figure}[h]
\centering
\begin{adjustbox}{max width=\linewidth}
\begin{tikzpicture}[node distance=0.8cm and 1.2cm]
  \node[datblk] (pairs) {\textbf{$(x, y)$} \\ (Curated Pairs)};
  \node[grayblk, right=of pairs] (concat) {\textbf{Concat} \\ ({[x ; y]})};
  \node[frzblk, right=of concat] (fwd) {\textbf{Transformer} \\ (Forward)};
  \node[lossblk, right=of fwd] (loss) {\textbf{Masked CE} \\ (Loss $\mathcal{L}_{\text{SFT}}$)};
  \node[trnblk, below=0.9cm of loss] (bwd) {\textbf{Backprop} \\ (response only)};
  \node[mrgblk, left=3.5cm of bwd] (out) {\textbf{Instruction-} \\ (Following $\theta^*$)};

  \draw[arr] (pairs)--(concat);
  \draw[arr] (concat)--(fwd);
  \draw[arr] (fwd)--(loss);
  \draw[arr] (loss)--(bwd);
  \draw[arr] (bwd)--(out);

  \node[red,font=\tiny,text width=2cm,align=center] at ($(fwd.south)+(0,-0.3)$) {prompt tokens\\masked in labels};
\end{tikzpicture}
\end{adjustbox}
\caption{SFT pipeline with prompt-masked loss: gradients flow only through response tokens.}
\end{figure}

\subsection{State Transition Table}
\begin{table}[h]\centering
\begin{tabular}{lll}\toprule
\textbf{Property} & \textbf{Before} & \textbf{After}\\\midrule
Behaviour & Completion only & Instruction following\\
Loss region & N/A & Response tokens only\\
Label masking & None & Prompt masked to $-100$\\
Data format & Raw text & (instruction, response) pairs\\
Output quality & Random completion & Task-aligned responses\\\bottomrule
\end{tabular}
\caption{State properties before and after supervised fine-tuning.}
\end{table}

\newpage
%% ------------------------------------------------------------
\section{Instruction Tuning}

\subsection{Mathematical Foundation}
Instruction Tuning exposes the model to multi-turn conversational patterns. A prompt template $T(I, x)$ converts an instruction $I$ and context $x$ into structured sequences. Turn boundaries are defined using chat-formatting systems (like ChatML). A single sequence layout is:
\begin{equation}
  S = \langle\text{user}\rangle I \parallel x \langle\text{assistant}\rangle y \langle\text{eos}\rangle
\end{equation}
Let $\mathcal{D}_{\text{inst}}$ consist of instructions spanning hundreds of tasks. The model parameters $\theta$ are optimized on the joint multi-task dataset:
\begin{equation}
  \mathcal{L}_{\text{IT}}(\theta) = -\mathbb{E}_{(I, x, y) \sim \mathcal{D}_{\text{inst}}} \sum_{t=1}^{|y|} \log P_\theta\bigl(y_t \mid T(I, x), y_{<t}\bigr)
\end{equation}
We mask all prompt-side tokens, including conversational format tags (like $\langle\text{user}\rangle$ and instruction text), ensuring parameters are trained only on synthesizing assistant responses:
\begin{equation}
  \mathcal{L}_{\text{IT}}(\theta) = -\frac{1}{|y|} \sum_{i=1}^{|T(I,x)|+|y|} M_i \log P_\theta(S_i \mid S_{<i}), \quad M_i = \mathbb{I}(i > |T(I, x)|)
\end{equation}

\subsection{Architecture Flow Diagram}

\begin{figure}[h]
\centering
\begin{adjustbox}{max width=\linewidth}
\begin{tikzpicture}[node distance=0.7cm and 1.1cm]
  \node[datblk] (t1) {\textbf{Summarise} \\ (Task $\mathcal{T}_1$)};
  \node[datblk, below=0.4cm of t1] (t2) {\textbf{Translate} \\ (Task $\mathcal{T}_2$)};
  \node[datblk, below=0.4cm of t2] (t3) {\textbf{QA / Code} \\ (Task $\mathcal{T}_K$)};
  \node[grayblk, right=1.6cm of t2] (tmpl) {\textbf{Chat Template} \\ (Formatter)};
  \node[frzblk, right=of tmpl] (model) {\textbf{LLM} \\ ($\theta$)};
  \node[lossblk, right=of model] (loss) {\textbf{SFT Loss} \\ (multi-task)};
  \node[mrgblk, below=0.8cm of loss] (out) {\textbf{Instruction-} \\ (Tuned $\theta^*$)};

  \draw[arr] (t1.east) -| (tmpl.north);
  \draw[arr] (t2)--(tmpl);
  \draw[arr] (t3.east) -| (tmpl.south);
  \draw[arr] (tmpl)--(model);
  \draw[arr] (model)--(loss);
  \draw[arr] (loss)--(out);
\end{tikzpicture}
\end{adjustbox}
\caption{Instruction tuning over diverse task types funnelled through a unified chat template.}
\end{figure}

\subsection{State Transition Table}
\begin{table}[h]\centering
\begin{tabular}{lll}\toprule
\textbf{Property} & \textbf{Before} & \textbf{After}\\\midrule
Zero-shot generalisation & Poor & Strong across task types\\
Task diversity in training & N/A & $K$ heterogeneous tasks\\
Input format & Free-form text & Structured chat template\\
Prompt sensitivity & High & Reduced\\
Held-out task performance & Weak & Strong (FLAN, InstructGPT)\\\bottomrule
\end{tabular}
\caption{State properties before and after instruction tuning.}
\end{table}

\newpage
%% ------------------------------------------------------------
\section{Multi-Task Training}

\subsection{Mathematical Foundation}
Multi-Task Training aggregates losses from $K$ distinct tasks. A simple summation can result in tasks with large gradient scales dominating the parameter updates. To resolve this, we use the GradNorm algorithm. Let $W$ be the shared parameter weights (e.g., the last transformer block weights). We define the gradient norm for task $k$ at training step $t$ as:
\begin{equation}
  G_k(t) = \|\nabla_W (w_k(t) \mathcal{L}_k(t))\|_2
\end{equation}
Let $\bar{G}(t) = \frac{1}{K} \sum_k G_k(t)$ be the average gradient norm. The target gradient norm for task $k$ balances task difficulties:
\begin{equation}
  \hat{G}_k(t) = \bar{G}(t) \times \left( \frac{\tilde{\mathcal{L}}_k(t)}{\frac{1}{K}\sum_j \tilde{\mathcal{L}}_j(t)} \right)^\alpha
\end{equation}
where $\tilde{\mathcal{L}}_k(t) = \mathcal{L}_k(t) / \mathcal{L}_k(0)$ is the task loss ratio indicating task training progress, and $\alpha$ is a hyperparameter managing structural gradient balancing. The task weights $w_k(t)$ are updated by minimizing the L1 distance between $G_k(t)$ and $\hat{G}_k(t)$:
\begin{equation}
  \mathcal{L}_{\text{grad}}(w_1, \dots, w_K; t) = \sum_{k=1}^K \left| G_k(t) - \hat{G}_k(t) \right|
\end{equation}
The final joint parameter objective is:
\begin{equation}
  \mathcal{L}_{\text{MT}}(\theta) = \sum_{k=1}^K w_k \mathcal{L}_k(\theta)
\end{equation}

\subsection{Architecture Flow Diagram}

\begin{figure}[h]
\centering
\begin{adjustbox}{max width=\linewidth}
\begin{tikzpicture}[node distance=0.6cm and 1.1cm]
  \node[datblk] (d1) {\textbf{Dataset $\mathcal{D}_1$}};
  \node[datblk, below=0.4cm of d1] (d2) {\textbf{Dataset $\mathcal{D}_2$}};
  \node[datblk, below=0.4cm of d2] (dk) {\textbf{Dataset $\mathcal{D}_K$}};
  \node[frzblk, right=1.8cm of d2, minimum height=3.2cm, minimum width=3.0cm] (backbone) {\textbf{Shared} \\ \textbf{Transformer} \\ (Backbone $\theta_{\text{shared}}$)};
  \node[trnblk, right=1.6cm of backbone] (h2) {\textbf{Head$_2$} \\ ($\theta_2$)};
  \node[trnblk] (h1) at (h2 |- d1) {\textbf{Head$_1$} \\ ($\theta_1$)};
  \node[trnblk] (hk) at (h2 |- dk) {\textbf{Head$_K$} \\ ($\theta_K$)};
  \node[lossblk, right=1.6cm of h2] (loss) {\textbf{Weighted} \\ (Loss $\mathcal{L}_{\text{MT}}$)};

  \draw[arr] (d1)--(backbone.west|-h1);
  \draw[arr] (d2)--(backbone);
  \draw[arr] (dk)--(backbone.west|-hk);
  \draw[arr] (backbone.east|-h1)--(h1);
  \draw[arr] (backbone.east|-h2)--(h2);
  \draw[arr] (backbone.east|-hk)--(hk);
  \draw[arr] (h1.east) -| (loss.north);
  \draw[arr] (h2)--(loss);
  \draw[arr] (hk.east) -| (loss.south);
\end{tikzpicture}
\end{adjustbox}
\caption{Multi-task training: each dataset feeds task-specific heads on a shared backbone.}
\end{figure}

\subsection{State Transition Table}
\begin{table}[h]\centering
\begin{tabular}{lll}\toprule
\textbf{Property} & \textbf{Before} & \textbf{After}\\\midrule
Backbone specialisation & Single task & $K$ tasks jointly\\
Task heads & 1 & $K$ heads\\
Generalisation & Narrow & Broad\\
Data efficiency & Per-task only & Cross-task transfer\\
Gradient conflicts & N/A & Managed via $\lambda_k$\\\bottomrule
\end{tabular}
\caption{State properties before and after multi-task training.}
\end{table}

\newpage
%% ------------------------------------------------------------
\section{Chain-of-Thought SFT}

\subsection{Mathematical Foundation}
CoT SFT trains models to output intermediate reasoning steps before generating the final answer. Let $x$ be the prompt, $c = (c_1, \dots, c_M)$ be the reasoning chain tokens, and $y = (y_1, \dots, y_U)$ be the final answer tokens. The joint probability factorization is:
\begin{equation}
  P(c, y \mid x) = \prod_{i=1}^M P_\theta(c_i \mid x, c_{<i}) \prod_{j=1}^U P_\theta(y_j \mid x, c, y_{<j})
\end{equation}
The model is trained by minimizing the joint loss:
\begin{equation}
  \mathcal{L}_{\text{CoT}}(\theta) = -\beta \sum_{i=1}^M \log P_\theta(c_i \mid x, c_{<i}) - (1 - \beta) \sum_{j=1}^U \log P_\theta(y_j \mid x, c, y_{<j})
\end{equation}
During autoregressive generation, we sample tokens using Temperature-scaled softmax and Nucleus (top-$p$) filtering. The logits $z_t$ are modified as:
\begin{equation}
  \tilde{z}_{t,i} = \frac{z_{t,i}}{T}
\end{equation}
where $T > 0$ is the temperature. The top-$p$ vocabulary subset $\mathcal{V}^{(p)}$ at step $t$ is the smallest set satisfying:
\begin{equation}
  \sum_{v \in \mathcal{V}^{(p)}} P(v \mid x, s_{<t}) \ge p
\end{equation}
Tokens outside $\mathcal{V}^{(p)}$ are assigned probability 0 before final selection.

\subsection{Architecture Flow Diagram}

\begin{figure}[h]
\centering
\begin{adjustbox}{max width=\linewidth}
\begin{tikzpicture}[node distance=0.6cm and 0.9cm]
  \node[datblk] (q) {\textbf{Question $x$}};
  \node[grayblk, right=of q] (cot) {\textbf{Reasoning} \\ (Chain $c_1..c_M$)};
  \node[trnblk, right=of cot] (ans) {\textbf{Final} \\ (Answer $y$)};
  \node[lossblk, below=0.8cm of cot] (lc) {\textbf{Chain Loss} \\ ($\mathcal{L}_c$)};
  \node[lossblk, below=0.8cm of ans] (la) {\textbf{Answer Loss} \\ ($\mathcal{L}_y$)};
  \node[mrgblk, below=0.8cm of $(lc)!0.5!(la)$] (total) {\textbf{Total Loss} \\ ($\mathcal{L}_{\text{CoT}}$)};

  \draw[arr] (q)--(cot);
  \draw[arr] (cot)--(ans);
  \draw[arr] (cot)--(lc);
  \draw[arr] (ans)--(la);
  \draw[arr] (lc)--(total);
  \draw[arr] (la)--(total);

  \draw[losscolor,thick,decorate,decoration={brace,amplitude=6pt,raise=4pt}]
    (cot.north west) -- (ans.north east) node[midway,above=10pt,font=\tiny]{loss applied to entire output};
\end{tikzpicture}
\end{adjustbox}
\caption{CoT SFT: loss is computed over both the reasoning chain and the final answer.}
\end{figure}

\subsection{State Transition Table}
\begin{table}[h]\centering
\begin{tabular}{lll}\toprule
\textbf{Property} & \textbf{Before} & \textbf{After}\\\midrule
Output format & Direct answer & Reasoning chain + answer\\
Training signal per example & $|y|$ tokens & $|c| + |y|$ tokens\\
Multi-step reasoning & Weak & Strong\\
Data annotation cost & Low & High (chain annotation)\\
Accuracy on complex tasks & Low & Substantially improved\\\bottomrule
\end{tabular}
\caption{State properties before and after Chain-of-Thought SFT.}
\end{table}

\newpage
%% ------------------------------------------------------------
\section{Step-by-Step SFT}

\subsection{Mathematical Foundation}
Step-by-Step SFT formats the reasoning chain into discrete steps separated by a boundary token $T_{\text{step}} = \text{\texttt{[STEP]}}$. Let the response contain $S$ steps: $s = (s_1, T_{\text{step}}, s_2, T_{\text{step}}, \dots, s_S, T_{\text{step}})$. The step-level loss is:
\begin{equation}
  \mathcal{L}_{\text{SBS}}(\theta) = -\sum_{k=1}^S \sum_{t=1}^{|s_k|} \log P_\theta(s_{k,t} \mid x, s_{<k}, s_{k,<t}) - \sum_{k=1}^S \log P_\theta(T_{\text{step}} \mid x, s_{\le k})
\end{equation}
We define the average step perplexity $\text{PPL}_k$ as:
\begin{equation}
  \text{PPL}_k = \exp\left( -\frac{1}{|s_k|} \sum_{t=1}^{|s_k|} \log P_\theta(s_{k,t} \mid x, s_{<k}, s_{k,<t}) \right)
\end{equation}
Steps with high perplexity ($\text{PPL}_k > \tau$, where $\tau$ is a threshold) indicate low model confidence and can be flagged for step-level revision during generation.
The model is trained to minimize the combined cross-entropy loss over all steps.

\subsection{Architecture Flow Diagram}

\begin{figure}[h]
\centering
\begin{adjustbox}{max width=\linewidth}
\begin{tikzpicture}[node distance=0.5cm and 0.8cm]
  \node[datblk,minimum width=1.8cm] (inp) {\textbf{Input $x$}};
  \node[trnblk,right=0.7cm of inp,minimum width=1.8cm] (s1) {\textbf{Step 1} \\ ($s_1$)};
  \node[trnblk,right=0.6cm of s1,minimum width=1.8cm] (s2) {\textbf{Step 2} \\ ($s_2$)};
  \node[grayblk,right=0.6cm of s2,minimum width=0.6cm] (dots) {\textbf{$\cdots$}};
  \node[trnblk,right=0.6cm of dots,minimum width=1.8cm] (sn) {\textbf{Step $N$} \\ ($s_N$)};
  \node[mrgblk,right=0.7cm of sn,minimum width=1.8cm] (out) {\textbf{Answer} \\ ($y$)};

  \foreach \a/\b in {inp/s1, s1/s2, s2/dots, dots/sn, sn/out}
    \draw[arr] (\a)--(\b);

  \node[losscolor,font=\tiny] at ($(s1.south)+(0,-0.45)$) {$\mathcal{L}_1$};
  \node[losscolor,font=\tiny] at ($(s2.south)+(0,-0.45)$) {$\mathcal{L}_2$};
  \node[losscolor,font=\tiny] at ($(sn.south)+(0,-0.45)$) {$\mathcal{L}_N$};
  \node[losscolor,font=\tiny] at ($(out.south)+(0,-0.45)$) {$\mathcal{L}_y$};

  \foreach \n in {s1,s2,sn,out}
    \draw[garr] (\n.south) -- ++(0,-0.35);
\end{tikzpicture}
\end{adjustbox}
\caption{Step-by-Step SFT: each numbered step and the final answer contribute to the total loss.}
\end{figure}

\subsection{State Transition Table}
\begin{table}[h]\centering
\begin{tabular}{lll}\toprule
\textbf{Property} & \textbf{Before} & \textbf{After}\\\midrule
Solution structure & Unstructured & Numbered procedural steps\\
Step delimiter tokens & Not in vocabulary & Learned as first-class tokens\\
Interpretability & Low & Step-level traceability\\
Annotation format & Raw answer & $(x, [s_1,...,s_N], y)$\\
Task decomposition ability & Implicit & Explicit\\\bottomrule
\end{tabular}
\caption{State properties before and after Step-by-Step SFT.}
\end{table}

\newpage
%% ------------------------------------------------------------
\section{Process Supervision}

\subsection{Mathematical Foundation}
Process Supervision trains a Process Reward Model (PRM) using step-level binary correctness labels. Let $x$ be the problem description, and $s_1, \dots, s_S$ be the generated steps. Human or model-based annotations assign correctness labels $r_k \in \{0, 1\}$ for each step $s_k$.
The PRM model outputs correctness probabilities by applying a logistic sigmoid function to the step boundary token representation:
\begin{equation}
  \hat{r}_k = \sigma\left( \text{PRM}_\phi(h_{t_k}) \right) = \frac{1}{1 + \exp\left( -\text{PRM}_\phi(h_{t_k}) \right)}
\end{equation}
where $h_{t_k}$ is the transformer hidden output at index $t_k$ representing the step boundary token. The PRM parameters $\phi$ are optimized using binary cross entropy:
\begin{equation}
  \mathcal{L}_{\text{PRM}}(\phi) = -\sum_{k=1}^S \left[ r_k \log \hat{r}_k + (1-r_k) \log(1 - \hat{r}_k) \right]
\end{equation}
During inference, candidate sequences are generated. The total sequence correctness score is calculated as the product of step correctness probabilities:
\begin{equation}
  \text{Score}(s_{\le S}) = \prod_{k=1}^S \hat{r}_k
\end{equation}
This score is used to rank candidate rollouts in Best-of-$N$ sampling or guide selection during Monte Carlo Tree Search (MCTS).

\subsection{Architecture Flow Diagram}

\begin{figure}[h]
\centering
\begin{adjustbox}{max width=\linewidth}
\begin{tikzpicture}[node distance=0.6cm and 1.0cm]
  \node[datblk] (sol) {\textbf{Solution} \\ (Steps $s_{1..N}$)};
  \node[grayblk, right=of sol] (prm) {\textbf{Process Reward} \\ (Annotator)};
  \node[trnblk, right=0.6cm of prm, minimum height=2.2cm, minimum width=2.2cm] (labels) {\textbf{Step} \\ \textbf{Labels} \\ ($r_1,..,r_N$)};
  \node[lossblk, right=1.6cm of labels] (loss) {\textbf{Weighted} \\ (SFT Loss)};
  \node[mrgblk, below=0.9cm of loss] (out) {\textbf{Process-Supervised} \\ (Model)};

  \node[green!60!black,draw,rounded corners,font=\tiny,minimum width=1.5cm] at ($(labels.east)+(0.8,0.5)$) {$r_n{=}{+}1$: train};
  \node[red!70!black,draw,rounded corners,font=\tiny,minimum width=1.5cm] at ($(labels.east)+(0.8,-0.5)$) {$r_n{=}{-}1$: skip};

  \draw[arr] (sol)--(prm);
  \draw[arr] (prm)--(labels);
  \draw[arr] (labels)--(loss);
  \draw[arr] (loss)--(out);
\end{tikzpicture}
\end{adjustbox}
\caption{Process supervision: a step annotator labels each reasoning step; only correct steps contribute to the SFT loss.}
\end{figure}

\subsection{State Transition Table}
\begin{table}[h]\centering
\begin{tabular}{lll}\toprule
\textbf{Property} & \textbf{Before} & \textbf{After}\\\midrule
Supervision granularity & Final answer only & Per-step labels\\
Error signal & Outcome-level & Step-level\\
Error localisation & None & Identifies which step failed\\
Annotation cost & Low & High (step-level labelling)\\
Reasoning quality & Outcome-driven & Step-correctness driven\\\bottomrule
\end{tabular}
\caption{State properties before and after process supervision.}
\end{table}

\newpage
%% ============================================================
\part{Parameter-Efficient Fine-Tuning (PEFT)}
%% ============================================================

\section{Low-Rank Adaptation (LoRA)}

\subsection{Mathematical Foundation}
LoRA represents the weight update matrix $\Delta W \in \mathbb{R}^{k \times d}$ for a linear projection layer as the product of two low-rank matrices $B \in \mathbb{R}^{k \times r}$ and $A \in \mathbb{R}^{r \times d}$, where the rank satisfies $r \ll \min(d, k)$:
\begin{equation}
  h = W_0 x + \Delta W x = W_0 x + \frac{\alpha}{r} B A x
\end{equation}
where $\alpha$ is a constant scaling parameter. The base weight matrix $W_0 \in \mathbb{R}^{k \times d}$ is frozen; only $B$ and $A$ are updated. Under a loss $\mathcal{L}$, the gradient propagation rules using the chain rule are derived as follows:
Let $z = A x \in \mathbb{R}^r$ be the intermediate bottleneck activation. Then $h = W_0 x + \frac{\alpha}{r} B z$.
The gradient with respect to output activation is $\nabla_h \mathcal{L} \in \mathbb{R}^{k \times 1}$ (represented as a column vector).
The gradient with respect to the intermediate bottleneck representation $z$:
\begin{equation}
  \nabla_z \mathcal{L} = \frac{\alpha}{r} B^\top \nabla_h \mathcal{L} \in \mathbb{R}^{r \times 1}
\end{equation}
The gradient with respect to the matrix $B$:
\begin{equation}
  \nabla_B \mathcal{L} = \frac{\alpha}{r} \nabla_h \mathcal{L} z^\top = \frac{\alpha}{r} \nabla_h \mathcal{L} x^\top A^\top \in \mathbb{R}^{k \times r}
\end{equation}
The gradient with respect to the matrix $A$:
\begin{equation}
  \nabla_A \mathcal{L} = \nabla_z \mathcal{L} x^\top = \frac{\alpha}{r} B^\top \nabla_h \mathcal{L} x^\top \in \mathbb{R}^{r \times d}
\end{equation}
To eliminate latency at deployment, the weights are merged directly into the base model:
\begin{equation}
  W_{\text{merged}} = W_0 + \frac{\alpha}{r} B A
\end{equation}

\subsection{Architecture Flow Diagram}

\begin{figure}[h]
\centering
\begin{adjustbox}{max width=\linewidth}
\begin{tikzpicture}[node distance=0.7cm and 1.1cm]
  \node[neuron] (x1) at (0,1.5) {$x_1$};
  \node[neuron] (x2) at (0,0.5) {$x_2$};
  \node[neuron] (xd) at (0,-0.8) {$x_d$};
  \node[font=\tiny] at (0,-0.15) {$\vdots$};

  \node[frzblk, minimum width=2.8cm, minimum height=2.0cm] (W0) at (3.8,1.5) {\textbf{Frozen} \\ ($W_0 \in \mathbb{R}^{k\times d}$)};
  \node[bneck] (z1) at (3.8,-0.2) {$z_1$};
  \node[bneck] (zr) at (3.8,-1.2) {$z_r$};
  \node[font=\tiny] at (3.8,-0.7) {$\vdots$};

  \node[outnrn] (h1) at (7.6,1.5) {$h_1$};
  \node[outnrn] (h2) at (7.6,0.5) {$h_2$};
  \node[outnrn] (hk) at (7.6,-0.8) {$h_k$};
  \node[font=\tiny] at (7.6,-0.15) {$\vdots$};

  \draw[arr] (x1.east) -- ($(W0.west)+(0,0.5)$);
  \draw[arr] (x2.east) -- ($(W0.west)+(0,0)$);
  \draw[arr] (xd.east) -- ($(W0.west)+(0,-0.5)$);
  \foreach \s in {x1,x2,xd} \foreach \t in {z1,zr}
    \draw[arr,green!70!black,thin] (\s.east)--(\t.west);
  \foreach \s in {z1,zr} \foreach \t in {h1,h2,hk}
    \draw[arr,green!70!black,thin] (\s.east)--(\t.west);
  \draw[arr] ($(W0.east)+(0,0.5)$)--(h1.west);
  \draw[arr] ($(W0.east)+(0,0.0)$)--(h2.west);
  \draw[arr] ($(W0.east)+(0,-0.5)$)--(hk.west);

  \node[green!60!black,font=\tiny] at (1.9,-0.9) {$A$ (train)};
  \node[green!60!black,font=\tiny] at (5.7,-0.9) {$B$ (train)};
  \draw[garr] (h1.east)--++(0.6,0) node[right,red,font=\tiny]{$g_1$};
  \draw[garr] (hk.east)--++(0.6,0) node[right,red,font=\tiny]{$g_k$};
\end{tikzpicture}
\end{adjustbox}
\caption{LoRA: input flows through frozen $W_0$ and trainable low-rank path $A \to B$; outputs are summed.}
\end{figure}

\subsection{State Transition Table}
\begin{table}[h]\centering
\begin{tabular}{lll}\toprule
\textbf{Property} & \textbf{Before} & \textbf{After}\\\midrule
Base weights $W_0$ & requires\_grad = True & Frozen (False)\\
Trainable params & $k \times d$ & $r(k+d) \ll kd$\\
Matrix $B$ & Does not exist & $\mathbb{R}^{k\times r}$, init $= 0$\\
Matrix $A$ & Does not exist & $\mathbb{R}^{r\times d}$, init $\sim \mathcal{N}$\\
Inference overhead & None & Extra BA multiply (or merged)\\\bottomrule
\end{tabular}
\caption{State properties before and after LoRA injection.}
\end{table}

\subsection{Memory Complexity Analysis}
In this section, we analyze the memory footprint of the optimizer states. We compare full fine-tuning of a base weight matrix $W_0 \in \mathbb{R}^{k \times d}$ against low-rank training of $A$ and $B$. Let $N_{\text{base}} = k \cdot d$ be the number of base parameters, and $N_{\text{LoRA}} = r(k + d)$ be the number of LoRA parameters.

AdamW tracks the following state elements for each trainable parameter:
\begin{enumerate}
  \item \textbf{Parameter value copy} (master weights in mixed-precision): $4$ bytes.
  \item \textbf{First moment vector $m$ (momentum)}: $4$ bytes.
  \item \textbf{Second moment vector $v$ (variance)}: $4$ bytes.
\end{enumerate}
This results in $M_{\text{opt}} = 12$ bytes of optimizer state memory overhead per trainable parameter.

\noindent The memory saving ratio $\Phi$ is defined as:
\begin{equation}
  \Phi = \frac{M_{\text{LoRA}}}{M_{\text{Full}}} = \frac{r(k+d)}{k \cdot d} = r \left( \frac{1}{k} + \frac{1}{d} \right)
\end{equation}
Assuming a standard linear layer in a transformer block where $k = d = 4096$, and adapter rank $r = 8$:
\begin{equation}
  \Phi = 8 \left( \frac{1}{4096} + \frac{1}{4096} \right) = \frac{16}{4096} = \frac{1}{256} \approx 0.39\%
\end{equation}
This indicates a \textbf{256-fold reduction} (a decrease of $99.61\%$) in the memory required to store optimizer states for this weight matrix.



\subsection{Detailed Operational Analysis of LoRA Variables}
To understand the execution of LoRA at a granular level, we detail the mathematical nature of each variable, how its values are derived in reality, its child-level control dial mental model, and its physical mechanism at the individual neuron level. We provide a full neuron-level highlight diagram for each variable.

\begin{enumerate}
  \item \textbf{Input Activation Vector ($x$)}:
    \begin{itemize}
      \item \textbf{Formula}: $x$ (output of previous layer)
      \item \textbf{Neuron Highlight Diagram}:

\begin{center}
\begin{adjustbox}{max width=0.85\linewidth}
\begin{tikzpicture}[node distance=1.0cm and 1.5cm, scale=0.85, every node/.style={transform shape}]
  \useasboundingbox (-1.8, -3.5) rectangle (12.8, 3.8);
  % Input Neurons
  \node[neuron] (x1) at (0, 2) {$x_1$};
  \node[neuron] (x2) at (0, 0.8) {$x_2$};
  \node[neuron] (xk) at (0, -0.8) {$x_d$};
  \node[font=\tiny] at (0, 0.1) {$\vdots$};
  \node[lbl, above=0.2cm of x1] {Input $x$};

  % Base Weight Matrix W_0
  \node[frzblk, minimum width=3.2cm, minimum height=1.5cm] (W0) at (3.5, 2.5) {\textbf{Frozen Base} \\ $W_0 \in \mathbb{R}^{k \times d}$};

  % Bottleneck Neurons
  \node[bneck] (z1) at (3.5, -0.3) {$z_1$};
  \node[bneck] (zr) at (3.5, -1.5) {$z_r$};
  \node[font=\tiny] at (3.5, -0.9) {$\vdots$};
  \node[lbl, above=0.2cm of z1] {Bottleneck $z$};

  % Connections A
  \foreach \s in {x1,x2,xk} \foreach \t in {z1,zr}
    \draw[arr, gray!60, thin] (\s.east)--(\t.west);
  \node[text=gray, font=\footnotesize] at (1.5, -0.25) {$A$};

  % Summing Junctions and Output Neurons
  \node[sumnode] (sum1) at (7.0, 2) {$+$};
  \node[sumnode] (sumd) at (7.0, 0.2) {$+$};
  \node[font=\tiny] at (7.0, 1.2) {$\vdots$};
  
  \node[outnrn] (h1) at (7.8, 2) {$h_1$};
  \node[outnrn] (hd) at (7.8, 0.2) {$h_k$};
  \node[font=\tiny] at (7.8, 1.2) {$\vdots$};
  \node[lbl, above=0.2cm of h1] {Output $h$};

  % Connections B
  \foreach \s in {z1,zr} \foreach \t in {sum1,sumd}
    \draw[arr, gray!60, thin] (\s.east)--(\t.west);
  \node[text=gray, font=\footnotesize] at (5.3, -0.55) {$B$};

  % Scaling Multiplier
  \node[draw, regular polygon, regular polygon sides=3, rotate=-90, fill=gray!10, minimum size=0.9cm, draw=gray!60] (scale) at (5.3, 0.8) {$\frac{\alpha}{r}$};

  % Forward path connections
  \draw[arr, gray!60] (x1.east) -- ($(W0.west)+(0,0.5)$);
  \draw[arr, gray!60] (xk.east) -- ($(W0.west)+(0,-0.5)$);
  \draw[arr, gray!60] ($(W0.east)+(0,0.5)$) -- (sum1.west);
  \draw[arr, gray!60] ($(W0.east)+(0,-0.5)$) -- (sumd.west);
  \draw[arr, gray!60] (sum1) -- (h1);
  \draw[arr, gray!60] (sumd) -- (hd);

  % Loss and Target
  \node[datblk] (y) at (11.2, 0.2) {Target $y$};
  \node[lossblk] (loss) at (11.2, 2.0) {Loss $\mathcal{L}$};
  \draw[arr, gray!60] (h1.east) -- (loss.west);
  \draw[arr, gray!60] (y.north) -- (loss.south);
  
  % Dashed base outlines fit
  \node[fit=(x1) (x2) (xk)] (fit_x) {};
  \node[fit=(W0)] (fit_w0) {};
  \node[fit=(z1) (zr)] (fit_z) {};
  \node[fit=(sum1) (sumd) (h1) (hd)] (fit_h) {};

  \node[draw=orange, dashed, rounded corners=4pt, fit=(x1) (x2) (xk), inner sep=5pt, line width=1.5pt] (circle_x) {};
  \node[text=orange, font=\bfseries\small, left=0.2cm of circle_x, align=center] {HIGHLIGHT:\\Input $x$};

\end{tikzpicture}
\end{adjustbox}
\vspace{1.5em}
\end{center}

      \item \textbf{Mental Model}: Think of this vector as a list of positions for a row of thousands of physical dials. When the model processes a word, it adjusts these dials to specific numbers (e.g., dial 1 is at 4.2, dial 2 is at -1.5) to capture the meaning of the word. The model does not think in words; it thinks in these thousands of dial settings.
      \item \textbf{Neuron-Level Physical Explanation}: A single neuron in the current layer receives $x$ as a set of electrical signals from the neurons of the previous layer. Each element $x_i$ represents the activation level (firing rate) of a previous neuron after its activation function (such as ReLU or GELU). The dendrites of the current neuron receive these signals.
      \item \textbf{Mathematical Representation}: $x \in \mathbb{R}^{d \times 1}$
      \item \textbf{Code Representation}: \texttt{x} (a PyTorch tensor of shape \texttt{(batch\_size, seq\_len, d)})
      \item \textbf{Relationships}: Fed into the base projection $W_0 x$ and compressed to $z = Ax$.
    \end{itemize}

  \item \textbf{Frozen Base Weight Matrix ($W_0$)}:
    \begin{itemize}
      \item \textbf{Formula}: $W_0$
      \item \textbf{Neuron Highlight Diagram}:

\begin{center}
\begin{adjustbox}{max width=0.85\linewidth}
\begin{tikzpicture}[node distance=1.0cm and 1.5cm, scale=0.85, every node/.style={transform shape}]
  \useasboundingbox (-1.8, -3.5) rectangle (12.8, 3.8);
  % Input Neurons
  \node[neuron] (x1) at (0, 2) {$x_1$};
  \node[neuron] (x2) at (0, 0.8) {$x_2$};
  \node[neuron] (xk) at (0, -0.8) {$x_d$};
  \node[font=\tiny] at (0, 0.1) {$\vdots$};
  \node[lbl, above=0.2cm of x1] {Input $x$};

  % Base Weight Matrix W_0
  \node[frzblk, minimum width=3.2cm, minimum height=1.5cm] (W0) at (3.5, 2.5) {\textbf{Frozen Base} \\ $W_0 \in \mathbb{R}^{k \times d}$};

  % Bottleneck Neurons
  \node[bneck] (z1) at (3.5, -0.3) {$z_1$};
  \node[bneck] (zr) at (3.5, -1.5) {$z_r$};
  \node[font=\tiny] at (3.5, -0.9) {$\vdots$};
  \node[lbl, above=0.2cm of z1] {Bottleneck $z$};

  % Connections A
  \foreach \s in {x1,x2,xk} \foreach \t in {z1,zr}
    \draw[arr, gray!60, thin] (\s.east)--(\t.west);
  \node[text=gray, font=\footnotesize] at (1.5, -0.25) {$A$};

  % Summing Junctions and Output Neurons
  \node[sumnode] (sum1) at (7.0, 2) {$+$};
  \node[sumnode] (sumd) at (7.0, 0.2) {$+$};
  \node[font=\tiny] at (7.0, 1.2) {$\vdots$};
  
  \node[outnrn] (h1) at (7.8, 2) {$h_1$};
  \node[outnrn] (hd) at (7.8, 0.2) {$h_k$};
  \node[font=\tiny] at (7.8, 1.2) {$\vdots$};
  \node[lbl, above=0.2cm of h1] {Output $h$};

  % Connections B
  \foreach \s in {z1,zr} \foreach \t in {sum1,sumd}
    \draw[arr, gray!60, thin] (\s.east)--(\t.west);
  \node[text=gray, font=\footnotesize] at (5.3, -0.55) {$B$};

  % Scaling Multiplier
  \node[draw, regular polygon, regular polygon sides=3, rotate=-90, fill=gray!10, minimum size=0.9cm, draw=gray!60] (scale) at (5.3, 0.8) {$\frac{\alpha}{r}$};

  % Forward path connections
  \draw[arr, gray!60] (x1.east) -- ($(W0.west)+(0,0.5)$);
  \draw[arr, gray!60] (xk.east) -- ($(W0.west)+(0,-0.5)$);
  \draw[arr, gray!60] ($(W0.east)+(0,0.5)$) -- (sum1.west);
  \draw[arr, gray!60] ($(W0.east)+(0,-0.5)$) -- (sumd.west);
  \draw[arr, gray!60] (sum1) -- (h1);
  \draw[arr, gray!60] (sumd) -- (hd);

  % Loss and Target
  \node[datblk] (y) at (11.2, 0.2) {Target $y$};
  \node[lossblk] (loss) at (11.2, 2.0) {Loss $\mathcal{L}$};
  \draw[arr, gray!60] (h1.east) -- (loss.west);
  \draw[arr, gray!60] (y.north) -- (loss.south);
  
  % Dashed base outlines fit
  \node[fit=(x1) (x2) (xk)] (fit_x) {};
  \node[fit=(W0)] (fit_w0) {};
  \node[fit=(z1) (zr)] (fit_z) {};
  \node[fit=(sum1) (sumd) (h1) (hd)] (fit_h) {};

  \node[draw=blue, dashed, rounded corners=4pt, fit=(W0), inner sep=3pt, line width=1.5pt] (circle_w0) {};
  \node[text=blue, font=\bfseries\small, above=0.15cm of circle_w0] {HIGHLIGHT: Base Weight $W_0$};

\end{tikzpicture}
\end{adjustbox}
\vspace{1.5em}
\end{center}

      \item \textbf{Mental Model}: This is a giant, permanent circuit board inside the control room. The wires are soldered in place and cannot be moved (frozen). They take the positions of the incoming dials ($x$), mix them together according to the permanent rules, and output a new set of general dial positions.
      \item \textbf{Neuron-Level Physical Explanation}: A single neuron's base connections are represented by a row of $W_0$. The neuron performs a dot product: it multiplies each incoming signal $x_i$ by the strength of its corresponding synaptic connection $W_{0, ji}$ and sums them up: $y_j = \sum_i W_{0, ji} x_i$. During LoRA, these synaptic weights are frozen and cannot change, meaning their chemical/electrical conductance is locked.
      \item \textbf{Mathematical Representation}: $W_0 \in \mathbb{R}^{k \times d}$
      \item \textbf{Code Representation}: \texttt{W\_0} (a PyTorch tensor with \texttt{W\_0.requires\_grad = False} to disable gradient computation)
      \item \textbf{Relationships}: Multiplied by the input to compute the baseline representation: $h_{\text{base}} = W_0 x$.
    \end{itemize}

  \item \textbf{Down-projection Adapter ($A$)}:
    \begin{itemize}
      \item \textbf{Formula}: $A$
      \item \textbf{Neuron Highlight Diagram}:

\begin{center}
\begin{adjustbox}{max width=0.85\linewidth}
\begin{tikzpicture}[node distance=1.0cm and 1.5cm, scale=0.85, every node/.style={transform shape}]
  \useasboundingbox (-1.8, -3.5) rectangle (12.8, 3.8);
  % Input Neurons
  \node[neuron] (x1) at (0, 2) {$x_1$};
  \node[neuron] (x2) at (0, 0.8) {$x_2$};
  \node[neuron] (xk) at (0, -0.8) {$x_d$};
  \node[font=\tiny] at (0, 0.1) {$\vdots$};
  \node[lbl, above=0.2cm of x1] {Input $x$};

  % Base Weight Matrix W_0
  \node[frzblk, minimum width=3.2cm, minimum height=1.5cm] (W0) at (3.5, 2.5) {\textbf{Frozen Base} \\ $W_0 \in \mathbb{R}^{k \times d}$};

  % Bottleneck Neurons
  \node[bneck] (z1) at (3.5, -0.3) {$z_1$};
  \node[bneck] (zr) at (3.5, -1.5) {$z_r$};
  \node[font=\tiny] at (3.5, -0.9) {$\vdots$};
  \node[lbl, above=0.2cm of z1] {Bottleneck $z$};

  % Connections A
  \foreach \s in {x1,x2,xk} \foreach \t in {z1,zr}
    \draw[arr, gray!60, thin] (\s.east)--(\t.west);
  \node[text=gray, font=\footnotesize] at (1.5, -0.25) {$A$};

  % Summing Junctions and Output Neurons
  \node[sumnode] (sum1) at (7.0, 2) {$+$};
  \node[sumnode] (sumd) at (7.0, 0.2) {$+$};
  \node[font=\tiny] at (7.0, 1.2) {$\vdots$};
  
  \node[outnrn] (h1) at (7.8, 2) {$h_1$};
  \node[outnrn] (hd) at (7.8, 0.2) {$h_k$};
  \node[font=\tiny] at (7.8, 1.2) {$\vdots$};
  \node[lbl, above=0.2cm of h1] {Output $h$};

  % Connections B
  \foreach \s in {z1,zr} \foreach \t in {sum1,sumd}
    \draw[arr, gray!60, thin] (\s.east)--(\t.west);
  \node[text=gray, font=\footnotesize] at (5.3, -0.55) {$B$};

  % Scaling Multiplier
  \node[draw, regular polygon, regular polygon sides=3, rotate=-90, fill=gray!10, minimum size=0.9cm, draw=gray!60] (scale) at (5.3, 0.8) {$\frac{\alpha}{r}$};

  % Forward path connections
  \draw[arr, gray!60] (x1.east) -- ($(W0.west)+(0,0.5)$);
  \draw[arr, gray!60] (xk.east) -- ($(W0.west)+(0,-0.5)$);
  \draw[arr, gray!60] ($(W0.east)+(0,0.5)$) -- (sum1.west);
  \draw[arr, gray!60] ($(W0.east)+(0,-0.5)$) -- (sumd.west);
  \draw[arr, gray!60] (sum1) -- (h1);
  \draw[arr, gray!60] (sumd) -- (hd);

  % Loss and Target
  \node[datblk] (y) at (11.2, 0.2) {Target $y$};
  \node[lossblk] (loss) at (11.2, 2.0) {Loss $\mathcal{L}$};
  \draw[arr, gray!60] (h1.east) -- (loss.west);
  \draw[arr, gray!60] (y.north) -- (loss.south);
  
  % Dashed base outlines fit
  \node[fit=(x1) (x2) (xk)] (fit_x) {};
  \node[fit=(W0)] (fit_w0) {};
  \node[fit=(z1) (zr)] (fit_z) {};
  \node[fit=(sum1) (sumd) (h1) (hd)] (fit_h) {};

  \foreach \s in {x1,x2,xk} \foreach \t in {z1,zr}
    \draw[arr, green!80!black, line width=1.5pt] (\s.east)--(\t.west);
  \node[text=green!80!black, font=\bfseries\small, draw=green!80!black, dashed, inner sep=3pt, rounded corners=2pt, fill=white] (lbl_A) at (1.8, -2.5) {HIGHLIGHT: Down-proj $A$};
  \draw[dashed, green!80!black, line width=1pt, ->] (lbl_A.north) -- (1.5, -0.5);

\end{tikzpicture}
\end{adjustbox}
\vspace{1.5em}
\end{center}

      \item \textbf{Mental Model}: This is a smaller circuit board that acts as a compression funnel. It looks at all the thousands of input dials and summarizes them into just a few dials (e.g., summarizing $4096$ dials into just $8$ dials). It only listens to the dial changes that are relevant to the new task we are teaching the model.
      \item \textbf{Neuron-Level Physical Explanation}: Inside the adapter path, a single neuron in the bottleneck layer corresponds to a row of matrix $A$. This neuron listens to all input signals $x_i$, multiplies them by its trainable weights $A_{mi}$, and sums them up: $z_m = \sum_i A_{mi} x_i$. This bottleneck neuron acts as a feature detector that extracts a specific combination of input signals.
      \item \textbf{Mathematical Representation}: $A \in \mathbb{R}^{r \times d}$
      \item \textbf{Code Representation}: \texttt{A} (initialized as \texttt{nn.Linear(}d, r, bias=False) with weight values drawn from a Gaussian distribution: \texttt{nn.init.kaiming\_uniform\_(A.weight, a=math.sqrt(5))})
      \item \textbf{Relationships}: Compresses $x$ into the bottleneck space via $z = Ax$.
    \end{itemize}

  \item \textbf{Up-projection Adapter ($B$)}:
    \begin{itemize}
      \item \textbf{Formula}: $B$
      \item \textbf{Neuron Highlight Diagram}:

\begin{center}
\begin{adjustbox}{max width=0.85\linewidth}
\begin{tikzpicture}[node distance=1.0cm and 1.5cm, scale=0.85, every node/.style={transform shape}]
  \useasboundingbox (-1.8, -3.5) rectangle (12.8, 3.8);
  % Input Neurons
  \node[neuron] (x1) at (0, 2) {$x_1$};
  \node[neuron] (x2) at (0, 0.8) {$x_2$};
  \node[neuron] (xk) at (0, -0.8) {$x_d$};
  \node[font=\tiny] at (0, 0.1) {$\vdots$};
  \node[lbl, above=0.2cm of x1] {Input $x$};

  % Base Weight Matrix W_0
  \node[frzblk, minimum width=3.2cm, minimum height=1.5cm] (W0) at (3.5, 2.5) {\textbf{Frozen Base} \\ $W_0 \in \mathbb{R}^{k \times d}$};

  % Bottleneck Neurons
  \node[bneck] (z1) at (3.5, -0.3) {$z_1$};
  \node[bneck] (zr) at (3.5, -1.5) {$z_r$};
  \node[font=\tiny] at (3.5, -0.9) {$\vdots$};
  \node[lbl, above=0.2cm of z1] {Bottleneck $z$};

  % Connections A
  \foreach \s in {x1,x2,xk} \foreach \t in {z1,zr}
    \draw[arr, gray!60, thin] (\s.east)--(\t.west);
  \node[text=gray, font=\footnotesize] at (1.5, -0.25) {$A$};

  % Summing Junctions and Output Neurons
  \node[sumnode] (sum1) at (7.0, 2) {$+$};
  \node[sumnode] (sumd) at (7.0, 0.2) {$+$};
  \node[font=\tiny] at (7.0, 1.2) {$\vdots$};
  
  \node[outnrn] (h1) at (7.8, 2) {$h_1$};
  \node[outnrn] (hd) at (7.8, 0.2) {$h_k$};
  \node[font=\tiny] at (7.8, 1.2) {$\vdots$};
  \node[lbl, above=0.2cm of h1] {Output $h$};

  % Connections B
  \foreach \s in {z1,zr} \foreach \t in {sum1,sumd}
    \draw[arr, gray!60, thin] (\s.east)--(\t.west);
  \node[text=gray, font=\footnotesize] at (5.3, -0.55) {$B$};

  % Scaling Multiplier
  \node[draw, regular polygon, regular polygon sides=3, rotate=-90, fill=gray!10, minimum size=0.9cm, draw=gray!60] (scale) at (5.3, 0.8) {$\frac{\alpha}{r}$};

  % Forward path connections
  \draw[arr, gray!60] (x1.east) -- ($(W0.west)+(0,0.5)$);
  \draw[arr, gray!60] (xk.east) -- ($(W0.west)+(0,-0.5)$);
  \draw[arr, gray!60] ($(W0.east)+(0,0.5)$) -- (sum1.west);
  \draw[arr, gray!60] ($(W0.east)+(0,-0.5)$) -- (sumd.west);
  \draw[arr, gray!60] (sum1) -- (h1);
  \draw[arr, gray!60] (sumd) -- (hd);

  % Loss and Target
  \node[datblk] (y) at (11.2, 0.2) {Target $y$};
  \node[lossblk] (loss) at (11.2, 2.0) {Loss $\mathcal{L}$};
  \draw[arr, gray!60] (h1.east) -- (loss.west);
  \draw[arr, gray!60] (y.north) -- (loss.south);
  
  % Dashed base outlines fit
  \node[fit=(x1) (x2) (xk)] (fit_x) {};
  \node[fit=(W0)] (fit_w0) {};
  \node[fit=(z1) (zr)] (fit_z) {};
  \node[fit=(sum1) (sumd) (h1) (hd)] (fit_h) {};

  \foreach \s in {z1,zr} \foreach \t in {sum1,sumd}
    \draw[arr, green!80!black, line width=1.5pt] (\s.east)--(\t.west);
  \node[text=green!80!black, font=\bfseries\small, draw=green!80!black, dashed, inner sep=3pt, rounded corners=2pt, fill=white] (lbl_B) at (5.3, -2.5) {HIGHLIGHT: Up-proj $B$};
  \draw[dashed, green!80!black, line width=1pt, ->] (lbl_B.north) -- (5.3, -0.8);

\end{tikzpicture}
\end{adjustbox}
\vspace{1.5em}
\end{center}

      \item \textbf{Mental Model}: This is another small circuit board that acts as an expansion nozzle. It takes the tiny summary from the funnel (our $8$ dials) and expands it back into thousands of dials, showing how to nudge the final output dials to fit the new task.
      \item \textbf{Neuron-Level Physical Explanation}: A single output neuron receives contributions from the bottleneck neurons. The connection strength from bottleneck neuron $m$ to output neuron $j$ is the trainable weight $B_{jm}$. The output neuron sums the weighted signals from all bottleneck neurons: $u_j = \sum_m B_{jm} z_m$, scaling up the low-dimensional task features.
      \item \textbf{Mathematical Representation}: $B \in \mathbb{R}^{k \times r}$
      \item \textbf{Code Representation}: \texttt{B} (initialized as \texttt{nn.Linear(r, }k, bias=False) with weight values set to all zeros: \texttt{nn.init.zeros\_(B.weight)})
      \item \textbf{Relationships}: Expands the bottleneck activations via $B z$, scaled by $\frac{\alpha}{r}$.
    \end{itemize}

  \item \textbf{Bottleneck Rank ($r$)}:
    \begin{itemize}
      \item \textbf{Formula}: $r$
      \item \textbf{Neuron Highlight Diagram}:

\begin{center}
\begin{adjustbox}{max width=0.85\linewidth}
\begin{tikzpicture}[node distance=1.0cm and 1.5cm, scale=0.85, every node/.style={transform shape}]
  \useasboundingbox (-1.8, -3.5) rectangle (12.8, 3.8);
  % Input Neurons
  \node[neuron] (x1) at (0, 2) {$x_1$};
  \node[neuron] (x2) at (0, 0.8) {$x_2$};
  \node[neuron] (xk) at (0, -0.8) {$x_d$};
  \node[font=\tiny] at (0, 0.1) {$\vdots$};
  \node[lbl, above=0.2cm of x1] {Input $x$};

  % Base Weight Matrix W_0
  \node[frzblk, minimum width=3.2cm, minimum height=1.5cm] (W0) at (3.5, 2.5) {\textbf{Frozen Base} \\ $W_0 \in \mathbb{R}^{k \times d}$};

  % Bottleneck Neurons
  \node[bneck] (z1) at (3.5, -0.3) {$z_1$};
  \node[bneck] (zr) at (3.5, -1.5) {$z_r$};
  \node[font=\tiny] at (3.5, -0.9) {$\vdots$};
  \node[lbl, above=0.2cm of z1] {Bottleneck $z$};

  % Connections A
  \foreach \s in {x1,x2,xk} \foreach \t in {z1,zr}
    \draw[arr, gray!60, thin] (\s.east)--(\t.west);
  \node[text=gray, font=\footnotesize] at (1.5, -0.25) {$A$};

  % Summing Junctions and Output Neurons
  \node[sumnode] (sum1) at (7.0, 2) {$+$};
  \node[sumnode] (sumd) at (7.0, 0.2) {$+$};
  \node[font=\tiny] at (7.0, 1.2) {$\vdots$};
  
  \node[outnrn] (h1) at (7.8, 2) {$h_1$};
  \node[outnrn] (hd) at (7.8, 0.2) {$h_k$};
  \node[font=\tiny] at (7.8, 1.2) {$\vdots$};
  \node[lbl, above=0.2cm of h1] {Output $h$};

  % Connections B
  \foreach \s in {z1,zr} \foreach \t in {sum1,sumd}
    \draw[arr, gray!60, thin] (\s.east)--(\t.west);
  \node[text=gray, font=\footnotesize] at (5.3, -0.55) {$B$};

  % Scaling Multiplier
  \node[draw, regular polygon, regular polygon sides=3, rotate=-90, fill=gray!10, minimum size=0.9cm, draw=gray!60] (scale) at (5.3, 0.8) {$\frac{\alpha}{r}$};

  % Forward path connections
  \draw[arr, gray!60] (x1.east) -- ($(W0.west)+(0,0.5)$);
  \draw[arr, gray!60] (xk.east) -- ($(W0.west)+(0,-0.5)$);
  \draw[arr, gray!60] ($(W0.east)+(0,0.5)$) -- (sum1.west);
  \draw[arr, gray!60] ($(W0.east)+(0,-0.5)$) -- (sumd.west);
  \draw[arr, gray!60] (sum1) -- (h1);
  \draw[arr, gray!60] (sumd) -- (hd);

  % Loss and Target
  \node[datblk] (y) at (11.2, 0.2) {Target $y$};
  \node[lossblk] (loss) at (11.2, 2.0) {Loss $\mathcal{L}$};
  \draw[arr, gray!60] (h1.east) -- (loss.west);
  \draw[arr, gray!60] (y.north) -- (loss.south);
  
  % Dashed base outlines fit
  \node[fit=(x1) (x2) (xk)] (fit_x) {};
  \node[fit=(W0)] (fit_w0) {};
  \node[fit=(z1) (zr)] (fit_z) {};
  \node[fit=(sum1) (sumd) (h1) (hd)] (fit_h) {};

  \draw[decoration={brace,mirror,amplitude=6pt},decorate,line width=1.5pt,draw=purple] ($(zr.south west)+(-0.3,0)$) -- ($(z1.north west)+(-0.3,0)$);
  \node[text=purple, font=\bfseries\small] (lbl_r) at (3.5, -2.5) {HIGHLIGHT: Rank $r$};
  \draw[dashed, purple, line width=1pt, ->] (lbl_r.north) -- (3.2, -1.0);

\end{tikzpicture}
\end{adjustbox}
\vspace{1.5em}
\end{center}

      \item \textbf{Mental Model}: This is the exact number of summary dials inside the bottleneck. If $r=8$, it means we are forcing the model to summarize all its knowledge into just $8$ dials. A smaller number makes the funnel tighter, forcing the model to only keep the most important summaries.
      \item \textbf{Neuron-Level Physical Explanation}: This is the exact number of neurons present in the bottleneck layer of the adapter. Selecting $r=8$ means there are exactly 8 neurons in this intermediate layer. Each of these 8 neurons acts as a specialized task detector.
      \item \textbf{Mathematical Representation}: $r \in \mathbb{Z}^{+}$
      \item \textbf{Code Representation}: \texttt{r} (passed as an integer configuration parameter, e.g., \texttt{lora\_config = LoraConfig(r=8)})
      \item \textbf{Relationships}: Determines the inner dimension of matrices $A$ and $B$, and acts as a scaling denominator in $\frac{\alpha}{r}$.
    \end{itemize}

  \item \textbf{Scaling Constant ($\alpha$)}:
    \begin{itemize}
      \item \textbf{Formula}: $\alpha$
      \item \textbf{Neuron Highlight Diagram}:

\begin{center}
\begin{adjustbox}{max width=0.85\linewidth}
\begin{tikzpicture}[node distance=1.0cm and 1.5cm, scale=0.85, every node/.style={transform shape}]
  \useasboundingbox (-1.8, -3.5) rectangle (12.8, 3.8);
  % Input Neurons
  \node[neuron] (x1) at (0, 2) {$x_1$};
  \node[neuron] (x2) at (0, 0.8) {$x_2$};
  \node[neuron] (xk) at (0, -0.8) {$x_d$};
  \node[font=\tiny] at (0, 0.1) {$\vdots$};
  \node[lbl, above=0.2cm of x1] {Input $x$};

  % Base Weight Matrix W_0
  \node[frzblk, minimum width=3.2cm, minimum height=1.5cm] (W0) at (3.5, 2.5) {\textbf{Frozen Base} \\ $W_0 \in \mathbb{R}^{k \times d}$};

  % Bottleneck Neurons
  \node[bneck] (z1) at (3.5, -0.3) {$z_1$};
  \node[bneck] (zr) at (3.5, -1.5) {$z_r$};
  \node[font=\tiny] at (3.5, -0.9) {$\vdots$};
  \node[lbl, above=0.2cm of z1] {Bottleneck $z$};

  % Connections A
  \foreach \s in {x1,x2,xk} \foreach \t in {z1,zr}
    \draw[arr, gray!60, thin] (\s.east)--(\t.west);
  \node[text=gray, font=\footnotesize] at (1.5, -0.25) {$A$};

  % Summing Junctions and Output Neurons
  \node[sumnode] (sum1) at (7.0, 2) {$+$};
  \node[sumnode] (sumd) at (7.0, 0.2) {$+$};
  \node[font=\tiny] at (7.0, 1.2) {$\vdots$};
  
  \node[outnrn] (h1) at (7.8, 2) {$h_1$};
  \node[outnrn] (hd) at (7.8, 0.2) {$h_k$};
  \node[font=\tiny] at (7.8, 1.2) {$\vdots$};
  \node[lbl, above=0.2cm of h1] {Output $h$};

  % Connections B
  \foreach \s in {z1,zr} \foreach \t in {sum1,sumd}
    \draw[arr, gray!60, thin] (\s.east)--(\t.west);
  \node[text=gray, font=\footnotesize] at (5.3, -0.55) {$B$};

  % Scaling Multiplier
  \node[draw, regular polygon, regular polygon sides=3, rotate=-90, fill=gray!10, minimum size=0.9cm, draw=gray!60] (scale) at (5.3, 0.8) {$\frac{\alpha}{r}$};

  % Forward path connections
  \draw[arr, gray!60] (x1.east) -- ($(W0.west)+(0,0.5)$);
  \draw[arr, gray!60] (xk.east) -- ($(W0.west)+(0,-0.5)$);
  \draw[arr, gray!60] ($(W0.east)+(0,0.5)$) -- (sum1.west);
  \draw[arr, gray!60] ($(W0.east)+(0,-0.5)$) -- (sumd.west);
  \draw[arr, gray!60] (sum1) -- (h1);
  \draw[arr, gray!60] (sumd) -- (hd);

  % Loss and Target
  \node[datblk] (y) at (11.2, 0.2) {Target $y$};
  \node[lossblk] (loss) at (11.2, 2.0) {Loss $\mathcal{L}$};
  \draw[arr, gray!60] (h1.east) -- (loss.west);
  \draw[arr, gray!60] (y.north) -- (loss.south);
  
  % Dashed base outlines fit
  \node[fit=(x1) (x2) (xk)] (fit_x) {};
  \node[fit=(W0)] (fit_w0) {};
  \node[fit=(z1) (zr)] (fit_z) {};
  \node[fit=(sum1) (sumd) (h1) (hd)] (fit_h) {};

  \node[draw=purple, dashed, rounded corners=4pt, fit=(scale), inner sep=4pt, line width=1.5pt] (circle_scale) {};
  \node[text=purple, font=\bfseries\small] (lbl_scale) at (5.3, -2.5) {HIGHLIGHT: Scale $\frac{\alpha}{r}$};
  \draw[dashed, purple, line width=1pt, ->] (lbl_scale.north) -- (circle_scale.south);

\end{tikzpicture}
\end{adjustbox}
\vspace{1.5em}
\end{center}

      \item \textbf{Mental Model}: This is a volume control knob. It decides how loudly the changes from the new adapters ($A$ and $B$) should shout compared to the original, permanent circuit board ($W_0$). If we turn this knob up, the model listens more to its new training and less to its old pretrained knowledge.
      \item \textbf{Neuron-Level Physical Explanation}: At the neuron level, this acts as a synaptic gain control. The signal coming from the bottleneck via the adapter pathway is amplified or attenuated by a constant scale factor of $\frac{\alpha}{r}$ before it is added to the base neuron's output. This behaves like a biological neuromodulator that globally scales the synaptic sensitivity of the adapter path.
      \item \textbf{Mathematical Representation}: $\alpha \in \mathbb{R}^{+}$
      \item \textbf{Code Representation}: \texttt{lora\_alpha} (passed as a configuration parameter, e.g., \texttt{lora\_config = LoraConfig(lora\_alpha=16)})
      \item \textbf{Relationships}: Multiplies the adapter pathway output, yielding $\frac{\alpha}{r} B z$.
    \end{itemize}

  \item \textbf{Bottleneck Activation Vector ($z$)}:
    \begin{itemize}
      \item \textbf{Formula}: $z = Ax$
      \item \textbf{Neuron Highlight Diagram}:

\begin{center}
\begin{adjustbox}{max width=0.85\linewidth}
\begin{tikzpicture}[node distance=1.0cm and 1.5cm, scale=0.85, every node/.style={transform shape}]
  \useasboundingbox (-1.8, -3.5) rectangle (12.8, 3.8);
  % Input Neurons
  \node[neuron] (x1) at (0, 2) {$x_1$};
  \node[neuron] (x2) at (0, 0.8) {$x_2$};
  \node[neuron] (xk) at (0, -0.8) {$x_d$};
  \node[font=\tiny] at (0, 0.1) {$\vdots$};
  \node[lbl, above=0.2cm of x1] {Input $x$};

  % Base Weight Matrix W_0
  \node[frzblk, minimum width=3.2cm, minimum height=1.5cm] (W0) at (3.5, 2.5) {\textbf{Frozen Base} \\ $W_0 \in \mathbb{R}^{k \times d}$};

  % Bottleneck Neurons
  \node[bneck] (z1) at (3.5, -0.3) {$z_1$};
  \node[bneck] (zr) at (3.5, -1.5) {$z_r$};
  \node[font=\tiny] at (3.5, -0.9) {$\vdots$};
  \node[lbl, above=0.2cm of z1] {Bottleneck $z$};

  % Connections A
  \foreach \s in {x1,x2,xk} \foreach \t in {z1,zr}
    \draw[arr, gray!60, thin] (\s.east)--(\t.west);
  \node[text=gray, font=\footnotesize] at (1.5, -0.25) {$A$};

  % Summing Junctions and Output Neurons
  \node[sumnode] (sum1) at (7.0, 2) {$+$};
  \node[sumnode] (sumd) at (7.0, 0.2) {$+$};
  \node[font=\tiny] at (7.0, 1.2) {$\vdots$};
  
  \node[outnrn] (h1) at (7.8, 2) {$h_1$};
  \node[outnrn] (hd) at (7.8, 0.2) {$h_k$};
  \node[font=\tiny] at (7.8, 1.2) {$\vdots$};
  \node[lbl, above=0.2cm of h1] {Output $h$};

  % Connections B
  \foreach \s in {z1,zr} \foreach \t in {sum1,sumd}
    \draw[arr, gray!60, thin] (\s.east)--(\t.west);
  \node[text=gray, font=\footnotesize] at (5.3, -0.55) {$B$};

  % Scaling Multiplier
  \node[draw, regular polygon, regular polygon sides=3, rotate=-90, fill=gray!10, minimum size=0.9cm, draw=gray!60] (scale) at (5.3, 0.8) {$\frac{\alpha}{r}$};

  % Forward path connections
  \draw[arr, gray!60] (x1.east) -- ($(W0.west)+(0,0.5)$);
  \draw[arr, gray!60] (xk.east) -- ($(W0.west)+(0,-0.5)$);
  \draw[arr, gray!60] ($(W0.east)+(0,0.5)$) -- (sum1.west);
  \draw[arr, gray!60] ($(W0.east)+(0,-0.5)$) -- (sumd.west);
  \draw[arr, gray!60] (sum1) -- (h1);
  \draw[arr, gray!60] (sumd) -- (hd);

  % Loss and Target
  \node[datblk] (y) at (11.2, 0.2) {Target $y$};
  \node[lossblk] (loss) at (11.2, 2.0) {Loss $\mathcal{L}$};
  \draw[arr, gray!60] (h1.east) -- (loss.west);
  \draw[arr, gray!60] (y.north) -- (loss.south);
  
  % Dashed base outlines fit
  \node[fit=(x1) (x2) (xk)] (fit_x) {};
  \node[fit=(W0)] (fit_w0) {};
  \node[fit=(z1) (zr)] (fit_z) {};
  \node[fit=(sum1) (sumd) (h1) (hd)] (fit_h) {};

  \node[draw=teal!80!black, dashed, rounded corners=4pt, fit=(z1) (zr), inner sep=4pt, line width=1.5pt] (circle_z) {};
  \node[text=teal!80!black, font=\bfseries\small] (lbl_z) at (3.5, -2.5) {HIGHLIGHT: Bottleneck $z$};
  \draw[dashed, teal!80!black, line width=1pt, ->] (lbl_z.north) -- (circle_z.south);

\end{tikzpicture}
\end{adjustbox}
\vspace{1.5em}
\end{center}

      \item \textbf{Mental Model}: These are the actual values of our few summary dials inside the bottleneck. It represents the compressed task-specific message as it flows from the compression funnel to the expansion nozzle.
      \item \textbf{Neuron-Level Physical Explanation}: Each element $z_m$ is the electrical activation level (firing rate) of a single neuron in the bottleneck layer. It is calculated by that neuron summing its inputs: $z_m = \sum_i A_{mi} x_i$.
      \item \textbf{Mathematical Representation}: $z \in \mathbb{R}^{r \times 1}$
      \item \textbf{Code Representation}: \texttt{z = F.linear(x, A.weight)} (or \texttt{z = x @ A.weight.t()})
      \item \textbf{Relationships}: Created by $A$, consumed by $B$.
    \end{itemize}

  \item \textbf{Combined Output Vector ($h$)}:
    \begin{itemize}
      \item \textbf{Formula}: $h = W_0 x + \frac{\alpha}{r} B z$
      \item \textbf{Neuron Highlight Diagram}:

\begin{center}
\begin{adjustbox}{max width=0.85\linewidth}
\begin{tikzpicture}[node distance=1.0cm and 1.5cm, scale=0.85, every node/.style={transform shape}]
  \useasboundingbox (-1.8, -3.5) rectangle (12.8, 3.8);
  % Input Neurons
  \node[neuron] (x1) at (0, 2) {$x_1$};
  \node[neuron] (x2) at (0, 0.8) {$x_2$};
  \node[neuron] (xk) at (0, -0.8) {$x_d$};
  \node[font=\tiny] at (0, 0.1) {$\vdots$};
  \node[lbl, above=0.2cm of x1] {Input $x$};

  % Base Weight Matrix W_0
  \node[frzblk, minimum width=3.2cm, minimum height=1.5cm] (W0) at (3.5, 2.5) {\textbf{Frozen Base} \\ $W_0 \in \mathbb{R}^{k \times d}$};

  % Bottleneck Neurons
  \node[bneck] (z1) at (3.5, -0.3) {$z_1$};
  \node[bneck] (zr) at (3.5, -1.5) {$z_r$};
  \node[font=\tiny] at (3.5, -0.9) {$\vdots$};
  \node[lbl, above=0.2cm of z1] {Bottleneck $z$};

  % Connections A
  \foreach \s in {x1,x2,xk} \foreach \t in {z1,zr}
    \draw[arr, gray!60, thin] (\s.east)--(\t.west);
  \node[text=gray, font=\footnotesize] at (1.5, -0.25) {$A$};

  % Summing Junctions and Output Neurons
  \node[sumnode] (sum1) at (7.0, 2) {$+$};
  \node[sumnode] (sumd) at (7.0, 0.2) {$+$};
  \node[font=\tiny] at (7.0, 1.2) {$\vdots$};
  
  \node[outnrn] (h1) at (7.8, 2) {$h_1$};
  \node[outnrn] (hd) at (7.8, 0.2) {$h_k$};
  \node[font=\tiny] at (7.8, 1.2) {$\vdots$};
  \node[lbl, above=0.2cm of h1] {Output $h$};

  % Connections B
  \foreach \s in {z1,zr} \foreach \t in {sum1,sumd}
    \draw[arr, gray!60, thin] (\s.east)--(\t.west);
  \node[text=gray, font=\footnotesize] at (5.3, -0.55) {$B$};

  % Scaling Multiplier
  \node[draw, regular polygon, regular polygon sides=3, rotate=-90, fill=gray!10, minimum size=0.9cm, draw=gray!60] (scale) at (5.3, 0.8) {$\frac{\alpha}{r}$};

  % Forward path connections
  \draw[arr, gray!60] (x1.east) -- ($(W0.west)+(0,0.5)$);
  \draw[arr, gray!60] (xk.east) -- ($(W0.west)+(0,-0.5)$);
  \draw[arr, gray!60] ($(W0.east)+(0,0.5)$) -- (sum1.west);
  \draw[arr, gray!60] ($(W0.east)+(0,-0.5)$) -- (sumd.west);
  \draw[arr, gray!60] (sum1) -- (h1);
  \draw[arr, gray!60] (sumd) -- (hd);

  % Loss and Target
  \node[datblk] (y) at (11.2, 0.2) {Target $y$};
  \node[lossblk] (loss) at (11.2, 2.0) {Loss $\mathcal{L}$};
  \draw[arr, gray!60] (h1.east) -- (loss.west);
  \draw[arr, gray!60] (y.north) -- (loss.south);
  
  % Dashed base outlines fit
  \node[fit=(x1) (x2) (xk)] (fit_x) {};
  \node[fit=(W0)] (fit_w0) {};
  \node[fit=(z1) (zr)] (fit_z) {};
  \node[fit=(sum1) (sumd) (h1) (hd)] (fit_h) {};

  \node[draw=yellow!80!black, dashed, rounded corners=4pt, fit=(sum1) (sumd) (h1) (hd), inner sep=5pt, line width=1.5pt] (circle_h) {};
  \node[text=yellow!80!black, font=\bfseries\small] (lbl_h) at (7.5, -2.2) {HIGHLIGHT: Combined Output $h$};
  \draw[dashed, yellow!80!black, line width=1pt, ->] (lbl_h.north) -- (circle_h.south);

\end{tikzpicture}
\end{adjustbox}
\vspace{1.5em}
\end{center}

      \item \textbf{Mental Model}: This is the final row of output dials. It is created by taking the general dial settings computed by the permanent circuit board ($W_0$) and nudging them slightly using the task-specific dial settings computed by our adapters.
      \item \textbf{Neuron-Level Physical Explanation}: A single output neuron $j$ sums the signals from the frozen base connection path and the scaled trainable adapter path: $h_j = \sum_i W_{0, ji} x_i + \frac{\alpha}{r} \sum_m B_{jm} z_m$. The neuron acts as a physical summing junction, integrating the general-purpose signal and the specialized task signal.
      \item \textbf{Mathematical Representation}: $h \in \mathbb{R}^{k \times 1}$
      \item \textbf{Code Representation}: \texttt{h = base\_layer(x) + (alpha / r) * B(z)}
      \item \textbf{Relationships}: Fed as input to subsequent layers of the network.
    \end{itemize}

  \item \textbf{Scalar Loss Value ($\mathcal{L}$)}:
    \begin{itemize}
      \item \textbf{Formula}: $\mathcal{L}$ (e.g. Cross-Entropy loss)
      \item \textbf{Neuron Highlight Diagram}:

\begin{center}
\begin{adjustbox}{max width=0.85\linewidth}
\begin{tikzpicture}[node distance=1.0cm and 1.5cm, scale=0.85, every node/.style={transform shape}]
  \useasboundingbox (-1.8, -3.5) rectangle (12.8, 3.8);
  % Input Neurons
  \node[neuron] (x1) at (0, 2) {$x_1$};
  \node[neuron] (x2) at (0, 0.8) {$x_2$};
  \node[neuron] (xk) at (0, -0.8) {$x_d$};
  \node[font=\tiny] at (0, 0.1) {$\vdots$};
  \node[lbl, above=0.2cm of x1] {Input $x$};

  % Base Weight Matrix W_0
  \node[frzblk, minimum width=3.2cm, minimum height=1.5cm] (W0) at (3.5, 2.5) {\textbf{Frozen Base} \\ $W_0 \in \mathbb{R}^{k \times d}$};

  % Bottleneck Neurons
  \node[bneck] (z1) at (3.5, -0.3) {$z_1$};
  \node[bneck] (zr) at (3.5, -1.5) {$z_r$};
  \node[font=\tiny] at (3.5, -0.9) {$\vdots$};
  \node[lbl, above=0.2cm of z1] {Bottleneck $z$};

  % Connections A
  \foreach \s in {x1,x2,xk} \foreach \t in {z1,zr}
    \draw[arr, gray!60, thin] (\s.east)--(\t.west);
  \node[text=gray, font=\footnotesize] at (1.5, -0.25) {$A$};

  % Summing Junctions and Output Neurons
  \node[sumnode] (sum1) at (7.0, 2) {$+$};
  \node[sumnode] (sumd) at (7.0, 0.2) {$+$};
  \node[font=\tiny] at (7.0, 1.2) {$\vdots$};
  
  \node[outnrn] (h1) at (7.8, 2) {$h_1$};
  \node[outnrn] (hd) at (7.8, 0.2) {$h_k$};
  \node[font=\tiny] at (7.8, 1.2) {$\vdots$};
  \node[lbl, above=0.2cm of h1] {Output $h$};

  % Connections B
  \foreach \s in {z1,zr} \foreach \t in {sum1,sumd}
    \draw[arr, gray!60, thin] (\s.east)--(\t.west);
  \node[text=gray, font=\footnotesize] at (5.3, -0.55) {$B$};

  % Scaling Multiplier
  \node[draw, regular polygon, regular polygon sides=3, rotate=-90, fill=gray!10, minimum size=0.9cm, draw=gray!60] (scale) at (5.3, 0.8) {$\frac{\alpha}{r}$};

  % Forward path connections
  \draw[arr, gray!60] (x1.east) -- ($(W0.west)+(0,0.5)$);
  \draw[arr, gray!60] (xk.east) -- ($(W0.west)+(0,-0.5)$);
  \draw[arr, gray!60] ($(W0.east)+(0,0.5)$) -- (sum1.west);
  \draw[arr, gray!60] ($(W0.east)+(0,-0.5)$) -- (sumd.west);
  \draw[arr, gray!60] (sum1) -- (h1);
  \draw[arr, gray!60] (sumd) -- (hd);

  % Loss and Target
  \node[datblk] (y) at (11.2, 0.2) {Target $y$};
  \node[lossblk] (loss) at (11.2, 2.0) {Loss $\mathcal{L}$};
  \draw[arr, gray!60] (h1.east) -- (loss.west);
  \draw[arr, gray!60] (y.north) -- (loss.south);
  
  % Dashed base outlines fit
  \node[fit=(x1) (x2) (xk)] (fit_x) {};
  \node[fit=(W0)] (fit_w0) {};
  \node[fit=(z1) (zr)] (fit_z) {};
  \node[fit=(sum1) (sumd) (h1) (hd)] (fit_h) {};

  \node[draw=red, dashed, rounded corners=4pt, fit=(loss), inner sep=4pt, line width=1.5pt] (circle_loss) {};
  \node[text=red, font=\bfseries\small, right=0.4cm of circle_loss, align=left] (lbl_loss) {HIGHLIGHT:\\Loss $\mathcal{L}$};
  \draw[dashed, red, line width=1pt, ->] (lbl_loss.west) -- (circle_loss.east);

\end{tikzpicture}
\end{adjustbox}
\vspace{1.5em}
\end{center}

      \item \textbf{Mental Model}: This is a single scorecard number. It tells us how far off the final output dials are from what they should be. If the score is high, it means the model is making bad guesses. The goal of training is to adjust the new dials until the scorecard reads zero.
      \item \textbf{Neuron-Level Physical Explanation}: The loss $\mathcal{L}$ measures how far the firing rates of the output neurons in the final layer differ from the desired target firing rates. The target represents the correct next word, and the error signal is generated by comparing the target with the actual outputs.
      \item \textbf{Mathematical Representation}: $\mathcal{L} \in \mathbb{R}$
      \item \textbf{Code Representation}: \texttt{loss = F.cross\_entropy(logits, targets)}
      \item \textbf{Relationships}: Traversed in reverse during backpropagation; all gradients are derivatives of $\mathcal{L}$.
    \end{itemize}

  \item \textbf{Output Layer Gradient ($\nabla_h \mathcal{L}$)}:
    \begin{itemize}
      \item \textbf{Formula}: $\nabla_h \mathcal{L} = \frac{\partial \mathcal{L}}{\partial h}$
      \item \textbf{Neuron Highlight Diagram}:

\begin{center}
\begin{adjustbox}{max width=0.85\linewidth}
\begin{tikzpicture}[node distance=1.0cm and 1.5cm, scale=0.85, every node/.style={transform shape}]
  \useasboundingbox (-1.8, -3.5) rectangle (12.8, 3.8);
  % Input Neurons
  \node[neuron] (x1) at (0, 2) {$x_1$};
  \node[neuron] (x2) at (0, 0.8) {$x_2$};
  \node[neuron] (xk) at (0, -0.8) {$x_d$};
  \node[font=\tiny] at (0, 0.1) {$\vdots$};
  \node[lbl, above=0.2cm of x1] {Input $x$};

  % Base Weight Matrix W_0
  \node[frzblk, minimum width=3.2cm, minimum height=1.5cm] (W0) at (3.5, 2.5) {\textbf{Frozen Base} \\ $W_0 \in \mathbb{R}^{k \times d}$};

  % Bottleneck Neurons
  \node[bneck] (z1) at (3.5, -0.3) {$z_1$};
  \node[bneck] (zr) at (3.5, -1.5) {$z_r$};
  \node[font=\tiny] at (3.5, -0.9) {$\vdots$};
  \node[lbl, above=0.2cm of z1] {Bottleneck $z$};

  % Connections A
  \foreach \s in {x1,x2,xk} \foreach \t in {z1,zr}
    \draw[arr, gray!60, thin] (\s.east)--(\t.west);
  \node[text=gray, font=\footnotesize] at (1.5, -0.25) {$A$};

  % Summing Junctions and Output Neurons
  \node[sumnode] (sum1) at (7.0, 2) {$+$};
  \node[sumnode] (sumd) at (7.0, 0.2) {$+$};
  \node[font=\tiny] at (7.0, 1.2) {$\vdots$};
  
  \node[outnrn] (h1) at (7.8, 2) {$h_1$};
  \node[outnrn] (hd) at (7.8, 0.2) {$h_k$};
  \node[font=\tiny] at (7.8, 1.2) {$\vdots$};
  \node[lbl, above=0.2cm of h1] {Output $h$};

  % Connections B
  \foreach \s in {z1,zr} \foreach \t in {sum1,sumd}
    \draw[arr, gray!60, thin] (\s.east)--(\t.west);
  \node[text=gray, font=\footnotesize] at (5.3, -0.55) {$B$};

  % Scaling Multiplier
  \node[draw, regular polygon, regular polygon sides=3, rotate=-90, fill=gray!10, minimum size=0.9cm, draw=gray!60] (scale) at (5.3, 0.8) {$\frac{\alpha}{r}$};

  % Forward path connections
  \draw[arr, gray!60] (x1.east) -- ($(W0.west)+(0,0.5)$);
  \draw[arr, gray!60] (xk.east) -- ($(W0.west)+(0,-0.5)$);
  \draw[arr, gray!60] ($(W0.east)+(0,0.5)$) -- (sum1.west);
  \draw[arr, gray!60] ($(W0.east)+(0,-0.5)$) -- (sumd.west);
  \draw[arr, gray!60] (sum1) -- (h1);
  \draw[arr, gray!60] (sumd) -- (hd);

  % Loss and Target
  \node[datblk] (y) at (11.2, 0.2) {Target $y$};
  \node[lossblk] (loss) at (11.2, 2.0) {Loss $\mathcal{L}$};
  \draw[arr, gray!60] (h1.east) -- (loss.west);
  \draw[arr, gray!60] (y.north) -- (loss.south);
  
  % Dashed base outlines fit
  \node[fit=(x1) (x2) (xk)] (fit_x) {};
  \node[fit=(W0)] (fit_w0) {};
  \node[fit=(z1) (zr)] (fit_z) {};
  \node[fit=(sum1) (sumd) (h1) (hd)] (fit_h) {};

  \draw[arr, red, dashed, line width=1.5pt, ->] (loss.west) to[bend right=15] node[midway, above, red, font=\bfseries\small] {$\nabla_h \mathcal{L}$} (h1.east);
  \draw[arr, red, dashed, line width=1.5pt, ->] (loss.west) to[bend left=15] node[midway, below, red, font=\bfseries\small] {$\nabla_h \mathcal{L}$} (hd.east);
  \node[text=red, font=\bfseries\small] (lbl_gradh) at (7.8, -2.5) {HIGHLIGHT: Output Gradient $\nabla_h \mathcal{L}$};
  \draw[dashed, red, line width=1pt, ->] (lbl_gradh.north) -- (7.8, -0.2);

\end{tikzpicture}
\end{adjustbox}
\vspace{1.5em}
\end{center}

      \item \textbf{Mental Model}: This is a set of correction instructions for the final output dials. It says: "dial 1 needs to turn clockwise by $0.2$, and dial 2 needs to turn counter-clockwise by $1.5$" to lower the error score.
      \item \textbf{Neuron-Level Physical Explanation}: For a single output neuron $j$, the element $(\nabla_h \mathcal{L})_j = \frac{\partial \mathcal{L}}{\partial h_j}$ represents the error feedback signal. It tells the neuron whether it fired too strongly (positive gradient) or too weakly (negative gradient) and by how much it needs to adjust its total output to minimize the error.
      \item \textbf{Mathematical Representation}: $\nabla_h \mathcal{L} \in \mathbb{R}^{k \times 1}$
      \item \textbf{Code Representation}: automatically tracked by PyTorch's \texttt{loss.backward()}, represented as the grad tensor associated with output \texttt{h.grad}
      \item \textbf{Relationships}: Directly drives the gradient computation for $B$ and $z$.
    \end{itemize}

  \item \textbf{Bottleneck Layer Gradient ($\nabla_z \mathcal{L}$)}:
    \begin{itemize}
      \item \textbf{Formula}: $\nabla_z \mathcal{L} = \frac{\alpha}{r} B^\top \nabla_h \mathcal{L}$
      \item \textbf{Neuron Highlight Diagram}:

\begin{center}
\begin{adjustbox}{max width=0.85\linewidth}
\begin{tikzpicture}[node distance=1.0cm and 1.5cm, scale=0.85, every node/.style={transform shape}]
  \useasboundingbox (-1.8, -3.5) rectangle (12.8, 3.8);
  % Input Neurons
  \node[neuron] (x1) at (0, 2) {$x_1$};
  \node[neuron] (x2) at (0, 0.8) {$x_2$};
  \node[neuron] (xk) at (0, -0.8) {$x_d$};
  \node[font=\tiny] at (0, 0.1) {$\vdots$};
  \node[lbl, above=0.2cm of x1] {Input $x$};

  % Base Weight Matrix W_0
  \node[frzblk, minimum width=3.2cm, minimum height=1.5cm] (W0) at (3.5, 2.5) {\textbf{Frozen Base} \\ $W_0 \in \mathbb{R}^{k \times d}$};

  % Bottleneck Neurons
  \node[bneck] (z1) at (3.5, -0.3) {$z_1$};
  \node[bneck] (zr) at (3.5, -1.5) {$z_r$};
  \node[font=\tiny] at (3.5, -0.9) {$\vdots$};
  \node[lbl, above=0.2cm of z1] {Bottleneck $z$};

  % Connections A
  \foreach \s in {x1,x2,xk} \foreach \t in {z1,zr}
    \draw[arr, gray!60, thin] (\s.east)--(\t.west);
  \node[text=gray, font=\footnotesize] at (1.5, -0.25) {$A$};

  % Summing Junctions and Output Neurons
  \node[sumnode] (sum1) at (7.0, 2) {$+$};
  \node[sumnode] (sumd) at (7.0, 0.2) {$+$};
  \node[font=\tiny] at (7.0, 1.2) {$\vdots$};
  
  \node[outnrn] (h1) at (7.8, 2) {$h_1$};
  \node[outnrn] (hd) at (7.8, 0.2) {$h_k$};
  \node[font=\tiny] at (7.8, 1.2) {$\vdots$};
  \node[lbl, above=0.2cm of h1] {Output $h$};

  % Connections B
  \foreach \s in {z1,zr} \foreach \t in {sum1,sumd}
    \draw[arr, gray!60, thin] (\s.east)--(\t.west);
  \node[text=gray, font=\footnotesize] at (5.3, -0.55) {$B$};

  % Scaling Multiplier
  \node[draw, regular polygon, regular polygon sides=3, rotate=-90, fill=gray!10, minimum size=0.9cm, draw=gray!60] (scale) at (5.3, 0.8) {$\frac{\alpha}{r}$};

  % Forward path connections
  \draw[arr, gray!60] (x1.east) -- ($(W0.west)+(0,0.5)$);
  \draw[arr, gray!60] (xk.east) -- ($(W0.west)+(0,-0.5)$);
  \draw[arr, gray!60] ($(W0.east)+(0,0.5)$) -- (sum1.west);
  \draw[arr, gray!60] ($(W0.east)+(0,-0.5)$) -- (sumd.west);
  \draw[arr, gray!60] (sum1) -- (h1);
  \draw[arr, gray!60] (sumd) -- (hd);

  % Loss and Target
  \node[datblk] (y) at (11.2, 0.2) {Target $y$};
  \node[lossblk] (loss) at (11.2, 2.0) {Loss $\mathcal{L}$};
  \draw[arr, gray!60] (h1.east) -- (loss.west);
  \draw[arr, gray!60] (y.north) -- (loss.south);
  
  % Dashed base outlines fit
  \node[fit=(x1) (x2) (xk)] (fit_x) {};
  \node[fit=(W0)] (fit_w0) {};
  \node[fit=(z1) (zr)] (fit_z) {};
  \node[fit=(sum1) (sumd) (h1) (hd)] (fit_h) {};

  \draw[arr, red, dashed, line width=1.5pt, ->] (h1.west) to[bend right=15] node[midway, above, red, font=\bfseries\small] {$\nabla_z \mathcal{L}$} (z1.east);
  \draw[arr, red, dashed, line width=1.5pt, ->] (hd.west) to[bend left=15] node[midway, below, red, font=\bfseries\small] {$\nabla_z \mathcal{L}$} (zr.east);
  \node[text=red, font=\bfseries\small] (lbl_gradz) at (3.5, -2.5) {HIGHLIGHT: Bottleneck Gradient $\nabla_z \mathcal{L}$};
  \draw[dashed, red, line width=1pt, ->] (lbl_gradz.north) -- (3.5, -1.8);

\end{tikzpicture}
\end{adjustbox}
\vspace{1.5em}
\end{center}

      \item \textbf{Mental Model}: This translates the correction instructions from the output room back down into the bottleneck. It tells the summary dials inside the funnel how they need to adjust to fix the output dials.
      \item \textbf{Neuron-Level Physical Explanation}: For a bottleneck neuron $m$, this is the backpropagated error signal. It is computed by summing the error feedback signals from all output neurons, weighted by their connections $B_{jm}$ and scaled: $(\nabla_z \mathcal{L})_m = \frac{\alpha}{r} \sum_j B_{jm} (\nabla_h \mathcal{L})_j$. It tells the bottleneck neuron how its activation should have changed.
      \item \textbf{Mathematical Representation}: $\nabla_z \mathcal{L} \in \mathbb{R}^{r \times 1}$
      \item \textbf{Code Representation}: \texttt{grad\_z = (alpha / r) * torch.matmul(grad\_output, B.weight)}
      \item \textbf{Relationships}: Drives the gradient computation for $A$.
    \end{itemize}

  \item \textbf{Update Gradient for Matrix $B$ ($\nabla_B \mathcal{L}$)}:
    \begin{itemize}
      \item \textbf{Formula}: $\nabla_B \mathcal{L} = \frac{\alpha}{r} (\nabla_h \mathcal{L}) z^\top$
      \item \textbf{Neuron Highlight Diagram}:

\begin{center}
\begin{adjustbox}{max width=0.85\linewidth}
\begin{tikzpicture}[node distance=1.0cm and 1.5cm, scale=0.85, every node/.style={transform shape}]
  \useasboundingbox (-1.8, -3.5) rectangle (12.8, 3.8);
  % Input Neurons
  \node[neuron] (x1) at (0, 2) {$x_1$};
  \node[neuron] (x2) at (0, 0.8) {$x_2$};
  \node[neuron] (xk) at (0, -0.8) {$x_d$};
  \node[font=\tiny] at (0, 0.1) {$\vdots$};
  \node[lbl, above=0.2cm of x1] {Input $x$};

  % Base Weight Matrix W_0
  \node[frzblk, minimum width=3.2cm, minimum height=1.5cm] (W0) at (3.5, 2.5) {\textbf{Frozen Base} \\ $W_0 \in \mathbb{R}^{k \times d}$};

  % Bottleneck Neurons
  \node[bneck] (z1) at (3.5, -0.3) {$z_1$};
  \node[bneck] (zr) at (3.5, -1.5) {$z_r$};
  \node[font=\tiny] at (3.5, -0.9) {$\vdots$};
  \node[lbl, above=0.2cm of z1] {Bottleneck $z$};

  % Connections A
  \foreach \s in {x1,x2,xk} \foreach \t in {z1,zr}
    \draw[arr, gray!60, thin] (\s.east)--(\t.west);
  \node[text=gray, font=\footnotesize] at (1.5, -0.25) {$A$};

  % Summing Junctions and Output Neurons
  \node[sumnode] (sum1) at (7.0, 2) {$+$};
  \node[sumnode] (sumd) at (7.0, 0.2) {$+$};
  \node[font=\tiny] at (7.0, 1.2) {$\vdots$};
  
  \node[outnrn] (h1) at (7.8, 2) {$h_1$};
  \node[outnrn] (hd) at (7.8, 0.2) {$h_k$};
  \node[font=\tiny] at (7.8, 1.2) {$\vdots$};
  \node[lbl, above=0.2cm of h1] {Output $h$};

  % Connections B
  \foreach \s in {z1,zr} \foreach \t in {sum1,sumd}
    \draw[arr, gray!60, thin] (\s.east)--(\t.west);
  \node[text=gray, font=\footnotesize] at (5.3, -0.55) {$B$};

  % Scaling Multiplier
  \node[draw, regular polygon, regular polygon sides=3, rotate=-90, fill=gray!10, minimum size=0.9cm, draw=gray!60] (scale) at (5.3, 0.8) {$\frac{\alpha}{r}$};

  % Forward path connections
  \draw[arr, gray!60] (x1.east) -- ($(W0.west)+(0,0.5)$);
  \draw[arr, gray!60] (xk.east) -- ($(W0.west)+(0,-0.5)$);
  \draw[arr, gray!60] ($(W0.east)+(0,0.5)$) -- (sum1.west);
  \draw[arr, gray!60] ($(W0.east)+(0,-0.5)$) -- (sumd.west);
  \draw[arr, gray!60] (sum1) -- (h1);
  \draw[arr, gray!60] (sumd) -- (hd);

  % Loss and Target
  \node[datblk] (y) at (11.2, 0.2) {Target $y$};
  \node[lossblk] (loss) at (11.2, 2.0) {Loss $\mathcal{L}$};
  \draw[arr, gray!60] (h1.east) -- (loss.west);
  \draw[arr, gray!60] (y.north) -- (loss.south);
  
  % Dashed base outlines fit
  \node[fit=(x1) (x2) (xk)] (fit_x) {};
  \node[fit=(W0)] (fit_w0) {};
  \node[fit=(z1) (zr)] (fit_z) {};
  \node[fit=(sum1) (sumd) (h1) (hd)] (fit_h) {};

  \foreach \s in {z1,zr} \foreach \t in {sum1,sumd}
    \draw[arr, red, line width=1.5pt] (\s.east)--(\t.west);
  \node[draw=red, rectangle, fill=red!10, line width=1.5pt, font=\small, align=center] (gradB) at (5.3, -2.5) {HIGHLIGHT: Update $B$ \\ $\nabla_B \mathcal{L} = \frac{\alpha}{r} (\nabla_h \mathcal{L}) z^\top$};
  \draw[dashed, red, line width=1.2pt, ->] (zr.south) to[bend right=15] (gradB.west);
  \draw[dashed, red, line width=1.2pt, ->] (hd.south) to[bend left=15] (gradB.east);
  \draw[dashed, red, line width=1.5pt, ->] (gradB.north) -- (5.3, -0.8) node[midway, right, red, font=\bfseries\tiny] {Updates};

\end{tikzpicture}
\end{adjustbox}
\vspace{1.5em}
\end{center}

      \item \textbf{Mental Model}: This is the instruction sheet showing how to rewire the dials of the expansion nozzle ($B$) so that they do a better job of translating the summary back to the output dials.
      \item \textbf{Neuron-Level Physical Explanation}: For the synapse connecting bottleneck neuron $m$ to output neuron $j$, the gradient is calculated as $(\nabla_B \mathcal{L})_{jm} = \frac{\alpha}{r} (\nabla_h \mathcal{L})_j z_m$. This follows Hebbian learning: the strength of the update is proportional to the product of the pre-synaptic activation ($z_m$) and the post-synaptic error signal ($(\nabla_h \mathcal{L})_j$).
      \item \textbf{Mathematical Representation}: $\nabla_B \mathcal{L} \in \mathbb{R}^{k \times r}$
      \item \textbf{Code Representation}: \texttt{B.weight.grad = (alpha / r) * torch.matmul(grad\_output.t(), z)}
      \item \textbf{Relationships}: Used by the optimizer to update the values of $B$.
    \end{itemize}

  \item \textbf{Update Gradient for Matrix $A$ ($\nabla_A \mathcal{L}$)}:
    \begin{itemize}
      \item \textbf{Formula}: $\nabla_A \mathcal{L} = (\nabla_z \mathcal{L}) x^\top$
      \item \textbf{Neuron Highlight Diagram}:

\begin{center}
\begin{adjustbox}{max width=0.85\linewidth}
\begin{tikzpicture}[node distance=1.0cm and 1.5cm, scale=0.85, every node/.style={transform shape}]
  \useasboundingbox (-1.8, -3.5) rectangle (12.8, 3.8);
  % Input Neurons
  \node[neuron] (x1) at (0, 2) {$x_1$};
  \node[neuron] (x2) at (0, 0.8) {$x_2$};
  \node[neuron] (xk) at (0, -0.8) {$x_d$};
  \node[font=\tiny] at (0, 0.1) {$\vdots$};
  \node[lbl, above=0.2cm of x1] {Input $x$};

  % Base Weight Matrix W_0
  \node[frzblk, minimum width=3.2cm, minimum height=1.5cm] (W0) at (3.5, 2.5) {\textbf{Frozen Base} \\ $W_0 \in \mathbb{R}^{k \times d}$};

  % Bottleneck Neurons
  \node[bneck] (z1) at (3.5, -0.3) {$z_1$};
  \node[bneck] (zr) at (3.5, -1.5) {$z_r$};
  \node[font=\tiny] at (3.5, -0.9) {$\vdots$};
  \node[lbl, above=0.2cm of z1] {Bottleneck $z$};

  % Connections A
  \foreach \s in {x1,x2,xk} \foreach \t in {z1,zr}
    \draw[arr, gray!60, thin] (\s.east)--(\t.west);
  \node[text=gray, font=\footnotesize] at (1.5, -0.25) {$A$};

  % Summing Junctions and Output Neurons
  \node[sumnode] (sum1) at (7.0, 2) {$+$};
  \node[sumnode] (sumd) at (7.0, 0.2) {$+$};
  \node[font=\tiny] at (7.0, 1.2) {$\vdots$};
  
  \node[outnrn] (h1) at (7.8, 2) {$h_1$};
  \node[outnrn] (hd) at (7.8, 0.2) {$h_k$};
  \node[font=\tiny] at (7.8, 1.2) {$\vdots$};
  \node[lbl, above=0.2cm of h1] {Output $h$};

  % Connections B
  \foreach \s in {z1,zr} \foreach \t in {sum1,sumd}
    \draw[arr, gray!60, thin] (\s.east)--(\t.west);
  \node[text=gray, font=\footnotesize] at (5.3, -0.55) {$B$};

  % Scaling Multiplier
  \node[draw, regular polygon, regular polygon sides=3, rotate=-90, fill=gray!10, minimum size=0.9cm, draw=gray!60] (scale) at (5.3, 0.8) {$\frac{\alpha}{r}$};

  % Forward path connections
  \draw[arr, gray!60] (x1.east) -- ($(W0.west)+(0,0.5)$);
  \draw[arr, gray!60] (xk.east) -- ($(W0.west)+(0,-0.5)$);
  \draw[arr, gray!60] ($(W0.east)+(0,0.5)$) -- (sum1.west);
  \draw[arr, gray!60] ($(W0.east)+(0,-0.5)$) -- (sumd.west);
  \draw[arr, gray!60] (sum1) -- (h1);
  \draw[arr, gray!60] (sumd) -- (hd);

  % Loss and Target
  \node[datblk] (y) at (11.2, 0.2) {Target $y$};
  \node[lossblk] (loss) at (11.2, 2.0) {Loss $\mathcal{L}$};
  \draw[arr, gray!60] (h1.east) -- (loss.west);
  \draw[arr, gray!60] (y.north) -- (loss.south);
  
  % Dashed base outlines fit
  \node[fit=(x1) (x2) (xk)] (fit_x) {};
  \node[fit=(W0)] (fit_w0) {};
  \node[fit=(z1) (zr)] (fit_z) {};
  \node[fit=(sum1) (sumd) (h1) (hd)] (fit_h) {};

  \foreach \s in {x1,x2,xk} \foreach \t in {z1,zr}
    \draw[arr, red, line width=1.5pt] (\s.east)--(\t.west);
  \node[draw=red, rectangle, fill=red!10, line width=1.5pt, font=\small, align=center] (gradA) at (1.5, -2.5) {HIGHLIGHT: Update $A$ \\ $\nabla_A \mathcal{L} = (\nabla_z \mathcal{L}) x^\top$};
  \draw[dashed, red, line width=1.2pt, ->] (xk.south) to[bend right=15] (gradA.west);
  \draw[dashed, red, line width=1.2pt, ->] (zr.south) to[bend left=15] (gradA.east);
  \draw[dashed, red, line width=1.5pt, ->] (gradA.north) -- (1.5, -0.5) node[midway, left, red, font=\bfseries\tiny] {Updates};

\end{tikzpicture}
\end{adjustbox}
\vspace{1.5em}
\end{center}

      \item \textbf{Mental Model}: This is the instruction sheet showing how to rewire the dials of the compression funnel ($A$) so that they do a better job of compressing the input dials into the summary.
      \item \textbf{Neuron-Level Physical Explanation}: For the synapse connecting input neuron $i$ to bottleneck neuron $m$, the gradient is $(\nabla_A \mathcal{L})_{mi} = (\nabla_z \mathcal{L})_m x_i$. This also follows Hebbian learning: it multiplies the pre-synaptic input signal ($x_i$) by the post-synaptic bottleneck error signal ($(\nabla_z \mathcal{L})_m$) to determine how much the synapse should strengthen or weaken.
      \item \textbf{Mathematical Representation}: $\nabla_A \mathcal{L} \in \mathbb{R}^{r \times d}$
      \item \textbf{Code Representation}: \texttt{A.weight.grad = torch.matmul(grad\_z.t(), x)}
      \item \textbf{Relationships}: Used by the optimizer to update the values of $A$.
    \end{itemize}
\end{enumerate}

\subsection{The Physical and Philosophical Foundations of LoRA}

\subsubsection{Philosophical Foundations: Occam's Razor and the Manifold Hypothesis}
The success of Low-Rank Adaptation (LoRA) is not merely an empirical coincidence, but is deeply rooted in two foundational philosophical principles of complexity. We analyze each principle through a dialectic framework of Arguments, Counter-Arguments, and Rebuttals:

\paragraph{1. The Manifold Hypothesis}
\begin{itemize}
  \item \textbf{Argument}: This hypothesis posits that high-dimensional data (such as the activations in a neural network) actually lies on a much lower-dimensional, highly structured manifold. Pretraining constructs this manifold by mapping the topological structure of natural language.
  
  In the context of Large Language Models, the embedding space has an extrinsic dimension of d = 4096 (e.g., $d_{\text{model}} = 4096$). However, the set of all grammatically correct and semantically coherent sentences does not fill this space uniformly; instead, it is confined to a lower-dimensional manifold. When we adapt the model to a downstream task (e.g., teaching it Python syntax or chat style), we do not need to construct a new manifold. We only need to compute a small trajectory or translation vector along the existing surface. This translation vector corresponds to a tangent vector in the manifold's tangent bundle. By confining our search to this tangent space, we avoid distorting the global topology of the language manifold, preventing catastrophic forgetting.
  
  \item \textbf{Counter-Argument}: If the downstream task represents an extreme domain shift (e.g., shifting from English literature to highly specialized biochemistry nomenclature or low-level assembly code), the target representations may lie completely outside the manifold mapped during pretraining. In such cases, the translation vector cannot lie along the existing manifold surface, and restricting the updates to a low-rank tangent space will lead to underfitting. The model would fail to acquire the new features because the required coordinates do not exist in the pretrained column space.
  
  \item \textbf{Rebuttal / Synthesis}: While severe domain shifts can strain low-rank constraints, empirical evidence shows that even highly specialized knowledge maps back to general-purpose linguistic structures (syntax, logic, context). Furthermore, if the target lies slightly off-manifold, a small rank $r$ (e.g., 8 to 16) is still sufficient to project the input into the nearest point on the task manifold. If needed, the user can scale the rank $r$ to increase the degrees of freedom, which mathematically expands the coordinate projection space, allowing the model to capture features that lie slightly outside the pretraining manifold while still regularizing the update to prevent arbitrary drift.
\end{itemize}

\paragraph{2. Occam's Razor (Lex Parsimoniae)}
\begin{itemize}
  \item \textbf{Argument}: Plurality should not be posited without necessity. Full fine-tuning searches an extremely large, redundant parameter space ($\mathbb{R}^{k \times d}$), which contains millions of redundant degrees of freedom. This massive parameter space allows the model to find arbitrary, highly complex fitting functions that overfit the downstream training data and overwrite pretrained general knowledge.
  
  By restricting the learning search to a compressed subspace ($\mathbb{R}^{r \times (d+k)}$), LoRA acts as a mathematical regularization agent. It forces the model to find the simplest, most parsimonious path of adaptation. This acts as a geometric constraint (projection regularization), restricting updates to only the principal directions of adaptation, thereby preserving the model's core foundational capabilities.
  
  \item \textbf{Counter-Argument}: By restricting the parameters to a low-rank product $BA$, we might force the model to find an oversimplified adaptation path that fails to capture the subtle, non-linear correlations required for complex reasoning tasks. Furthermore, Occam's Razor assumes that simpler models generalize better, but in deep learning, the phenomenon of "double descent" shows that highly overparameterized models can generalize better than simpler ones by finding smooth interpolating functions in high-dimensional spaces.
  
  \item \textbf{Rebuttal / Synthesis}: Double descent occurs when the entire network is overparameterized and trained on vast, diverse datasets where it can interpolate smoothly. However, in downstream fine-tuning, the dataset is small and highly focused. In this regime, overparameterization leads to catastrophic forgetting and representation collapse. Constraining the updates to a low-rank subspace restricts the model from modifying the general-purpose features, while allowing sufficient capacity to learn the task. Thus, LoRA combines the benefits of high-dimensional pretraining (smooth interpolation on general data) with low-dimensional adaptation (strict regularization on task-specific data).
\end{itemize}

\subsubsection{Internal Mechanics and Representation Space Geometry}
To understand why LoRA works under the hood, we must analyze the geometry of representation space and the intrinsic dimension of neural networks:
\begin{itemize}
  \item \textbf{Intrinsic vs. Extrinsic Dimension}: A pretrained model has a very high extrinsic dimension (the number of weights $W_0 \in \mathbb{R}^{k \times d}$). However, empirical work by Aghajanyan et al. (2020) demonstrated that the parameter space of pretrained language models has a much lower \textit{intrinsic dimension}. That is, the directions in weight space that significantly affect the loss are highly concentrated on a low-dimensional subspace. This is why we can optimize a tiny fraction of parameters and achieve the same or better task performance.
  \item \textbf{Subspace Alignment and Trajectory Geometry}: The activation vectors $x$ represent coordinates on the manifold. The base weight matrix $W_0$ maps these coordinates to the next representation space. The update matrix $\Delta W = \frac{\alpha}{r} B A$ acts as a linear projection operator that projects the activations into a tiny subspace of dimension $r$ via $A$, rotates/scales them in this bottleneck space, and then projects them back to the output space via $B$. This operation restricts the adaptation to a set of $r$ directional vectors. Instead of shifting the entire coordinate frame, the adapter only redirects the activation trajectory along the principal directions relevant to the downstream task.
  
  Geometrically, the column space of $\Delta W$ is spanned by the columns of $B$, which is an $r$-dimensional subspace of the output space $\mathbb{R}^{k}$. The row space of $\Delta W$ is spanned by the rows of $A$, which is an $r$-dimensional subspace of the input space $\mathbb{R}^{d}$. Thus, the adaptation updates are strictly constrained to these coordinates, preserving all other dimensions from degradation.
\end{itemize}

\subsubsection{Visualizing the Calculations}
We can visualize how the information bottleneck constrains the updates in the representation space during the forward and backward passes:

\begin{figure}[h]
\centering
\begin{adjustbox}{max width=0.9\linewidth}
\begin{tikzpicture}[node distance=1.0cm and 1.5cm, scale=0.95, every node/.style={transform shape}]
  \node[datblk] (input) {Input Space $x$\\\\( $d$ dimensions)};
  \node[grayblk, right=of input, minimum width=2.4cm] (project) {Compression\\\\$z = Ax$};
  \node[trnblk, right=of project, minimum width=2.4cm] (bottleneck) {Bottleneck Space $z$\\\\( $r$ dimensions)};
  \node[grayblk, right=of bottleneck, minimum width=2.4cm] (expand) {Expansion\\\\$B z$};
  \node[datblk, right=of expand, minimum width=2.4cm] (out_lora) {Adapter Output\\\\( $k$ dimensions)};
  
  \node[frzblk, below=1.2cm of bottleneck, minimum width=4cm] (base) {Frozen Base Path\\\\$W_0 x$};
  \node[sumnode, right=1.2cm of out_lora] (sum) {$+$};
  \node[outnrn, right=1.0cm of sum] (output) {Final output $h$};

  \draw[arr] (input) -- (project);
  \draw[arr] (project) -- (bottleneck);
  \draw[arr] (bottleneck) -- (expand);
  \draw[arr] (expand) -- (out_lora);
  \draw[arr] (out_lora) -- (sum);
  \draw[arr] (sum) -- (output);
  
  \draw[arr] (input.south) |- (base.west);
  \draw[arr] (base.east) -| (sum.south);
  
  \node[font=\footnotesize\itshape, align=center, above=0.3cm of bottleneck] {Compressed\\\\Information Canal};
\end{tikzpicture}
\end{adjustbox}
\caption{Visualization of the information bottleneck: high-dimensional input is projected into a low-dimensional canal where task-specific tuning forces are applied, then projected back.}
\end{figure}

\begin{itemize}
  \item \textbf{Forward Pass Compression}: The activation vector $x \in \mathbb{R}^{d \times 1}$ is projected down through a narrow bottleneck of rank $r$. Since $r \ll \min(d, k)$, the projection filters out dimensions irrelevant to the current task, creating a compressed task representation $z \in \mathbb{R}^{r \times 1}$.
  \item \textbf{Backward Pass Gradient Flow}: Error gradients $\nabla_h \mathcal{L}$ flow backwards. Because $W_0$ is frozen and lacks a gradient buffer, no gradient update calculations occur on the base weight matrix. Instead, the gradients flow down the low-rank adapters $B$ and $A$, computing updates specifically for the low-rank subspace and avoiding parameter updates across the whole network.
\end{itemize}

\subsection{The Physical and Philosophical Foundations of LoRA}

\subsubsection{Philosophical Foundations: Occam's Razor and the Manifold Hypothesis}
The success of Low-Rank Adaptation (LoRA) is not merely an empirical coincidence, but is deeply rooted in two foundational philosophical principles of complexity. We analyze each principle through a dialectic framework of Arguments, Counter-Arguments, and Rebuttals:

\paragraph{1. The Manifold Hypothesis}
\begin{itemize}
  \item \textbf{Argument}: This hypothesis posits that high-dimensional data (such as the activations in a neural network) actually lies on a much lower-dimensional, highly structured manifold. Pretraining constructs this manifold by mapping the topological structure of natural language.
  
  In the context of Large Language Models, the embedding space has an extrinsic dimension of d = 4096 (e.g., $d_{\text{model}} = 4096$). However, the set of all grammatically correct and semantically coherent sentences does not fill this space uniformly; instead, it is confined to a lower-dimensional manifold. When we adapt the model to a downstream task (e.g., teaching it Python syntax or chat style), we do not need to construct a new manifold. We only need to compute a small trajectory or translation vector along the existing surface. This translation vector corresponds to a tangent vector in the manifold's tangent bundle. By confining our search to this tangent space, we avoid distorting the global topology of the language manifold, preventing catastrophic forgetting.
  
  \item \textbf{Counter-Argument}: If the downstream task represents an extreme domain shift (e.g., shifting from English literature to highly specialized biochemistry nomenclature or low-level assembly code), the target representations may lie completely outside the manifold mapped during pretraining. In such cases, the translation vector cannot lie along the existing manifold surface, and restricting the updates to a low-rank tangent space will lead to underfitting. The model would fail to acquire the new features because the required coordinates do not exist in the pretrained column space.
  
  \item \textbf{Rebuttal / Synthesis}: While severe domain shifts can strain low-rank constraints, empirical evidence shows that even highly specialized knowledge maps back to general-purpose linguistic structures (syntax, logic, context). Furthermore, if the target lies slightly off-manifold, a small rank $r$ (e.g., 8 to 16) is still sufficient to project the input into the nearest point on the task manifold. If needed, the user can scale the rank $r$ to increase the degrees of freedom, which mathematically expands the coordinate projection space, allowing the model to capture features that lie slightly outside the pretraining manifold while still regularizing the update to prevent arbitrary drift.
\end{itemize}

\paragraph{2. Occam's Razor (Lex Parsimoniae)}
\begin{itemize}
  \item \textbf{Argument}: Plurality should not be posited without necessity. Full fine-tuning searches an extremely large, redundant parameter space ($\mathbb{R}^{k \times d}$), which contains millions of redundant degrees of freedom. This massive parameter space allows the model to find arbitrary, highly complex fitting functions that overfit the downstream training data and overwrite pretrained general knowledge.
  
  By restricting the learning search to a compressed subspace ($\mathbb{R}^{r \times (d+k)}$), LoRA acts as a mathematical regularization agent. It forces the model to find the simplest, most parsimonious path of adaptation. This acts as a geometric constraint (projection regularization), restricting updates to only the principal directions of adaptation, thereby preserving the model's core foundational capabilities.
  
  \item \textbf{Counter-Argument}: By restricting the parameters to a low-rank product $BA$, we might force the model to find an oversimplified adaptation path that fails to capture the subtle, non-linear correlations required for complex reasoning tasks. Furthermore, Occam's Razor assumes that simpler models generalize better, but in deep learning, the phenomenon of "double descent" shows that highly overparameterized models can generalize better than simpler ones by finding smooth interpolating functions in high-dimensional spaces.
  
  \item \textbf{Rebuttal / Synthesis}: Double descent occurs when the entire network is overparameterized and trained on vast, diverse datasets where it can interpolate smoothly. However, in downstream fine-tuning, the dataset is small and highly focused. In this regime, overparameterization leads to catastrophic forgetting and representation collapse. Constraining the updates to a low-rank subspace restricts the model from modifying the general-purpose features, while allowing sufficient capacity to learn the task. Thus, LoRA combines the benefits of high-dimensional pretraining (smooth interpolation on general data) with low-dimensional adaptation (strict regularization on task-specific data).
\end{itemize}

\subsubsection{Internal Mechanics and Representation Space Geometry}
To understand why LoRA works under the hood, we must analyze the geometry of representation space and the intrinsic dimension of neural networks:
\begin{itemize}
  \item \textbf{Intrinsic vs. Extrinsic Dimension}: A pretrained model has a very high extrinsic dimension (the number of weights $W_0 \in \mathbb{R}^{k \times d}$). However, empirical work by Aghajanyan et al. (2020) demonstrated that the parameter space of pretrained language models has a much lower \textit{intrinsic dimension}. That is, the directions in weight space that significantly affect the loss are highly concentrated on a low-dimensional subspace. This is why we can optimize a tiny fraction of parameters and achieve the same or better task performance.
  \item \textbf{Subspace Alignment and Trajectory Geometry}: The activation vectors $x$ represent coordinates on the manifold. The base weight matrix $W_0$ maps these coordinates to the next representation space. The update matrix $\Delta W = \frac{\alpha}{r} B A$ acts as a linear projection operator that projects the activations into a tiny subspace of dimension $r$ via $A$, rotates/scales them in this bottleneck space, and then projects them back to the output space via $B$. This operation restricts the adaptation to a set of $r$ directional vectors. Instead of shifting the entire coordinate frame, the adapter only redirects the activation trajectory along the principal directions relevant to the downstream task.
  
  Geometrically, the column space of $\Delta W$ is spanned by the columns of $B$, which is an $r$-dimensional subspace of the output space $\mathbb{R}^{k}$. The row space of $\Delta W$ is spanned by the rows of $A$, which is an $r$-dimensional subspace of the input space $\mathbb{R}^{d}$. Thus, the adaptation updates are strictly constrained to these coordinates, preserving all other dimensions from degradation.
\end{itemize}

\subsubsection{Visualizing the Calculations}
We can visualize how the information bottleneck constrains the updates in the representation space during the forward and backward passes:

\begin{figure}[h]
\centering
\begin{adjustbox}{max width=0.9\linewidth}
\begin{tikzpicture}[node distance=1.0cm and 1.5cm, scale=0.95, every node/.style={transform shape}]
  \node[datblk] (input) {Input Space $x$\\\\( $d$ dimensions)};
  \node[grayblk, right=of input, minimum width=2.4cm] (project) {Compression\\\\$z = Ax$};
  \node[trnblk, right=of project, minimum width=2.4cm] (bottleneck) {Bottleneck Space $z$\\\\( $r$ dimensions)};
  \node[grayblk, right=of bottleneck, minimum width=2.4cm] (expand) {Expansion\\\\$B z$};
  \node[datblk, right=of expand, minimum width=2.4cm] (out_lora) {Adapter Output\\\\( $k$ dimensions)};
  
  \node[frzblk, below=1.2cm of bottleneck, minimum width=4cm] (base) {Frozen Base Path\\\\$W_0 x$};
  \node[sumnode, right=1.2cm of out_lora] (sum) {$+$};
  \node[outnrn, right=1.0cm of sum] (output) {Final output $h$};

  \draw[arr] (input) -- (project);
  \draw[arr] (project) -- (bottleneck);
  \draw[arr] (bottleneck) -- (expand);
  \draw[arr] (expand) -- (out_lora);
  \draw[arr] (out_lora) -- (sum);
  \draw[arr] (sum) -- (output);
  
  \draw[arr] (input.south) |- (base.west);
  \draw[arr] (base.east) -| (sum.south);
  
  \node[font=\footnotesize\itshape, align=center, above=0.3cm of bottleneck] {Compressed\\\\Information Canal};
\end{tikzpicture}
\end{adjustbox}
\caption{Visualization of the information bottleneck: high-dimensional input is projected into a low-dimensional canal where task-specific tuning forces are applied, then projected back.}
\end{figure}

\begin{itemize}
  \item \textbf{Forward Pass Compression}: The activation vector $x \in \mathbb{R}^{d \times 1}$ is projected down through a narrow bottleneck of rank $r$. Since $r \ll \min(d, k)$, the projection filters out dimensions irrelevant to the current task, creating a compressed task representation $z \in \mathbb{R}^{r \times 1}$.
  \item \textbf{Backward Pass Gradient Flow}: Error gradients $\nabla_h \mathcal{L}$ flow backwards. Because $W_0$ is frozen and lacks a gradient buffer, no gradient update calculations occur on the base weight matrix. Instead, the gradients flow down the low-rank adapters $B$ and $A$, computing updates specifically for the low-rank subspace and avoiding parameter updates across the whole network.
\end{itemize}

\subsection{The Physical and Philosophical Foundations of LoRA}

\subsubsection{Philosophical Foundations: Occam's Razor and the Manifold Hypothesis}
The success of Low-Rank Adaptation (LoRA) is not merely an empirical coincidence, but is deeply rooted in two foundational philosophical principles of complexity. We analyze each principle through a dialectic framework of Arguments, Counter-Arguments, and Rebuttals:

\paragraph{1. The Manifold Hypothesis}
\begin{itemize}
  \item \textbf{Argument}: This hypothesis posits that high-dimensional data (such as the activations in a neural network) actually lies on a much lower-dimensional, highly structured manifold. Pretraining constructs this manifold by mapping the topological structure of natural language.
  
  In the context of Large Language Models, the embedding space has an extrinsic dimension of d = 4096 (e.g., $d_{\text{model}} = 4096$). However, the set of all grammatically correct and semantically coherent sentences does not fill this space uniformly; instead, it is confined to a lower-dimensional manifold. When we adapt the model to a downstream task (e.g., teaching it Python syntax or chat style), we do not need to construct a new manifold. We only need to compute a small trajectory or translation vector along the existing surface. This translation vector corresponds to a tangent vector in the manifold's tangent bundle. By confining our search to this tangent space, we avoid distorting the global topology of the language manifold, preventing catastrophic forgetting.
  
  \item \textbf{Counter-Argument}: If the downstream task represents an extreme domain shift (e.g., shifting from English literature to highly specialized biochemistry nomenclature or low-level assembly code), the target representations may lie completely outside the manifold mapped during pretraining. In such cases, the translation vector cannot lie along the existing manifold surface, and restricting the updates to a low-rank tangent space will lead to underfitting. The model would fail to acquire the new features because the required coordinates do not exist in the pretrained column space.
  
  \item \textbf{Rebuttal / Synthesis}: While severe domain shifts can strain low-rank constraints, empirical evidence shows that even highly specialized knowledge maps back to general-purpose linguistic structures (syntax, logic, context). Furthermore, if the target lies slightly off-manifold, a small rank $r$ (e.g., 8 to 16) is still sufficient to project the input into the nearest point on the task manifold. If needed, the user can scale the rank $r$ to increase the degrees of freedom, which mathematically expands the coordinate projection space, allowing the model to capture features that lie slightly outside the pretraining manifold while still regularizing the update to prevent arbitrary drift.
\end{itemize}

\paragraph{2. Occam's Razor (Lex Parsimoniae)}
\begin{itemize}
  \item \textbf{Argument}: Plurality should not be posited without necessity. Full fine-tuning searches an extremely large, redundant parameter space ($\mathbb{R}^{k \times d}$), which contains millions of redundant degrees of freedom. This massive parameter space allows the model to find arbitrary, highly complex fitting functions that overfit the downstream training data and overwrite pretrained general knowledge.
  
  By restricting the learning search to a compressed subspace ($\\mathbb{R}^{r \\times (d+k)}$), LoRA acts as a mathematical regularization agent. It forces the model to find the simplest, most parsimonious path of adaptation. This acts as a geometric constraint (projection regularization), restricting updates to only the principal directions of adaptation, thereby preserving the model's core foundational capabilities.
  
  \item \textbf{Counter-Argument}: By restricting the parameters to a low-rank product $BA$, we might force the model to find an oversimplified adaptation path that fails to capture the subtle, non-linear correlations required for complex reasoning tasks. Furthermore, Occam's Razor assumes that simpler models generalize better, but in deep learning, the phenomenon of "double descent" shows that highly overparameterized models can generalize better than simpler ones by finding smooth interpolating functions in high-dimensional spaces.
  
  \item \textbf{Rebuttal / Synthesis}: Double descent occurs when the entire network is overparameterized and trained on vast, diverse datasets where it can interpolate smoothly. However, in downstream fine-tuning, the dataset is small and highly focused. In this regime, overparameterization leads to catastrophic forgetting and representation collapse. Constraining the updates to a low-rank subspace restricts the model from modifying the general-purpose features, while allowing sufficient capacity to learn the task. Thus, LoRA combines the benefits of high-dimensional pretraining (smooth interpolation on general data) with low-dimensional adaptation (strict regularization on task-specific data).
\end{itemize}

\subsubsection{Internal Mechanics and Representation Space Geometry}
To understand why LoRA works under the hood, we must analyze the geometry of representation space and the intrinsic dimension of neural networks:
\begin{itemize}
  \item \textbf{Intrinsic vs. Extrinsic Dimension}: A pretrained model has a very high extrinsic dimension (the number of weights $W_0 \in \mathbb{R}^{k \times d}$). However, empirical work by Aghajanyan et al. (2020) demonstrated that the parameter space of pretrained language models has a much lower \textit{intrinsic dimension}. That is, the directions in weight space that significantly affect the loss are highly concentrated on a low-dimensional subspace. This is why we can optimize a tiny fraction of parameters and achieve the same or better task performance.
  \item \textbf{Subspace Alignment and Trajectory Geometry}: The activation vectors $x$ represent coordinates on the manifold. The base weight matrix $W_0$ maps these coordinates to the next representation space. The update matrix $\Delta W = \frac{\alpha}{r} B A$ acts as a linear projection operator that projects the activations into a tiny subspace of dimension $r$ via $A$, rotates/scales them in this bottleneck space, and then projects them back to the output space via $B$. This operation restricts the adaptation to a set of $r$ directional vectors. Instead of shifting the entire coordinate frame, the adapter only redirects the activation trajectory along the principal directions relevant to the downstream task.
  
  Geometrically, the column space of $\Delta W$ is spanned by the columns of $B$, which is an $r$-dimensional subspace of the output space $\mathbb{R}^{k}$. The row space of $\Delta W$ is spanned by the rows of $A$, which is an $r$-dimensional subspace of the input space $\mathbb{R}^{d}$. Thus, the adaptation updates are strictly constrained to these coordinates, preserving all other dimensions from degradation.
\end{itemize}

\subsubsection{Visualizing the Calculations}
We can visualize how the information bottleneck constrains the updates in the representation space during the forward and backward passes:

\begin{figure}[h]
\centering
\begin{adjustbox}{max width=0.9\linewidth}
\begin{tikzpicture}[node distance=1.0cm and 1.5cm, scale=0.95, every node/.style={transform shape}]
  \node[datblk] (input) {Input Space $x$\\\\( $d$ dimensions)};
  \node[grayblk, right=of input, minimum width=2.4cm] (project) {Compression\\\\$z = Ax$};
  \node[trnblk, right=of project, minimum width=2.4cm] (bottleneck) {Bottleneck Space $z$\\\\( $r$ dimensions)};
  \node[grayblk, right=of bottleneck, minimum width=2.4cm] (expand) {Expansion\\\\$B z$};
  \node[datblk, right=of expand, minimum width=2.4cm] (out_lora) {Adapter Output\\\\( $k$ dimensions)};
  
  \node[frzblk, below=1.2cm of bottleneck, minimum width=4cm] (base) {Frozen Base Path\\\\$W_0 x$};
  \node[sumnode, right=1.2cm of out_lora] (sum) {$+$};
  \node[outnrn, right=1.0cm of sum] (output) {Final output $h$};

  \draw[arr] (input) -- (project);
  \draw[arr] (project) -- (bottleneck);
  \draw[arr] (bottleneck) -- (expand);
  \draw[arr] (expand) -- (out_lora);
  \draw[arr] (out_lora) -- (sum);
  \draw[arr] (sum) -- (output);
  
  \draw[arr] (input.south) |- (base.west);
  \draw[arr] (base.east) -| (sum.south);
  
  \node[font=\footnotesize\itshape, align=center, above=0.3cm of bottleneck] {Compressed\\\\Information Canal};
\end{tikzpicture}
\end{adjustbox}
\caption{Visualization of the information bottleneck: high-dimensional input is projected into a low-dimensional canal where task-specific tuning forces are applied, then projected back.}
\end{figure}

\begin{itemize}
  \item \textbf{Forward Pass Compression}: The activation vector $x \in \mathbb{R}^{d \times 1}$ is projected down through a narrow bottleneck of rank $r$. Since $r \ll \min(d, k)$, the projection filters out dimensions irrelevant to the current task, creating a compressed task representation $z \in \mathbb{R}^{r \times 1}$.
  \item \textbf{Backward Pass Gradient Flow}: Error gradients $\nabla_h \mathcal{L}$ flow backwards. Because $W_0$ is frozen and lacks a gradient buffer, no gradient update calculations occur on the base weight matrix. Instead, the gradients flow down the low-rank adapters $B$ and $A$, computing updates specifically for the low-rank subspace and avoiding parameter updates across the whole network.
\end{itemize}

\subsection{The Physical and Philosophical Foundations of LoRA}

\subsubsection{Philosophical Foundations: Occam's Razor and the Manifold Hypothesis}
The success of Low-Rank Adaptation (LoRA) is not merely an empirical coincidence, but is deeply rooted in two foundational philosophical principles of complexity. We analyze each principle through a dialectic framework of Arguments, Counter-Arguments, and Rebuttals:

\paragraph{1. The Manifold Hypothesis}
\begin{itemize}
  \item \textbf{Argument}: This hypothesis posits that high-dimensional data (such as the activations in a neural network) actually lies on a much lower-dimensional, highly structured manifold. Pretraining constructs this manifold by mapping the topological structure of natural language.
  
  In the context of Large Language Models, the embedding space has an extrinsic dimension of d = 4096 (e.g., $d_{\text{model}} = 4096$). However, the set of all grammatically correct and semantically coherent sentences does not fill this space uniformly; instead, it is confined to a lower-dimensional manifold. When we adapt the model to a downstream task (e.g., teaching it Python syntax or chat style), we do not need to construct a new manifold. We only need to compute a small trajectory or translation vector along the existing surface. This translation vector corresponds to a tangent vector in the manifold's tangent bundle. By confining our search to this tangent space, we avoid distorting the global topology of the language manifold, preventing catastrophic forgetting.
  
  \item \textbf{Counter-Argument}: If the downstream task represents an extreme domain shift (e.g., shifting from English literature to highly specialized biochemistry nomenclature or low-level assembly code), the target representations may lie completely outside the manifold mapped during pretraining. In such cases, the translation vector cannot lie along the existing manifold surface, and restricting the updates to a low-rank tangent space will lead to underfitting. The model would fail to acquire the new features because the required coordinates do not exist in the pretrained column space.
  
  \item \textbf{Rebuttal / Synthesis}: While severe domain shifts can strain low-rank constraints, empirical evidence shows that even highly specialized knowledge maps back to general-purpose linguistic structures (syntax, logic, context). Furthermore, if the target lies slightly off-manifold, a small rank $r$ (e.g., 8 to 16) is still sufficient to project the input into the nearest point on the task manifold. If needed, the user can scale the rank $r$ to increase the degrees of freedom, which mathematically expands the coordinate projection space, allowing the model to capture features that lie slightly outside the pretraining manifold while still regularizing the update to prevent arbitrary drift.
\end{itemize}

\paragraph{2. Occam's Razor (Lex Parsimoniae)}
\begin{itemize}
  \item \textbf{Argument}: Plurality should not be posited without necessity. Full fine-tuning searches an extremely large, redundant parameter space ($\mathbb{R}^{k \times d}$), which contains millions of redundant degrees of freedom. This massive parameter space allows the model to find arbitrary, highly complex fitting functions that overfit the downstream training data and overwrite pretrained general knowledge.
  
  By restricting the learning search to a compressed subspace ($\\mathbb{R}^{r \\times (d+k)}$), LoRA acts as a mathematical regularization agent. It forces the model to find the simplest, most parsimonious path of adaptation. This acts as a geometric constraint (projection regularization), restricting updates to only the principal directions of adaptation, thereby preserving the model's core foundational capabilities.
  
  \item \textbf{Counter-Argument}: By restricting the parameters to a low-rank product $BA$, we might force the model to find an oversimplified adaptation path that fails to capture the subtle, non-linear correlations required for complex reasoning tasks. Furthermore, Occam's Razor assumes that simpler models generalize better, but in deep learning, the phenomenon of "double descent" shows that highly overparameterized models can generalize better than simpler ones by finding smooth interpolating functions in high-dimensional spaces.
  
  \item \textbf{Rebuttal / Synthesis}: Double descent occurs when the entire network is overparameterized and trained on vast, diverse datasets where it can interpolate smoothly. However, in downstream fine-tuning, the dataset is small and highly focused. In this regime, overparameterization leads to catastrophic forgetting and representation collapse. Constraining the updates to a low-rank subspace restricts the model from modifying the general-purpose features, while allowing sufficient capacity to learn the task. Thus, LoRA combines the benefits of high-dimensional pretraining (smooth interpolation on general data) with low-dimensional adaptation (strict regularization on task-specific data).
\end{itemize}

\subsubsection{Internal Mechanics and Representation Space Geometry}
To understand why LoRA works under the hood, we must analyze the geometry of representation space and the intrinsic dimension of neural networks:
\begin{itemize}
  \item \textbf{Intrinsic vs. Extrinsic Dimension}: A pretrained model has a very high extrinsic dimension (the number of weights $W_0 \in \mathbb{R}^{k \times d}$). However, empirical work by Aghajanyan et al. (2020) demonstrated that the parameter space of pretrained language models has a much lower \textit{intrinsic dimension}. That is, the directions in weight space that significantly affect the loss are highly concentrated on a low-dimensional subspace. This is why we can optimize a tiny fraction of parameters and achieve the same or better task performance.
  \item \textbf{Subspace Alignment and Trajectory Geometry}: The activation vectors $x$ represent coordinates on the manifold. The base weight matrix $W_0$ maps these coordinates to the next representation space. The update matrix $\Delta W = \frac{\alpha}{r} B A$ acts as a linear projection operator that projects the activations into a tiny subspace of dimension $r$ via $A$, rotates/scales them in this bottleneck space, and then projects them back to the output space via $B$. This operation restricts the adaptation to a set of $r$ directional vectors. Instead of shifting the entire coordinate frame, the adapter only redirects the activation trajectory along the principal directions relevant to the downstream task.
  
  Geometrically, the column space of $\Delta W$ is spanned by the columns of $B$, which is an $r$-dimensional subspace of the output space $\mathbb{R}^{k}$. The row space of $\Delta W$ is spanned by the rows of $A$, which is an $r$-dimensional subspace of the input space $\mathbb{R}^{d}$. Thus, the adaptation updates are strictly constrained to these coordinates, preserving all other dimensions from degradation.
\end{itemize}

\subsubsection{Visualizing the Calculations}
We can visualize how the information bottleneck constrains the updates in the representation space during the forward and backward passes:

\begin{figure}[h]
\centering
\begin{adjustbox}{max width=0.9\linewidth}
\begin{tikzpicture}[node distance=1.0cm and 1.5cm, scale=0.95, every node/.style={transform shape}]
  \node[datblk] (input) {Input Space $x$\\\\( $d$ dimensions)};
  \node[grayblk, right=of input, minimum width=2.4cm] (project) {Compression\\\\$z = Ax$};
  \node[trnblk, right=of project, minimum width=2.4cm] (bottleneck) {Bottleneck Space $z$\\\\( $r$ dimensions)};
  \node[grayblk, right=of bottleneck, minimum width=2.4cm] (expand) {Expansion\\\\$B z$};
  \node[datblk, right=of expand, minimum width=2.4cm] (out_lora) {Adapter Output\\\\( $k$ dimensions)};
  
  \node[frzblk, below=1.2cm of bottleneck, minimum width=4cm] (base) {Frozen Base Path\\\\$W_0 x$};
  \node[sumnode, right=1.2cm of out_lora] (sum) {$+$};
  \node[outnrn, right=1.0cm of sum] (output) {Final output $h$};

  \draw[arr] (input) -- (project);
  \draw[arr] (project) -- (bottleneck);
  \draw[arr] (bottleneck) -- (expand);
  \draw[arr] (expand) -- (out_lora);
  \draw[arr] (out_lora) -- (sum);
  \draw[arr] (sum) -- (output);
  
  \draw[arr] (input.south) |- (base.west);
  \draw[arr] (base.east) -| (sum.south);
  
  \node[font=\footnotesize\itshape, align=center, above=0.3cm of bottleneck] {Compressed\\\\Information Canal};
\end{tikzpicture}
\end{adjustbox}
\caption{Visualization of the information bottleneck: high-dimensional input is projected into a low-dimensional canal where task-specific tuning forces are applied, then projected back.}
\end{figure}

\begin{itemize}
  \item \textbf{Forward Pass Compression}: The activation vector $x \in \mathbb{R}^{d \times 1}$ is projected down through a narrow bottleneck of rank $r$. Since $r \ll \min(d, k)$, the projection filters out dimensions irrelevant to the current task, creating a compressed task representation $z \in \mathbb{R}^{r \times 1}$.
  \item \textbf{Backward Pass Gradient Flow}: Error gradients $\nabla_h \mathcal{L}$ flow backwards. Because $W_0$ is frozen and lacks a gradient buffer, no gradient update calculations occur on the base weight matrix. Instead, the gradients flow down the low-rank adapters $B$ and $A$, computing updates specifically for the low-rank subspace and avoiding parameter updates across the whole network.
\end{itemize}

\subsection{The Physical and Philosophical Foundations of LoRA}

\subsubsection{Philosophical Foundations: Occam's Razor and the Manifold Hypothesis}
The success of Low-Rank Adaptation (LoRA) is not merely an empirical coincidence, but is deeply rooted in two foundational philosophical principles of complexity. We analyze each principle through a dialectic framework of Arguments, Counter-Arguments, and Rebuttals:

\paragraph{1. The Manifold Hypothesis}
\begin{itemize}
  \item \textbf{Argument}: This hypothesis posits that high-dimensional data (such as the activations in a neural network) actually lies on a much lower-dimensional, highly structured manifold. Pretraining constructs this manifold by mapping the topological structure of natural language.
  
  In the context of Large Language Models, the embedding space has an extrinsic dimension of d = 4096 (e.g., $d_{\text{model}} = 4096$). However, the set of all grammatically correct and semantically coherent sentences does not fill this space uniformly; instead, it is confined to a lower-dimensional manifold. When we adapt the model to a downstream task (e.g., teaching it Python syntax or chat style), we do not need to construct a new manifold. We only need to compute a small trajectory or translation vector along the existing surface. This translation vector corresponds to a tangent vector in the manifold's tangent bundle. By confining our search to this tangent space, we avoid distorting the global topology of the language manifold, preventing catastrophic forgetting.
  
  \item \textbf{Counter-Argument}: If the downstream task represents an extreme domain shift (e.g., shifting from English literature to highly specialized biochemistry nomenclature or low-level assembly code), the target representations may lie completely outside the manifold mapped during pretraining. In such cases, the translation vector cannot lie along the existing manifold surface, and restricting the updates to a low-rank tangent space will lead to underfitting. The model would fail to acquire the new features because the required coordinates do not exist in the pretrained column space.
  
  \item \textbf{Rebuttal / Synthesis}: While severe domain shifts can strain low-rank constraints, empirical evidence shows that even highly specialized knowledge maps back to general-purpose linguistic structures (syntax, logic, context). Furthermore, if the target lies slightly off-manifold, a small rank $r$ (e.g., 8 to 16) is still sufficient to project the input into the nearest point on the task manifold. If needed, the user can scale the rank $r$ to increase the degrees of freedom, which mathematically expands the coordinate projection space, allowing the model to capture features that lie slightly outside the pretraining manifold while still regularizing the update to prevent arbitrary drift.
\end{itemize}

\paragraph{2. Occam's Razor (Lex Parsimoniae)}
\begin{itemize}
  \item \textbf{Argument}: Plurality should not be posited without necessity. Full fine-tuning searches an extremely large, redundant parameter space ($\mathbb{R}^{k \times d}$), which contains millions of redundant degrees of freedom. This massive parameter space allows the model to find arbitrary, highly complex fitting functions that overfit the downstream training data and overwrite pretrained general knowledge.
  
  By restricting the learning search to a compressed subspace ($\mathbb{R}^{r \times (d+k)}$), LoRA acts as a mathematical regularization agent. It forces the model to find the simplest, most parsimonious path of adaptation. This acts as a geometric constraint (projection regularization), restricting updates to only the principal directions of adaptation, thereby preserving the model's core foundational capabilities.
  
  \item \textbf{Counter-Argument}: By restricting the parameters to a low-rank product $BA$, we might force the model to find an oversimplified adaptation path that fails to capture the subtle, non-linear correlations required for complex reasoning tasks. Furthermore, Occam's Razor assumes that simpler models generalize better, but in deep learning, the phenomenon of "double descent" shows that highly overparameterized models can generalize better than simpler ones by finding smooth interpolating functions in high-dimensional spaces.
  
  \item \textbf{Rebuttal / Synthesis}: Double descent occurs when the entire network is overparameterized and trained on vast, diverse datasets where it can interpolate smoothly. However, in downstream fine-tuning, the dataset is small and highly focused. In this regime, overparameterization leads to catastrophic forgetting and representation collapse. Constraining the updates to a low-rank subspace restricts the model from modifying the general-purpose features, while allowing sufficient capacity to learn the task. Thus, LoRA combines the benefits of high-dimensional pretraining (smooth interpolation on general data) with low-dimensional adaptation (strict regularization on task-specific data).
\end{itemize}

\subsubsection{Internal Mechanics and Representation Space Geometry}
To understand why LoRA works under the hood, we must analyze the geometry of representation space and the intrinsic dimension of neural networks:
\begin{itemize}
  \item \textbf{Intrinsic vs. Extrinsic Dimension}: A pretrained model has a very high extrinsic dimension (the number of weights $W_0 \in \mathbb{R}^{k \times d}$). However, empirical work by Aghajanyan et al. (2020) demonstrated that the parameter space of pretrained language models has a much lower \textit{intrinsic dimension}. That is, the directions in weight space that significantly affect the loss are highly concentrated on a low-dimensional subspace. This is why we can optimize a tiny fraction of parameters and achieve the same or better task performance.
  \item \textbf{Subspace Alignment and Trajectory Geometry}: The activation vectors $x$ represent coordinates on the manifold. The base weight matrix $W_0$ maps these coordinates to the next representation space. The update matrix $\Delta W = \frac{\alpha}{r} B A$ acts as a linear projection operator that projects the activations into a tiny subspace of dimension $r$ via $A$, rotates/scales them in this bottleneck space, and then projects them back to the output space via $B$. This operation restricts the adaptation to a set of $r$ directional vectors. Instead of shifting the entire coordinate frame, the adapter only redirects the activation trajectory along the principal directions relevant to the downstream task.
  
  Geometrically, the column space of $\Delta W$ is spanned by the columns of $B$, which is an $r$-dimensional subspace of the output space $\mathbb{R}^{k}$. The row space of $\Delta W$ is spanned by the rows of $A$, which is an $r$-dimensional subspace of the input space $\mathbb{R}^{d}$. Thus, the adaptation updates are strictly constrained to these coordinates, preserving all other dimensions from degradation.
\end{itemize}

\subsubsection{Visualizing the Calculations}
We can visualize how the information bottleneck constrains the updates in the representation space during the forward and backward passes:

\begin{figure}[h]
\centering
\begin{adjustbox}{max width=0.9\linewidth}
\begin{tikzpicture}[node distance=1.0cm and 1.5cm, scale=0.95, every node/.style={transform shape}]
  \node[datblk] (input) {Input Space $x$\\\\( $d$ dimensions)};
  \node[grayblk, right=of input, minimum width=2.4cm] (project) {Compression\\\\$z = Ax$};
  \node[trnblk, right=of project, minimum width=2.4cm] (bottleneck) {Bottleneck Space $z$\\\\( $r$ dimensions)};
  \node[grayblk, right=of bottleneck, minimum width=2.4cm] (expand) {Expansion\\\\$B z$};
  \node[datblk, right=of expand, minimum width=2.4cm] (out_lora) {Adapter Output\\\\( $k$ dimensions)};
  
  \node[frzblk, below=1.2cm of bottleneck, minimum width=4cm] (base) {Frozen Base Path\\\\$W_0 x$};
  \node[sumnode, right=1.2cm of out_lora] (sum) {$+$};
  \node[outnrn, right=1.0cm of sum] (output) {Final output $h$};

  \draw[arr] (input) -- (project);
  \draw[arr] (project) -- (bottleneck);
  \draw[arr] (bottleneck) -- (expand);
  \draw[arr] (expand) -- (out_lora);
  \draw[arr] (out_lora) -- (sum);
  \draw[arr] (sum) -- (output);
  
  \draw[arr] (input.south) |- (base.west);
  \draw[arr] (base.east) -| (sum.south);
  
  \node[font=\footnotesize\itshape, align=center, above=0.3cm of bottleneck] {Compressed\\\\Information Canal};
\end{tikzpicture}
\end{adjustbox}
\caption{Visualization of the information bottleneck: high-dimensional input is projected into a low-dimensional canal where task-specific tuning forces are applied, then projected back.}
\end{figure}

\begin{itemize}
  \item \textbf{Forward Pass Compression}: The activation vector $x \in \mathbb{R}^{d \times 1}$ is projected down through a narrow bottleneck of rank $r$. Since $r \ll \min(d, k)$, the projection filters out dimensions irrelevant to the current task, creating a compressed task representation $z \in \mathbb{R}^{r \times 1}$.
  \item \textbf{Backward Pass Gradient Flow}: Error gradients $\nabla_h \mathcal{L}$ flow backwards. Because $W_0$ is frozen and lacks a gradient buffer, no gradient update calculations occur on the base weight matrix. Instead, the gradients flow down the low-rank adapters $B$ and $A$, computing updates specifically for the low-rank subspace and avoiding parameter updates across the whole network.
\end{itemize}

\subsection{The Physical and Philosophical Foundations of LoRA}

\subsubsection{Philosophical Foundations: Occam's Razor and the Manifold Hypothesis}
The success of Low-Rank Adaptation (LoRA) is not merely an empirical coincidence, but is deeply rooted in two foundational philosophical principles of complexity. We analyze each principle through a dialectic framework of Arguments, Counter-Arguments, and Rebuttals:

\paragraph{1. The Manifold Hypothesis}
\begin{itemize}
  \item \textbf{Argument}: This hypothesis posits that high-dimensional data (such as the activations in a neural network) actually lies on a much lower-dimensional, highly structured manifold. Pretraining constructs this manifold by mapping the topological structure of natural language.
  
  In the context of Large Language Models, the embedding space has an extrinsic dimension of $d = 4096$ (e.g., $d_{\text{model}} = 4096$). However, the set of all grammatically correct and semantically coherent sentences does not fill this space uniformly; instead, it is confined to a lower-dimensional manifold. When we adapt the model to a downstream task (e.g., teaching it Python syntax or chat style), we do not need to construct a new manifold. We only need to compute a small trajectory or translation vector along the existing surface. This translation vector corresponds to a tangent vector in the manifold's tangent bundle. By confining our search to this tangent space, we avoid distorting the global topology of the language manifold, preventing catastrophic forgetting.
  
  \item \textbf{Counter-Argument}: If the downstream task represents an extreme domain shift (e.g., shifting from English literature to highly specialized biochemistry nomenclature or low-level assembly code), the target representations may lie completely outside the manifold mapped during pretraining. In such cases, the translation vector cannot lie along the existing manifold surface, and restricting the updates to a low-rank tangent space will lead to underfitting. The model would fail to acquire the new features because the required coordinates do not exist in the pretrained column space.
  
  \item \textbf{Rebuttal / Synthesis}: While severe domain shifts can strain low-rank constraints, empirical evidence shows that even highly specialized knowledge maps back to general-purpose linguistic structures (syntax, logic, context). Furthermore, if the target lies slightly off-manifold, a small rank $r$ (e.g., 8 to 16) is still sufficient to project the input into the nearest point on the task manifold. If needed, the user can scale the rank $r$ to increase the degrees of freedom, which mathematically expands the coordinate projection space, allowing the model to capture features that lie slightly outside the pretraining manifold while still regularizing the update to prevent arbitrary drift.
\end{itemize}

\paragraph{2. Occam's Razor (Lex Parsimoniae)}
\begin{itemize}
  \item \textbf{Argument}: Plurality should not be posited without necessity. Full fine-tuning searches an extremely large, redundant parameter space ($\mathbb{R}^{k \times d}$), which contains millions of redundant degrees of freedom. This massive parameter space allows the model to find arbitrary, highly complex fitting functions that overfit the downstream training data and overwrite pretrained general knowledge.
  
  By restricting the learning search to a compressed subspace ($\mathbb{R}^{r \times (d+k)}$), LoRA acts as a mathematical regularization agent. It forces the model to find the simplest, most parsimonious path of adaptation. This acts as a geometric constraint (projection regularization), restricting updates to only the principal directions of adaptation, thereby preserving the model's core foundational capabilities.
  
  \item \textbf{Counter-Argument}: By restricting the parameters to a low-rank product $BA$, we might force the model to find an oversimplified adaptation path that fails to capture the subtle, non-linear correlations required for complex reasoning tasks. Furthermore, Occam's Razor assumes that simpler models generalize better, but in deep learning, the phenomenon of "double descent" shows that highly overparameterized models can generalize better than simpler ones by finding smooth interpolating functions in high-dimensional spaces.
  
  \item \textbf{Rebuttal / Synthesis}: Double descent occurs when the entire network is overparameterized and trained on vast, diverse datasets where it can interpolate smoothly. However, in downstream fine-tuning, the dataset is small and highly focused. In this regime, overparameterization leads to catastrophic forgetting and representation collapse. Constraining the updates to a low-rank subspace restricts the model from modifying the general-purpose features, while allowing sufficient capacity to learn the task. Thus, LoRA combines the benefits of high-dimensional pretraining (smooth interpolation on general data) with low-dimensional adaptation (strict regularization on task-specific data).
\end{itemize}

\subsubsection{Internal Mechanics and Representation Space Geometry}
To understand why LoRA works under the hood, we must analyze the geometry of representation space and the intrinsic dimension of neural networks:
\begin{itemize}
  \item \textbf{Intrinsic vs. Extrinsic Dimension}: A pretrained model has a very high extrinsic dimension (the number of weights $W_0 \in \mathbb{R}^{k \times d}$). However, empirical work by Aghajanyan et al. (2020) demonstrated that the parameter space of pretrained language models has a much lower \textit{intrinsic dimension}. That is, the directions in weight space that significantly affect the loss are highly concentrated on a low-dimensional subspace. This is why we can optimize a tiny fraction of parameters and achieve the same or better task performance.
  \item \textbf{Subspace Alignment and Trajectory Geometry}: The activation vectors $x$ represent coordinates on the manifold. The base weight matrix $W_0$ maps these coordinates to the next representation space. The update matrix $\Delta W = \frac{\alpha}{r} B A$ acts as a linear projection operator that projects the activations into a tiny subspace of dimension $r$ via $A$, rotates/scales them in this bottleneck space, and then projects them back to the output space via $B$. This operation restricts the adaptation to a set of $r$ directional vectors. Instead of shifting the entire coordinate frame, the adapter only redirects the activation trajectory along the principal directions relevant to the downstream task.
  
  Geometrically, the column space of $\Delta W$ is spanned by the columns of $B$, which is an $r$-dimensional subspace of the output space $\mathbb{R}^{k}$. The row space of $\Delta W$ is spanned by the rows of $A$, which is an $r$-dimensional subspace of the input space $\mathbb{R}^{d}$. Thus, the adaptation updates are strictly constrained to these coordinates, preserving all other dimensions from degradation.
\end{itemize}

\subsubsection{Visualizing the Calculations}
We can visualize how the information bottleneck constrains the updates in the representation space during the forward and backward passes:

\begin{figure}[h]
\centering
\begin{adjustbox}{max width=0.9\linewidth}
\begin{tikzpicture}[node distance=1.0cm and 1.5cm]
  \node[datblk] (input) {Input Space $x$\\\\( $d$ dimensions)};
  \node[grayblk, right=of input, minimum width=2.4cm] (project) {Compression\\\\$z = Ax$};
  \node[trnblk, right=of project, minimum width=2.4cm] (bottleneck) {Bottleneck Space $z$\\\\( $r$ dimensions)};
  \node[grayblk, right=of bottleneck, minimum width=2.4cm] (expand) {Expansion\\\\$B z$};
  \node[datblk, right=of expand, minimum width=2.4cm] (out_lora) {Adapter Output\\\\( $k$ dimensions)};
  
  \node[frzblk, below=1.2cm of bottleneck, minimum width=4cm] (base) {Frozen Base Path\\\\$W_0 x$};
  \node[sumnode, right=1.2cm of out_lora] (sum) {$+$};
  \node[outnrn, right=1.0cm of sum] (output) {Final output $h$};

  \draw[arr] (input) -- (project);
  \draw[arr] (project) -- (bottleneck);
  \draw[arr] (bottleneck) -- (expand);
  \draw[arr] (expand) -- (out_lora);
  \draw[arr] (out_lora) -- (sum);
  \draw[arr] (sum) -- (output);
  
  \draw[arr] (input.south) |- (base.west);
  \draw[arr] (base.east) -| (sum.south);
  
  \node[font=\footnotesize\itshape, align=center, above=0.3cm of bottleneck] {Compressed\\\\Information Canal};
\end{tikzpicture}
\end{adjustbox}
\caption{Visualization of the information bottleneck: high-dimensional input is projected into a low-dimensional canal where task-specific tuning forces are applied, then projected back.}
\end{figure}

\begin{itemize}
  \item \textbf{Forward Pass Compression}: The activation vector $x \in \mathbb{R}^{d \times 1}$ is projected down through a narrow bottleneck of rank $r$. Since $r \ll \min(d, k)$, the projection filters out dimensions irrelevant to the current task, creating a compressed task representation $z \in \mathbb{R}^{r \times 1}$.
  \item \textbf{Backward Pass Gradient Flow}: Error gradients $\nabla_h \mathcal{L}$ flow backwards. Because $W_0$ is frozen and lacks a gradient buffer, no gradient update calculations occur on the base weight matrix. Instead, the gradients flow down the low-rank adapters $B$ and $A$, computing updates specifically for the low-rank subspace and avoiding parameter updates across the whole network.
\end{itemize}

\subsection{The Physical and Philosophical Foundations of LoRA}

\subsubsection{Philosophical Foundations: Occam's Razor and the Manifold Hypothesis}
The success of Low-Rank Adaptation (LoRA) is not merely an empirical coincidence, but is deeply rooted in two foundational philosophical principles of complexity. We analyze each principle through a dialectic framework of Arguments, Counter-Arguments, and Rebuttals:

\paragraph{1. The Manifold Hypothesis}
\begin{itemize}
  \item \textbf{Argument}: This hypothesis posits that high-dimensional data (such as the activations in a neural network) actually lies on a much lower-dimensional, highly structured manifold. Pretraining constructs this manifold by mapping the topological structure of natural language.
  
  In the context of Large Language Models, the embedding space has an extrinsic dimension of $d = 4096$ (e.g., $d_{\text{model}} = 4096$). However, the set of all grammatically correct and semantically coherent sentences does not fill this space uniformly; instead, it is confined to a lower-dimensional manifold. When we adapt the model to a downstream task (e.g., teaching it Python syntax or chat style), we do not need to construct a new manifold. We only need to compute a small trajectory or translation vector along the existing surface. This translation vector corresponds to a tangent vector in the manifold's tangent bundle. By confining our search to this tangent space, we avoid distorting the global topology of the language manifold, preventing catastrophic forgetting.
  
  \item \textbf{Counter-Argument}: If the downstream task represents an extreme domain shift (e.g., shifting from English literature to highly specialized biochemistry nomenclature or low-level assembly code), the target representations may lie completely outside the manifold mapped during pretraining. In such cases, the translation vector cannot lie along the existing manifold surface, and restricting the updates to a low-rank tangent space will lead to underfitting. The model would fail to acquire the new features because the required coordinates do not exist in the pretrained column space.
  
  \item \textbf{Rebuttal / Synthesis}: While severe domain shifts can strain low-rank constraints, empirical evidence shows that even highly specialized knowledge maps back to general-purpose linguistic structures (syntax, logic, context). Furthermore, if the target lies slightly off-manifold, a small rank $r$ (e.g., 8 to 16) is still sufficient to project the input into the nearest point on the task manifold. If needed, the user can scale the rank $r$ to increase the degrees of freedom, which mathematically expands the coordinate projection space, allowing the model to capture features that lie slightly outside the pretraining manifold while still regularizing the update to prevent arbitrary drift.
\end{itemize}

\paragraph{2. Occam's Razor (Lex Parsimoniae)}
\begin{itemize}
  \item \textbf{Argument}: Plurality should not be posited without necessity. Full fine-tuning searches an extremely large, redundant parameter space ($\mathbb{R}^{k \times d}$), which contains millions of redundant degrees of freedom. This massive parameter space allows the model to find arbitrary, highly complex fitting functions that overfit the downstream training data and overwrite pretrained general knowledge.
  
  By restricting the learning search to a compressed subspace ($\mathbb{R}^{r \times (d+k)}$), LoRA acts as a mathematical regularization agent. It forces the model to find the simplest, most parsimonious path of adaptation. This acts as a geometric constraint (projection regularization), restricting updates to only the principal directions of adaptation, thereby preserving the model's core foundational capabilities.
  
  \item \textbf{Counter-Argument}: By restricting the parameters to a low-rank product $BA$, we might force the model to find an oversimplified adaptation path that fails to capture the subtle, non-linear correlations required for complex reasoning tasks. Furthermore, Occam's Razor assumes that simpler models generalize better, but in deep learning, the phenomenon of "double descent" shows that highly overparameterized models can generalize better than simpler ones by finding smooth interpolating functions in high-dimensional spaces.
  
  \item \textbf{Rebuttal / Synthesis}: Double descent occurs when the entire network is overparameterized and trained on vast, diverse datasets where it can interpolate smoothly. However, in downstream fine-tuning, the dataset is small and highly focused. In this regime, overparameterization leads to catastrophic forgetting and representation collapse. Constraining the updates to a low-rank subspace restricts the model from modifying the general-purpose features, while allowing sufficient capacity to learn the task. Thus, LoRA combines the benefits of high-dimensional pretraining (smooth interpolation on general data) with low-dimensional adaptation (strict regularization on task-specific data).
\end{itemize}

\subsubsection{Internal Mechanics and Representation Space Geometry}
To understand why LoRA works under the hood, we must analyze the geometry of representation space and the intrinsic dimension of neural networks:
\begin{itemize}
  \item \textbf{Intrinsic vs. Extrinsic Dimension}: A pretrained model has a very high extrinsic dimension (the number of weights $W_0 \in \mathbb{R}^{k \times d}$). However, empirical work by Aghajanyan et al. (2020) demonstrated that the parameter space of pretrained language models has a much lower \textit{intrinsic dimension}. That is, the directions in weight space that significantly affect the loss are highly concentrated on a low-dimensional subspace. This is why we can optimize a tiny fraction of parameters and achieve the same or better task performance.
  \item \textbf{Subspace Alignment and Trajectory Geometry}: The activation vectors $x$ represent coordinates on the manifold. The base weight matrix $W_0$ maps these coordinates to the next representation space. The update matrix $\Delta W = \frac{\alpha}{r} B A$ acts as a linear projection operator that projects the activations into a tiny subspace of dimension $r$ via $A$, rotates/scales them in this bottleneck space, and then projects them back to the output space via $B$. This operation restricts the adaptation to a set of $r$ directional vectors. Instead of shifting the entire coordinate frame, the adapter only redirects the activation trajectory along the principal directions relevant to the downstream task.
  
  Geometrically, the column space of $\Delta W$ is spanned by the columns of $B$, which is an $r$-dimensional subspace of the output space $\mathbb{R}^{k}$. The row space of $\Delta W$ is spanned by the rows of $A$, which is an $r$-dimensional subspace of the input space $\mathbb{R}^{d}$. Thus, the adaptation updates are strictly constrained to these coordinates, preserving all other dimensions from degradation.
\end{itemize}

\subsubsection{Visualizing the Calculations}
We can visualize how the information bottleneck constrains the updates in the representation space during the forward and backward passes:

\begin{figure}[h]
\centering
\begin{adjustbox}{max width=0.9\linewidth}
\begin{tikzpicture}[node distance=1.0cm and 1.5cm]
  \node[datblk] (input) {Input Space $x$\\( $d$ dimensions)};
  \node[grayblk, right=of input, minimum width=2.4cm] (project) {Compression\\$z = Ax$};
  \node[trnblk, right=of project, minimum width=2.4cm] (bottleneck) {Bottleneck Space $z$\\( $r$ dimensions)};
  \node[grayblk, right=of bottleneck, minimum width=2.4cm] (expand) {Expansion\\$B z$};
  \node[datblk, right=of expand, minimum width=2.4cm] (out_lora) {Adapter Output\\( $k$ dimensions)};
  
  \node[frzblk, below=1.2cm of bottleneck, minimum width=4cm] (base) {Frozen Base Path\\$W_0 x$};
  \node[sumnode, right=1.2cm of out_lora] (sum) {$+$};
  \node[outnrn, right=1.0cm of sum] (output) {Final output $h$};

  \draw[arr] (input) -- (project);
  \draw[arr] (project) -- (bottleneck);
  \draw[arr] (bottleneck) -- (expand);
  \draw[arr] (expand) -- (out_lora);
  \draw[arr] (out_lora) -- (sum);
  \draw[arr] (sum) -- (output);
  
  \draw[arr] (input.south) |- (base.west);
  \draw[arr] (base.east) -| (sum.south);
  
  \node[font=\footnotesize\itshape, align=center, above=0.3cm of bottleneck] {Compressed\\Information Canal};
\end{tikzpicture}
\end{adjustbox}
\caption{Visualization of the information bottleneck: high-dimensional input is projected into a low-dimensional canal where task-specific tuning forces are applied, then projected back.}
\end{figure}

\begin{itemize}
  \item \textbf{Forward Pass Compression}: The activation vector $x \in \mathbb{R}^{d \times 1}$ is projected down through a narrow bottleneck of rank $r$. Since $r \ll \min(d, k)$, the projection filters out dimensions irrelevant to the current task, creating a compressed task representation $z \in \mathbb{R}^{r \times 1}$.
  \item \textbf{Backward Pass Gradient Flow}: Error gradients $\nabla_h \mathcal{L}$ flow backwards. Because $W_0$ is frozen and lacks a gradient buffer, no gradient update calculations occur on the base weight matrix. Instead, the gradients flow down the low-rank adapters $B$ and $A$, computing updates specifically for the low-rank subspace and avoiding parameter updates across the whole network.
\end{itemize}

\subsection{LoRA CLRS-Style Algorithms}
This section contains the CLRS-style mathematical algorithms for the forward pass, manual backward pass, parameter merging, and AdamW optimizer update step.

\begin{figure}[h]
\centering
\begin{minipage}{0.8\linewidth}
\begin{tabbing}
\quad \= \quad \= \quad \= \kill
\textsc{Lora-Forward}($x, W_0, A, B, \alpha, r$) \\
\> $h_{\text{base}} = W_0 x$ \\
\> $z = A x$ \\
\> $z_{\text{out}} = B z$ \\
\> $h = h_{\text{base}} + \frac{\alpha}{r} z_{\text{out}}$ \\
\> \textbf{return} $h$
\end{tabbing}
\end{minipage}
\caption{\textsc{Lora-Forward} algorithm: computes base output and scaled low-rank adapter output.}
\end{figure}

\subsubsection{Analysis of \textsc{Lora-Forward}}
\begin{itemize}
  \item \textbf{What is achieved}: Computes the layer output $h$ by parallelizing the computation of the frozen base weight path ($W_0 x$) and the trainable low-rank adapter path ($B(A x)$), and scaling the adapter's contribution.
  \item \textbf{What internally changed}: Instantiates intermediate activation tensors $z = Ax \in \mathbb{R}^{r \times 1}$ and $z_{\text{out}} = Bz \in \mathbb{R}^{k \times 1}$ in memory. During training, these tensors are kept in GPU memory (VRAM) because their values are required to compute the gradients in the backward pass.
  \item \textbf{User-modifiable parameters \& impact}:
    \begin{itemize}
      \item \textbf{Rank $r$} (hyperparameter): Modifies the bottleneck dimension. A larger $r$ (e.g., 32) captures more task-specific information but increases memory requirements and latency. A smaller $r$ (e.g., 8) is highly memory-efficient but might underfit complex tasks.
      \item \textbf{Scaling constant $\alpha$} (hyperparameter): Adjusts the strength of the adapter's adjustments. Increasing $\alpha$ scales up the adapter's updates. Keep this constant when tuning $r$ to stabilize learning.
    \end{itemize}
\end{itemize}

\vspace{1.5em}

\begin{figure}[h]
\centering
\begin{minipage}{0.8\linewidth}
\begin{tabbing}
\quad \= \quad \= \quad \= \kill
\textsc{Lora-Backward}($x, A, B, \nabla_h \mathcal{L}, \alpha, r$) \\
\> $z = A x$ \\
\> $\nabla_B \mathcal{L} = \frac{\alpha}{r} (\nabla_h \mathcal{L}) z^\top$ \\
\> $\nabla_z \mathcal{L} = \frac{\alpha}{r} B^\top \nabla_h \mathcal{L}$ \\
\> $\nabla_A \mathcal{L} = (\nabla_z \mathcal{L}) x^\top$ \\
\> $\nabla_{x,\text{adapters}} \mathcal{L} = A^\top \nabla_z \mathcal{L}$ \\
\> \textbf{return} $\nabla_A \mathcal{L}, \nabla_B \mathcal{L}, \nabla_{x,\text{adapters}} \mathcal{L}$
\end{tabbing}
\end{minipage}
\caption{\textsc{Lora-Backward} algorithm: derives gradients for adapters and backpropagated input gradient.}
\end{figure}

\subsubsection{Analysis of \textsc{Lora-Backward}}
\begin{itemize}
  \item \textbf{What is achieved}: Calculates the exact gradients $\nabla_A \mathcal{L}$ and $\nabla_B \mathcal{L}$ for the adapter weights, and computes the backpropagated input gradient $\nabla_{x,\text{adapters}} \mathcal{L}$ to pass back to previous layers.
  \item \textbf{What internally changed}: Populates the gradient buffers \texttt{A.grad} and \texttt{B.grad} in VRAM. The frozen weights $W_0$ do not have gradient buffers, saving $d \times k \times 2$ or $4$ bytes of memory.
  \item \textbf{User-modifiable parameters \& impact}:
    \begin{itemize}
      \item \textbf{Rank $r$ and Scaling $\alpha$}: Directly scale the gradients via the factor $\frac{\alpha}{r}$. Modifying them scales the gradient step size, affecting convergence stability.
    \end{itemize}
\end{itemize}

\newpage

\begin{figure}[h]
\centering
\begin{minipage}{0.8\linewidth}
\begin{tabbing}
\quad \= \quad \= \quad \= \kill
\textsc{Lora-Merge}($W_0, A, B, \alpha, r, \text{\textit{mode}}$) \\
\> $\Delta W = \frac{\alpha}{r} B A$ \\
\> \textbf{if} $\text{\textit{mode}} == \text{``merge''}$ \\
\> \> $W = W_0 + \Delta W$ \\
\> \textbf{else if} $\text{\textit{mode}} == \text{``demerge''}$ \\
\> \> $W = W_0 - \Delta W$ \\
\> \textbf{return} $W$
\end{tabbing}
\end{minipage}
\caption{\textsc{Lora-Merge} algorithm: handles parameter folding and unfolding for zero-latency inference.}
\end{figure}

\subsubsection{Analysis of \textsc{Lora-Merge}}
\begin{itemize}
  \item \textbf{What is achieved}: Merges the low-rank adapters directly into the base weights ($W = W_0 \pm \Delta W$) for zero-latency inference, or subtracts them to restore the base weights.
  \item \textbf{What internally changed}: In-place updates are applied to the base weight tensor $W_0$ in memory. Once merged, the adapter weights $A$ and $B$ can be discarded, and the model behaves as a single unified feedforward network, eliminating the need to compute separate low-rank paths during inference.
  \item \textbf{User-modifiable parameters \& impact}:
    \begin{itemize}
      \item \textbf{Mode}: User can choose \texttt{"merge"} to prepare the model for high-throughput deployment, or \texttt{"demerge"} to unpack the adapters and swap in a different set of adapter weights.
    \end{itemize}
\end{itemize}

\vspace{1.5em}

\begin{figure}[h]
\centering
\begin{minipage}{0.8\linewidth}
\begin{tabbing}
\quad \= \quad \= \quad \= \quad \= \kill
\textsc{Lora-AdamW-Update}($\theta, g, m, v, \eta, \lambda, \beta_1, \beta_2, \epsilon, t$) \\
\> \textbf{for} \textbf{each} parameter matrix $P \in \theta$ \textbf{do} \\
\> \> $P = P - \eta \lambda P$ \\
\> \> $m_P = \beta_1 m_P + (1 - \beta_1) g_P$ \\
\> \> $v_P = \beta_2 v_P + (1 - \beta_2) g_P^2$ \\
\> \> $\hat{m}_P = m_P / (1 - \beta_1^t)$ \\
\> \> $\hat{v}_P = v_P / (1 - \beta_2^t)$ \\
\> \> $P = P - \eta \frac{\hat{m}_P}{\sqrt{\hat{v}_P} + \epsilon}$ \\
\> \textbf{return} $\theta, m, v$
\end{tabbing}
\end{minipage}
\caption{\textsc{Lora-AdamW-Update} algorithm: updates parameters with decoupled weight decay and momentum/variance tracking.}
\end{figure}

\subsubsection{Analysis of \textsc{Lora-AdamW-Update}}
\begin{itemize}
  \item \textbf{What is achieved}: Applies decoupled weight decay to the adapter parameters and updates their values using momentum and variance of gradients.
  \item \textbf{What internally changed}: Trainable parameters $A$ and $B$ are updated. Additionally, the first moment vectors $m_A, m_B$ and second moment vectors $v_A, v_B$ are updated in GPU memory. Because $W_0$ is frozen, no optimizer states are created or updated for $W_0$, resulting in huge VRAM savings.
  \item \textbf{User-modifiable parameters \& impact}:
    \begin{itemize}
      \item \textbf{Learning Rate $\eta$}: Controls the step size. High values speed up training but risk divergence; low values ensure stability but slow down convergence.
      \item \textbf{Weight Decay $\lambda$}: Regularizes the weights to prevent overfitting.
      \item \textbf{Momentum coefficients $\beta_1, \beta_2$}: Control the memory of past gradients.
    \end{itemize}
\end{itemize}


\newpage
%% ------------------------------------------------------------
\section{Quantized LoRA (QLoRA)}

\subsection{Mathematical Foundation}
QLoRA quantizes the base weight $W_0$ to 4-bit NormalFloat (NF4) and uses Double Quantization (DQ) to minimize memory requirements.
NF4 divides the standard normal distribution $\mathcal{N}(0, 1)$ into $2^k=16$ equal-area intervals. The quantization levels $q_i$ are calculated using the standard normal quantile function $Q_X$:
\begin{equation}
  q_i = \frac{1}{2} \left( Q_X\left(\frac{i}{2^k}\right) + Q_X\left(\frac{i+1}{2^k}\right) \right)
\end{equation}
The base weights are normalized to $[-1, 1]$ using block-wise scaling. Let the block size be $B=64$. For a block $W_b$, the scale is $c_b = \max(|W_b|)$. The quantized values are:
\begin{equation}
  q_t = \arg\min_{i \in \{0, \dots, 15\}} \left| \frac{W_{b,t}}{c_b} - q_i \right|
\end{equation}
Double Quantization (DQ) quantizes the scaling factors $c_b$ themselves to 8-bit floats ($c_b^{\text{FP8}}$) with a second scale block size of 256, reducing scale memory overhead from $32/64 = 0.5$ bits/weight to $8/64 + 32/(64 \times 256) \approx 0.127$ bits/weight.
During backpropagation, the NF4 weights are dequantized to FP16/BF16 to compute outputs, and gradients are computed only for the low-rank parameters:
\begin{equation}
  h = \text{dequant}(c_b, W^{\text{NF4}}) x + \frac{\alpha}{r} B A x
\end{equation}

\subsection{Architecture Flow Diagram}

\begin{figure}[h]
\centering
\begin{adjustbox}{max width=\linewidth}
\begin{tikzpicture}[node distance=0.7cm and 1.2cm]
  \node[datblk] (x) {\textbf{Input} \\ ($x$ (BF16))};
  \node[frzblk, right=of x] (nf4) {\textbf{$W_0^{\text{NF4}}$} \\ (4-bit frozen)};
  \node[grayblk, right=of nf4] (deq) {\textbf{Dequantise} \\ (to BF16)};
  \node[sumnode, right=of deq] (sum) {$\oplus$};
  \node[trnblk, above=0.8cm of sum] (lora) {\textbf{LoRA} \\ ($BA$ (BF16))};
  \node[mrgblk, right=of sum] (out) {\textbf{Output} \\ ($h$ (BF16))};

  \draw[arr] (x)--(nf4);
  \draw[arr] (nf4)--(deq);
  \draw[arr] (deq)--(sum);
  \draw[arr] (lora)--(sum);
  \draw[arr] (x.north) |- (lora.west);
  \draw[arr] (sum)--(out);

  \node[red,font=\tiny,text width=2.5cm,align=center,below=0.3cm of nf4] {4-bit: $\sim$4 GB\\vs BF16: $\sim$16 GB};
\end{tikzpicture}
\end{adjustbox}
\caption{QLoRA: 4-bit frozen base dequantised at runtime; LoRA adapters trained in BF16.}
\end{figure}

\subsection{State Transition Table}
\begin{table}[h]\centering
\begin{tabular}{lll}\toprule
\textbf{Property} & \textbf{Before} & \textbf{After}\\\midrule
$W_0$ precision & BF16 / FP16 & NF4 (4-bit)\\
GPU memory for 70B model & ${\sim}140$ GB & ${\sim}35$ GB\\
LoRA adapter precision & N/A & BF16\\
Training fidelity & Full & Near-full (straight-through)\\
Inference latency & Baseline & $+$dequant overhead\\\bottomrule
\end{tabular}
\caption{State properties before and after QLoRA.}
\end{table}

\newpage
%% ------------------------------------------------------------
\section{Adaptive LoRA (AdaLoRA)}

\subsection{Mathematical Foundation}
AdaLoRA parameterizes incremental weight updates using a singular value decomposition (SVD) representation to dynamically adjust the adapter rank during training:
\begin{equation}
  \Delta W = P \Lambda Q
\end{equation}
where $P \in \mathbb{R}^{d \times r}$ and $Q \in \mathbb{R}^{r \times k}$ are left and right singular vectors, and $\Lambda = \text{diag}(\lambda_1, \dots, \lambda_r) \in \mathbb{R}^{r \times r}$ contains singular values.
To prevent SVD divergence and maintain parameter stability, we enforce orthogonality constraints:
\begin{equation}
  \mathcal{L}_{\text{orth}}(P, Q) = \gamma \left( \|P^\top P - I\|_F^2 + \|Q Q^\top - I\|_F^2 \right)
\end{equation}
We estimate the training importance of each singular triplet $(P_{*,i}, \lambda_i, Q_{i,*})$ using gradient-weight product magnitude:
\begin{equation}
  S_i = s(\lambda_i) + \frac{1}{d} \sum_{j=1}^d s(P_{j,i}) + \frac{1}{k} \sum_{j=1}^k s(Q_{i,j})
\end{equation}
where $s(\theta) = |\theta \cdot \nabla_\theta \mathcal{L}|$. At designated intervals, triplets with the lowest importance scores $S_i$ are pruned by zeroing out their singular values $\lambda_i$, dynamically reducing the adapter rank budget.

\subsection{Architecture Flow Diagram}

\begin{figure}[h]
\centering
\begin{adjustbox}{max width=\linewidth}
\begin{tikzpicture}[node distance=0.7cm and 1.1cm]
  \node[datblk] (x) {\textbf{Input $x$}};
  \node[frzblk, right=1.5cm of x] (W0) {\textbf{Frozen $W_0$}};
  
  % Parallel SVD Adapter path
  \node[trnblk, above=1.0cm of W0] (Q) {\textbf{$Q$ (down-proj)}};
  \node[trnblk, right=1.2cm of Q] (L) {\textbf{$\Lambda$ (diag)}};
  \node[trnblk, right=1.2cm of L] (P) {\textbf{$P$ (up-proj)}};
  
  \node[grayblk, above=0.6cm of L] (prune) {\textbf{Rank Pruner}};
  \node[sumnode, right=1.5cm of W0] (sum) {$\oplus$};
  \node[mrgblk, right=1.5cm of sum] (out) {\textbf{Output $h$}};

  % Connections
  \draw[arr] (x)--(W0);
  \draw[arr] (x.north) |- (Q.west);
  \draw[arr] (Q)--(L);
  \draw[arr] (L)--(P);
  \draw[arr] (P.east) -| (sum.north);
  \draw[arr] (W0)--(sum);
  \draw[arr] (sum)--(out);
  
  % Pruning monitor line
  \draw[arr,dashed,red] (prune)--node[right,font=\tiny]{prunes $\lambda_i$}(L);
  \node[red,font=\tiny,align=center] at ($(prune.north)+(0,0.25)$) {$S_i{<}\gamma \Rightarrow \lambda_i{=}0$};
\end{tikzpicture}
\end{adjustbox}
\caption{AdaLoRA: SVD-decomposed update with importance-based rank pruning per layer.}
\end{figure}

\subsection{State Transition Table}
\begin{table}[h]\centering
\begin{tabular}{lll}\toprule
\textbf{Property} & \textbf{Before} & \textbf{After}\\\midrule
Rank per layer & Fixed $r$ & Adaptive (varies by importance)\\
Update structure & $BA$ product & $P\Lambda Q$ SVD triplet\\
Parameter budget & Fixed & Redistributed across layers\\
Low-importance ranks & Included & Pruned ($\lambda_i = 0$)\\
Training stability & Standard & Constrained by SVD orthogonality\\\bottomrule
\end{tabular}
\caption{State properties before and after AdaLoRA.}
\end{table}

\newpage
%% ------------------------------------------------------------
\section{Weight-Decomposed LoRA (DoRA)}

\subsection{Mathematical Foundation}
DoRA decomposes the weight matrix $W \in \mathbb{R}^{k \times d}$ into its magnitude vector $m \in \mathbb{R}^k$ and direction matrix $V \in \mathbb{R}^{k \times d}$:
\begin{equation}
  W = m \cdot \frac{V}{\|V\|_r} = m \cdot \frac{W_0 + \Delta W}{\|W_0 + \Delta W\|_r}
\end{equation}
where $\|\cdot\|_r$ is the L2 norm computed along each row (corresponding to output channels). In parameter-efficient training, we freeze the direction matrix to the base weight $V = W_0$ and introduce low-rank directional updates using adapters $\Delta V = \frac{\alpha}{r} B A$ (where $A \in \mathbb{R}^{r \times d}$ and $B \in \mathbb{R}^{k \times r}$):
\begin{equation}
  W = m \cdot \frac{W_0 + \frac{\alpha}{r} B A}{\|W_0 + \frac{\alpha}{r} B A\|_r}
\end{equation}
Only the magnitude vector $m$ and low-rank directional matrices $B$ and $A$ are trainable.
The gradient flow for the directional components is:
\begin{equation}
  \nabla_V \mathcal{L} = \frac{m}{\|V\|_c} \left( \nabla_W \mathcal{L} - \frac{V}{\|V\|_c^2} \odot \left( V^\top \nabla_W \mathcal{L} \right) \right)
\end{equation}
This mathematically separates the updates to magnitude and direction, mimicking full parameter fine-tuning behavior while retaining low parameter overhead.

\subsection{Architecture Flow Diagram}

\begin{figure}[h]
\centering
\begin{adjustbox}{max width=\linewidth}
\begin{tikzpicture}[node distance=0.7cm and 1.1cm]
  \node[frzblk] (W0) {\textbf{$W_0$ (frozen)}};
  \node[trnblk, above=1.0cm of W0] (ba) {\textbf{$BA$ (LoRA)}};
  
  \node[sumnode, right=1.2cm of W0] (sumV) {$\oplus$};
  
  \node[grayblk, right=1.0cm of sumV] (norm) {\textbf{Norm $\|\cdot\|_r$}};
  \node[sumnode, right=1.2cm of norm] (div) {$\div$};
  
  \node[trnblk, above=1.0cm of div] (m) {\textbf{Magnitude $m$}};
  \node[sumnode, right=1.2cm of div] (mul) {$\otimes$};
  \node[mrgblk, right=1.2cm of mul] (out) {\textbf{Weight $W'$}};

  % Connections
  \draw[arr] (W0)--(sumV);
  \draw[arr] (ba.east) -| (sumV.north);
  \draw[arr] (sumV)--(norm);
  \draw[arr] (sumV.east) -- ++(0.2,0) |- ($(div.west)+(0,0.2)$);
  \draw[arr] (norm)--(div);
  \draw[arr] (div)--(mul);
  \draw[arr] (m)--(mul);
  \draw[arr] (mul)--(out);

  \node[font=\tiny,below=0.3cm of norm] {direction component};
  \node[font=\tiny,above=0.3cm of m] {amplitude scaling};
\end{tikzpicture}
\end{adjustbox}
\caption{DoRA: weight decomposed into direction (LoRA-adapted) and magnitude (trainable scalar).}
\end{figure}

\subsection{State Transition Table}
\begin{table}[h]\centering
\begin{tabular}{lll}\toprule
\textbf{Property} & \textbf{Before (LoRA)} & \textbf{After (DoRA)}\\\midrule
Magnitude component & Implicit in $W_0$ & Explicit trainable $m$\\
Direction component & $W_0 + BA$ unnorm & Column-normalised\\
Full FT similarity & Lower & Higher\\
Trainable params & $r(d+k)$ & $r(d+k) + k$\\
Catastrophic forgetting & Low & Lower\\\bottomrule
\end{tabular}
\caption{State properties comparing LoRA and DoRA.}
\end{table}

\newpage
%% ------------------------------------------------------------
\section{IA$^3$ (Infused Adapter by Inhibiting and Amplifying Inner Activations)}

\subsection{Mathematical Foundation}
IA$^3$ introduces element-wise scaling vectors to directly modify inner activations in attention keys, values, and intermediate feed-forward layers. Let the scaling vectors be $l_K \in \mathbb{R}^d$, $l_V \in \mathbb{R}^d$, and $l_{FFN} \in \mathbb{R}^{d_{FFN}}$.
The attention layer activations are computed as:
\begin{align}
  \tilde{K} &= K \odot l_K = K \cdot \text{diag}(l_K) \\
  \tilde{V} &= V \odot l_V = V \cdot \text{diag}(l_V) \\
  \text{Attention}(Q, \tilde{K}, \tilde{V}) &= \text{softmax}\left( \frac{Q \tilde{K}^\top}{\sqrt{d_k}} \right) \tilde{V}
\end{align}
In the feed-forward network (FFN), let $W_1$ and $W_2$ be the weight matrices. The intermediate activations are scaled before the output projection:
\begin{equation}
  \text{FFN}(x) = \sigma\bigl(x W_1 \odot l_{FFN}\bigr) W_2 = \sigma\bigl(x W_1 \cdot \text{diag}(l_{FFN})\bigr) W_2
\end{equation}
The base model weights $W_Q, W_K, W_V, W_1, W_2$ are frozen. This diagonal matrix transformation scales the columns of the projection matrices with minimal parameter updates.

\subsection{Architecture Flow Diagram}

\begin{figure}[h]
\centering
\begin{adjustbox}{max width=\linewidth}
\begin{tikzpicture}[node distance=0.6cm and 1.0cm]
  \node[datblk] (x) {\textbf{Input $x$}};
  
  % Attention Block
  \node[frzblk, right=1.0cm of x] (qkv) {\textbf{Compute} \\ ($Q, K, V$)};
  \node[trnblk, above right=0.3cm and 0.8cm of qkv] (lk) {\textbf{$l_K \odot K$}};
  \node[trnblk, below right=0.3cm and 0.8cm of qkv] (lv) {\textbf{$l_V \odot V$}};
  \node[frzblk, right=2.8cm of qkv] (attn) {\textbf{Attention} \\ (Softmax)};
  
  % FFN Block
  \node[frzblk, right=1.0cm of attn] (w1) {\textbf{FFN $W_1$}};
  \node[trnblk, right=0.8cm of w1] (lffn) {\textbf{$l_{FFN} \odot$}};
  \node[frzblk, right=0.8cm of lffn] (w2) {\textbf{FFN $W_2$}};
  \node[mrgblk, right=1.0cm of w2] (out) {\textbf{Output}};

  % Connections
  \draw[arr] (x)--(qkv);
  \draw[arr] (qkv.east) -- ++(0.2,0) |- (lk.west);
  \draw[arr] (qkv.east) -- ++(0.2,0) |- (lv.west);
  \draw[arr] (qkv.east) -- ++(0.2,0) |- ($(attn.west)+(0,0.4)$); % Q direct
  \draw[arr] (lk.east) -| ($(attn.west)+(0,0.15)$);
  \draw[arr] (lv.east) -| ($(attn.west)+(0,-0.15)$);
  
  \draw[arr] (attn)--(w1);
  \draw[arr] (w1)--(lffn);
  \draw[arr] (lffn)--(w2);
  \draw[arr] (w2)--(out);
\end{tikzpicture}
\end{adjustbox}
\caption{IA$^3$: three tiny scaling vectors rescale keys, values, and FFN intermediate activations.}
\end{figure}

\subsection{State Transition Table}
\begin{table}[h]\centering
\begin{tabular}{lll}\toprule
\textbf{Property} & \textbf{Before} & \textbf{After}\\\midrule
Trainable params & 0 (frozen) & $3dL$ scalars\\
Init value of $l$ vectors & N/A & $\mathbf{1}$ (identity behaviour)\\
K, V, FFN activations & Unchanged & Element-wise scaled\\
Param overhead vs LoRA & N/A & ${\sim}10\times$ fewer params\\
Task adaptation strength & None & Moderate\\\bottomrule
\end{tabular}
\caption{State properties before and after IA$^3$ injection.}
\end{table}

\newpage
%% ------------------------------------------------------------
\section{Adapter Tuning}

\subsection{Mathematical Foundation}
Adapter tuning inserts small bottleneck layers into the transformer block, typically after attention and feed-forward modules. Let $h \in \mathbb{R}^d$ be the input hidden state.
In the Pfeiffer architecture, a single adapter is inserted per block. In the Houlsby architecture, two adapters are inserted per block. The mathematical transformation for an adapter block is:
\begin{equation}
  h' = h + \sigma(h W_{down}) W_{up}
\end{equation}
where $W_{down} \in \mathbb{R}^{d \times r}$ projects to a bottleneck dimension $r \ll d$, $\sigma$ is an activation function (typically GELU), and $W_{up} \in \mathbb{R}^{r \times d}$ projects back to the original dimension.
To ensure the adapter behaves as an identity mapping at the start of training, $W_{up}$ is initialized to zero:
\begin{equation}
  h' \bigr|_{t=0} = h + \sigma(h W_{down}) \cdot 0 = h
\end{equation}

\subsection{Architecture Flow Diagram}

\begin{figure}[h]
\centering
\begin{adjustbox}{max width=\linewidth}
\begin{tikzpicture}[node distance=0.5cm and 0.9cm]
  \node[datblk,minimum width=2.2cm] (h) {\textbf{Hidden $h$} \\ (from layer)};
  \node[trnblk,right=of h,minimum width=2.0cm] (dn) {\textbf{Down-proj} \\ ($W_{\text{dn}}$ ($d{\to}r$))};
  \node[grayblk,right=of dn,minimum width=1.8cm] (act) {\textbf{ReLU}};
  \node[trnblk,right=of act,minimum width=2.0cm] (up) {\textbf{Up-proj} \\ ($W_{\text{up}}$ ($r{\to}d$))};
  \node[sumnode,right=of up] (add) {$+$};
  \node[mrgblk,right=of add,minimum width=2.0cm] (out) {\textbf{Adapter} \\ (Output $h'$)};

  \draw[arr] (h)--(dn);
  \draw[arr] (dn)--(act);
  \draw[arr] (act)--(up);
  \draw[arr] (up)--(add);
  \draw[arr] (h.south) -- ++(0,-0.6) -| (add.south) node[midway,below,font=\tiny]{skip connection};
  \draw[arr] (add)--(out);

  \node[red,font=\tiny,below=0.2cm of dn] {bottleneck $r \ll d$};
\end{tikzpicture}
\end{adjustbox}
\caption{Adapter module: down-project to rank $r$, apply nonlinearity, up-project back, plus skip connection.}
\end{figure}

\subsection{State Transition Table}
\begin{table}[h]\centering
\begin{tabular}{lll}\toprule
\textbf{Property} & \textbf{Before} & \textbf{After}\\\midrule
Sub-layer output path & Direct pass-through & $+$adapter bottleneck\\
Transformer weights & Trainable & Frozen\\
New params per layer & 0 & $2dr + 2d$\\
Adapter init & N/A & Near-identity ($W_{\text{dn}} \sim \mathcal{N}$)\\
Inference latency & Baseline & $+$two extra matmuls per layer\\\bottomrule
\end{tabular}
\caption{State properties before and after adapter insertion.}
\end{table}

\newpage
%% ------------------------------------------------------------
\section{Prefix Tuning}

\subsection{Mathematical Foundation}
Prefix Tuning prepends continuous task-specific virtual prefix vectors $P_K, P_V \in \mathbb{R}^{l \times d}$ directly to the Key and Value matrices at every transformer layer. Let the original Key and Value matrices for sequence length $T$ be $K, V \in \mathbb{R}^{T \times d}$. The modified Key and Value matrices become:
\begin{align}
  \tilde{K} &= [P_K ; K] \in \mathbb{R}^{(l+T) \times d} \\
  \tilde{V} &= [P_V ; V] \in \mathbb{R}^{(l+T) \times d}
\end{align}
To stabilize optimization, during training the prefix is reparameterized using a Multilayer Perceptron (MLP) over a small source embedding $E_k$:
\begin{equation}
  P_K = \text{MLP}_K(E_k), \quad P_V = \text{MLP}_V(E_k)
\end{equation}
where $E_k \in \mathbb{R}^{l \times d_{\text{mid}}}$. The MLP is discarded at inference, and only the generated static prefix matrices $P_K, P_V$ are saved and used.
The attention output is computed as:
\begin{equation}
  \text{Attention}(Q, \tilde{K}, \tilde{V}) = \text{softmax}\left( \frac{Q [P_K ; K]^\top}{\sqrt{d_k}} \right) [P_V ; V]
\end{equation}

\subsection{Architecture Flow Diagram}

\begin{figure}[h]
\centering
\begin{adjustbox}{max width=\linewidth}
\begin{tikzpicture}[node distance=0.6cm and 1.0cm]
  \node[trnblk,minimum width=2.5cm] (prefix) {\textbf{Trainable Prefix} \\ ($P_k, P_v$ (per layer))};
  \node[datblk, below=1.2cm of prefix] (tokens) {\textbf{Input Token} \\ (Keys \& Values)};
  \node[grayblk, right=1.5cm of $(prefix)!0.5!(tokens)$] (concat) {\textbf{Concat} \\ ($[P_k;K],\,[P_v;V]$)};
  \node[frzblk, right=of concat] (attn) {\textbf{Frozen} \\ (Attention)};
  \node[mrgblk, right=of attn] (out) {\textbf{Contextual} \\ (Output)};

  \draw[arr] (prefix.east) -| (concat.north);
  \draw[arr] (tokens.east) -| (concat.south);
  \draw[arr] (concat)--(attn);
  \draw[arr] (attn)--(out);

  \node[font=\tiny,green!60!black] at ($(prefix.south)+(0,-0.2)$) {only these trained};
  \node[font=\tiny,blue!60!black] at ($(attn.south)+(0,-0.2)$) {all frozen};
\end{tikzpicture}
\end{adjustbox}
\caption{Prefix Tuning: trainable prefix vectors extend the key/value context at every attention layer.}
\end{figure}

\subsection{State Transition Table}
\begin{table}[h]\centering
\begin{tabular}{lll}\toprule
\textbf{Property} & \textbf{Before} & \textbf{After}\\\midrule
Attention context length & $T$ tokens & $T + l$ tokens (prefix prepended)\\
Transformer params & Trainable & Fully frozen\\
Trainable params & 0 & $2 \times l \times d_h \times L$\\
Prefix init & N/A & Via MLP reparameterisation\\
Task switch at inference & Reload model & Swap prefix only\\\bottomrule
\end{tabular}
\caption{State properties before and after prefix tuning.}
\end{table}

\newpage
%% ------------------------------------------------------------
\section{Prompt Tuning}

\subsection{Mathematical Foundation}
Prompt tuning prepends $p$ learnable embedding vectors $P \in \mathbb{R}^{p \times d_{emb}}$ to the input word representations:
\begin{equation}
  \tilde{X} = [P ; X_e] \in \mathbb{R}^{(p + T) \times d_{emb}}
\end{equation}
where $X_e \in \mathbb{R}^{T \times d_{emb}}$ is the original sequence word embedding matrix. The entire transformer backbone parameters $\theta_{base}$ remain frozen, and gradients are backpropagated only to the prompt parameters $P$:
\begin{equation}
  \nabla_P \mathcal{L} = \frac{\partial \mathcal{L}}{\partial \tilde{X}_{1..p}}
\end{equation}
During backpropagation, the gradients are passed through all layers of the frozen transformer and update only the prompt embeddings.

\subsection{Architecture Flow Diagram}

\begin{figure}[h]
\centering
\begin{adjustbox}{max width=\linewidth}
\begin{tikzpicture}[node distance=0.7cm and 1.0cm]
  \node[trnblk,minimum width=2.2cm] (soft) {\textbf{Soft Prompt} \\ \textbf{$P \in \mathbb{R}^{k\times d}$} \\ (trainable)};
  \node[datblk, right=of soft] (emb) {\textbf{Input} \\ \textbf{Embeddings} \\ ($E(x_{1..T})$)};
  \node[grayblk, right=of emb] (concat) {\textbf{Concat} \\ ($[\tilde{x}]$)};
  \node[frzblk, right=of concat] (lm) {\textbf{Frozen} \\ \textbf{Transformer} \\ ($f_\theta$)};
  \node[lossblk, right=of lm] (loss) {\textbf{SFT Loss} \\ (on $y$)};

  \draw[arr] (soft.east) -| (concat.north);
  \draw[arr] (emb)--(concat);
  \draw[arr] (concat)--(lm);
  \draw[arr] (lm)--(loss);

  \node[font=\tiny,green!60!black,align=center,below=0.25cm of soft] {only $P$\\is trained};
\end{tikzpicture}
\end{adjustbox}
\caption{Prompt Tuning: soft tokens prepended at the embedding layer; the full transformer is frozen.}
\end{figure}

\subsection{State Transition Table}
\begin{table}[h]\centering
\begin{tabular}{lll}\toprule
\textbf{Property} & \textbf{Before} & \textbf{After}\\\midrule
Trainable params & Full model & $k \times d_{\text{emb}}$ only\\
Prompt format & Discrete hand-crafted & Continuous soft vectors\\
Transformer weights & All trainable & All frozen\\
Task switching & Reload fine-tuned model & Swap $P$ only\\
Sensitivity to init & N/A & Strong (init from class labels helps)\\\bottomrule
\end{tabular}
\caption{State properties before and after prompt tuning.}
\end{table}

\newpage
%% ------------------------------------------------------------
\section{P-Tuning v2}

\subsection{Mathematical Foundation}
P-Tuning v2 extends prompt tuning by prepending learnable virtual tokens $P^{(l)} \in \mathbb{R}^{p \times 2 d_{attn}}$ directly to the attention key and value states at every layer $l \in [1, L]$ of the network. Unlike prefix tuning, P-Tuning v2 removes the MLP reparameterization network, directly optimizing the prefix parameters $P^{(l)}$:
\begin{align}
  \tilde{K}^{(l)} &= [P^{(l)}_K ; K^{(l)}] \\
  \tilde{V}^{(l)} &= [P^{(l)}_V ; V^{(l)}]
\end{align}
This deep prefix injection resolves the capacity limitation of shallow prompt tuning, allowing small models ($<10\text{B}$ parameters) to achieve performance comparable to full fine-tuning.

\subsection{Architecture Flow Diagram}

\begin{figure}[h]
\centering
\begin{adjustbox}{max width=\linewidth}
\begin{tikzpicture}[node distance=0.5cm and 0.9cm]
  \node[datblk,minimum width=2.0cm] (emb) {\textbf{Embedding} \\ (Layer)};
  \node[frzblk, right=0.9cm of emb,minimum width=2.0cm] (l1) {\textbf{Layer 1} \\ (Attn + FFN)};
  \node[frzblk, right=0.9cm of l1,minimum width=2.0cm] (l2) {\textbf{Layer 2} \\ (Attn + FFN)};
  \node[grayblk, right=0.6cm of l2,minimum width=0.6cm] (dots) {\textbf{$\cdots$}};
  \node[frzblk, right=0.6cm of dots,minimum width=2.0cm] (lL) {\textbf{Layer $L$} \\ (Attn + FFN)};

  \node[trnblk,minimum width=1.6cm,above=0.6cm of emb] (p0) {\textbf{$P^{(0)}$}};
  \node[trnblk,minimum width=1.6cm,above=0.6cm of l1] (p1) {\textbf{$P^{(1)}$}};
  \node[trnblk,minimum width=1.6cm,above=0.6cm of l2] (p2) {\textbf{$P^{(2)}$}};
  \node[trnblk,minimum width=1.6cm,above=0.6cm of lL] (pL) {\textbf{$P^{(L)}$}};

  \foreach \a/\b in {emb/l1, l1/l2, l2/dots, dots/lL}
    \draw[arr] (\a)--(\b);
  \foreach \p/\l in {p0/emb, p1/l1, p2/l2, pL/lL}
    \draw[arr,green!70!black] (\p)--(\l);

  \node[green!60!black,font=\tiny,align=center] at ($(p1.east)+(0.35,0)$) {deep\\prefix};
\end{tikzpicture}
\end{adjustbox}
\caption{P-Tuning v2: independent trainable prefix tokens injected at every transformer layer.}
\end{figure}

\subsection{State Transition Table}
\begin{table}[h]\centering
\begin{tabular}{lll}\toprule
\textbf{Property} & \textbf{Before (Prompt Tuning)} & \textbf{After (P-Tuning v2)}\\\midrule
Prefix layers & Embedding only & All $L$ layers\\
Trainable params & $k \times d$ & $L \times l \times d$\\
NLU performance & Degrades at small scale & Competitive at all scales\\
Conditioning depth & Shallow (first layer) & Deep (every layer)\\
Task-specific storage & Single $P$ matrix & $L$ prefix matrices\\\bottomrule
\end{tabular}
\caption{State properties comparing Prompt Tuning and P-Tuning v2.}
\end{table}

